\documentclass[11pt,a4paper]{article}

\usepackage[T1]{fontenc}
\usepackage[utf8]{inputenc}
\usepackage[italian]{babel}
\usepackage{lmodern}
\usepackage{geometry}
\geometry{margin=2.2cm}
\usepackage{graphicx}
\usepackage{float}
\usepackage{amsmath,amssymb,mathtools}
\usepackage{booktabs}
\usepackage{xcolor}
\usepackage{hyperref}
\usepackage{enumitem}
\usepackage{titlesec}
\usepackage{fancyhdr}
\usepackage[most]{tcolorbox}
\usepackage{caption}
\usepackage{subcaption}
\usepackage{microtype}
\usepackage{parskip}
\usepackage{longtable}
\usepackage{array}

\hypersetup{
    colorlinks=true,
    linkcolor=blue!50!black,
    urlcolor=blue!50!black,
    pdftitle={Tecnologie per il Controllo dei Robot - Appunti del corso},
    pdfauthor={OpenAI}
}

\definecolor{accent}{RGB}{82, 97, 170}
\definecolor{accentlight}{RGB}{235, 239, 252}
\definecolor{grayline}{RGB}{220, 223, 230}

\newcommand{\projectbasepath}{}
\IfFileExists{images/page-01.png}{}{%
    \renewcommand{\projectbasepath}{Appunti_Tecnologie_per_il_controllo_dei_Robot/}%
}

\titleformat{\section}{\Large\bfseries\color{accent}}{\thesection}{0.6em}{}
\titleformat{\subsection}{\large\bfseries\color{black}}{\thesubsection}{0.6em}{}
\titleformat{\subsubsection}{\normalsize\bfseries\color{black}}{\thesubsubsection}{0.6em}{}
\captionsetup{
    font=small,
    labelfont=bf,
    justification=raggedright,
    singlelinecheck=false
}

\setlength{\headheight}{14pt}
\pagestyle{fancy}
\fancyhf{}
\fancyhead[L]{Tecnologie per il Controllo dei Robot}
\fancyhead[R]{Appunti del corso}
\fancyfoot[C]{\thepage}
\setlength{\textfloatsep}{10pt plus 2pt minus 2pt}
\setlength{\floatsep}{10pt plus 2pt minus 2pt}
\setlength{\intextsep}{10pt plus 2pt minus 2pt}

\newtcolorbox{infobox}{
    colback=accentlight,
    colframe=accent,
    boxrule=0.6pt,
    arc=2mm,
    left=2mm,right=2mm,top=1.5mm,bottom=1.5mm
}

\newtcolorbox{focusbox}{
    colback=accentlight,
    colframe=accent,
    boxrule=0.6pt,
    arc=2mm,
    left=2mm,right=2mm,top=1.5mm,bottom=1.5mm
}

\newtcolorbox{keybox}{
    colback=grayline!30,
    colframe=black!70,
    boxrule=0.6pt,
    arc=2mm,
    left=2mm,right=2mm,top=1.5mm,bottom=1.5mm
}

\newcommand{\slidefig}[2]{
\begin{figure}[H]
\centering
\includegraphics[width=0.86\linewidth]{\projectbasepath#1}
\caption{#2}
\end{figure}
}

\newcommand{\slidepair}[3]{
\begin{figure}[H]
\centering
\begin{subfigure}[t]{0.48\linewidth}
\centering
\includegraphics[width=\linewidth]{\projectbasepath#1}
\end{subfigure}\hfill
\begin{subfigure}[t]{0.48\linewidth}
\centering
\includegraphics[width=\linewidth]{\projectbasepath#2}
\end{subfigure}
\caption{#3}
\end{figure}
}

\newcommand{\slidetriple}[4]{
\begin{figure}[H]
\centering
\begin{subfigure}[t]{0.32\linewidth}
\centering
\includegraphics[width=\linewidth]{\projectbasepath#1}
\end{subfigure}\hfill
\begin{subfigure}[t]{0.32\linewidth}
\centering
\includegraphics[width=\linewidth]{\projectbasepath#2}
\end{subfigure}\hfill
\begin{subfigure}[t]{0.32\linewidth}
\centering
\includegraphics[width=\linewidth]{\projectbasepath#3}
\end{subfigure}
\caption{#4}
\end{figure}
}

\begin{document}

\hypersetup{pageanchor=false}
\begin{titlepage}
\centering
\vspace*{2cm}
{\Huge\bfseries Tecnologie per il Controllo dei Robot \par}
\vspace{0.4cm}
{\LARGE Appunti del corso, traduzione ragionata e approfondimento\par}
\vspace{1cm}
{\large Basato sulle slide del corso di Lorenzo Pollini\\
LM in Ingegneria Robotica e dell'Automazione, Università di Pisa\par}
\vfill
\begin{infobox}
\textbf{Impostazione del documento.} Questi appunti seguono l’ordine delle slide originali, ma ne riorganizzano il contenuto in forma più leggibile: ogni blocco traduce il testo essenziale delle slide, ne chiarisce il significato tecnico, esplicita le equazioni e aggiunge il contesto necessario per lo studio e il ripasso.
\end{infobox}
\vfill
{\large Anno accademico 2025/2026\par}
\end{titlepage}
\hypersetup{pageanchor=true}

\tableofcontents
\clearpage

% =========================================================
% LEZIONE 1-2
% =========================================================

\section{Lezione 1-2}

\subsection*{Visione generale della lezione}
Le prime due lezioni introducono il ruolo dei sensori nei sistemi robotici e costruiscono il lessico di base necessario per tutto il corso: natura dell’informazione misurata, digitalizzazione, non idealità statiche, rumore e primi modelli di calibrazione. L’obiettivo non è ancora studiare sensori specifici in dettaglio, ma capire come nasce il dato che poi verrà usato dal controllore.

% =========================================================
\subsection{Introduzione, struttura del corso e ruolo dei sensori}

\slidepair{images/page-01.png}{images/page-02.png}{Inquadramento del corso e ruolo della misura nei sistemi robotici.}

Le prime due slide inquadrano il corso e chiariscono da subito il suo punto di vista: nel controllo dei robot non basta progettare un buon regolatore, bisogna anche capire come il sistema misura il mondo fisico. La qualità della retroazione dipende infatti dalla qualità della catena di misura.

\slidepair{images/page-03.png}{images/page-04.png}{La qualità del controllo dipende dalla qualità della catena di misura e della retroazione.}

Le slide successive spostano l’attenzione proprio su questo fatto: un controllore raffinato, alimentato da misure lente, rumorose o mal calibrate, può comportarsi peggio di un controllore più semplice ma costruito su dati affidabili.

\slidepair{images/page-05.png}{images/page-06.png}{Schema generale del sistema robotico: misura, elaborazione e azione.}

Per questo il corso parte dai sensori. In un sistema robotico, i sensori trasformano grandezze fisiche in segnali utilizzabili dal controllore, mentre gli attuatori realizzano sul mondo fisico le azioni calcolate dall’algoritmo. In forma schematica:
\[
\text{misura} \longrightarrow \text{elaborazione} \longrightarrow \text{azione}.
\]
Un altro punto importante è che il controllo non agisce mai sulla grandezza ``vera'', ma sulla sua misura o sulla sua stima. Se si indica con $y_{\text{true}}$ la variabile fisica reale e con $y_{\text{meas}}$ quella acquisita dal sistema, il controllore lavora in pratica su $y_{\text{meas}}$.

In un loop di retroazione il passaggio corretto da tenere in mente e` quindi
\[
e(t)=r(t)-y_{\text{meas}}(t),
\qquad
u(t)=C(s)e(t),
\qquad
y_{\text{meas}}(t)=H_s(s)\,y_{\text{true}}(t)+n(t),
\]
dove $H_s(s)$ raccoglie l’intera catena di misura e $n(t)$ rappresenta rumore e disturbi introdotti dal sensore. Questa scrittura fa vedere subito il punto chiave del corso: il sensore non e` un blocco neutro, ma un elemento che modifica direttamente l’errore visto dal controllore.

Dal punto di vista sistemistico, la catena di misura determina anche quali stati o grandezze del robot sono effettivamente osservabili in modo affidabile. Una grandezza fisica puo` essere presente nel sistema reale ma, se la misura e` rumorosa, lenta o non direttamente disponibile, il controllore dovra` lavorare su una stima filtrata o su una variabile correlata. E` per questo che sensori, stimatori e controllo vanno letti come parti dello stesso problema.

\begin{keybox}
Nel controllo robotico la qualità della misura è parte della qualità del loop di controllo, non un dettaglio separato.
\end{keybox}

% =========================================================
\subsection{Sensori attivi, passivi e distinzione tra sensore e trasduttore}

\slidepair{images/page-07.png}{images/page-08.png}{Sensori attivi e passivi: distinzione di principio e lettura applicativa.}

Le prime due slide introducono la distinzione tra sensori attivi e passivi. Un sensore attivo richiede una sorgente di alimentazione o di eccitazione per funzionare, come accade in molti sensori piezoresistivi o capacitivi. Un sensore passivo, invece, sfrutta direttamente l’energia del fenomeno fisico misurato, come nelle termocoppie o in alcuni dispositivi fotovoltaici. Nei sensori moderni, però, questa distinzione va letta con attenzione: il principio di trasduzione può essere passivo, mentre il modulo reale contiene comunque elettronica attiva.

\slidefig{images/page-09.png}{Distinzione tra trasduttore e sensore completo.}

La terza slide del blocco distingue poi \emph{trasduttore} e \emph{sensore}. Il trasduttore è l’elemento che realizza la conversione fisica. Il sensore è il dispositivo completo che può includere amplificazione, filtraggio, ADC e interfaccia di comunicazione. In modo schematico:
\[
y_{\text{meas}} = \mathcal{E}(\mathcal{T}(x)),
\]
dove $\mathcal{T}$ rappresenta la trasduzione fisica e $\mathcal{E}$ la catena elettronica e digitale che segue.

Questa distinzione non e` solo terminologica. In un progetto reale il trasduttore puo` essere eccellente dal punto di vista fisico, ma il sensore completo risultare mediocre a causa di loading elettrico, rumore dell’elettronica o conversione digitale insufficiente. Per questo, quando si valutano datasheet o specifiche, conviene sempre chiedersi dove finisca il principio di misura e dove inizino invece i limiti dell’interfacciamento.

Un modo equivalente di leggere il problema e` tramite impedenze e guadagni: il trasduttore genera un segnale sorgente con propria impedenza interna, l’elettronica successiva lo carica e lo condiziona, l’ADC lo discretizza. Se uno di questi passaggi e` mal progettato, la misura finale si degrada anche se il principio fisico di base sarebbe teoricamente molto buono.

\subsection{Sensori assoluti e relativi. Natura dell’informazione}

\slidepair{images/page-10.png}{images/page-11.png}{Sensori assoluti, sensori relativi e natura del dato misurato.}

Le prime due slide distinguono tra sensori assoluti e relativi. Un sensore assoluto fornisce il valore della grandezza rispetto a una scala fisica o a una codifica assoluta. Un sensore relativo, invece, misura differenze rispetto a un riferimento locale o a una condizione iniziale. In robotica questa distinzione è fondamentale, per esempio nel confronto tra encoder assoluti e incrementali.

\slidefig{images/page-12.png}{Segnali continui, discreti e disponibilità temporale dell’informazione.}

Viene poi introdotta un’altra distinzione: quella tra grandezze continue e discrete, e tra segnali continui e discreti nel tempo. Molte grandezze fisiche sono continue, ma l’informazione che il controllore riceve può essere analogica, digitale, logica o codificata. Inoltre, in un sistema digitale la misura è disponibile solo nei tempi campionati
\[
t_k = kT_s.
\]
Questa discretizzazione temporale è un passaggio centrale perché collega direttamente il tema dei sensori al controllo digitale.

Sul piano concettuale, un sensore relativo fornisce spesso una derivata, un incremento o una differenza locale che deve poi essere integrata o accumulata per ricostruire la grandezza d’interesse. Se ad esempio si misura una velocita` $\dot x$ invece della posizione $x$, la ricostruzione richiede
\[
x(t)=x(t_0)+\int_{t_0}^{t}\dot x(\tau)\,d\tau.
\]
Questo passaggio e` molto delicato, perche' ogni bias o rumore a bassa frequenza presente nella misura di $\dot x$ tende ad accumularsi nel tempo sotto forma di drift sulla stima di posizione.

Al contrario, un sensore assoluto riduce proprio questa dipendenza dalla storia del sistema: l’informazione utile e` gia` contenuta nel dato letto nell’istante corrente. Per questo, in robotica, la scelta tra misura assoluta e relativa va vista anche come scelta tra due modi diversi di gestire memoria, deriva e procedure di inizializzazione.

\begin{focusbox}
\textbf{Idea chiave.} Una grandezza fisica continua può essere convertita in formati molto diversi. Il formato del dato influisce direttamente su acquisizione, elaborazione e controllo.
\end{focusbox}

\subsection{Campionamento}

\slidefig{images/page-13.png}{Definizione del campionamento nel tempo.}

La prima slide sul campionamento chiarisce che il segnale non viene osservato in modo continuo, ma solo in una sequenza discreta di istanti:
\[
x[k] = x(kT_s).
\]
Questa operazione riduce l’informazione temporale disponibile e rende cruciale la scelta del periodo di campionamento. Se la frequenza di campionamento è troppo bassa rispetto al contenuto in frequenza del segnale, compaiono fenomeni di aliasing e il dato digitale non rappresenta più correttamente la grandezza fisica.

Nel caso ideale di segnale limitato in banda, il criterio di Nyquist-Shannon richiede
\[
f_s > 2f_{\max},
\]
dove $f_{\max}$ e` la massima frequenza significativa del segnale analogico prima del campionamento. Nella pratica ingegneristica, pero`, non ci si limita quasi mai a stare appena sopra il doppio della banda utile: si introduce un margine, si usa un filtro anti-alias analogico e si sceglie spesso
\[
f_s \approx 5\text{--}20\,f_B,
\]
dove $f_B$ e` la banda che si vuole osservare o controllare con buona fedelta`.

\slidepair{images/page-14.png}{images/page-15.png}{Effetti del campionamento e relazione tra periodo di campionamento e dinamica osservata.}

Dal punto di vista del controllo robotico, il campionamento non è solo una questione implementativa. Un campionamento lento aumenta il ritardo percepito dal controllore e peggiora la qualità della retroazione. Un campionamento molto alto aumenta invece il carico computazionale e può rendere più evidente il rumore. La scelta di $T_s$ va quindi fatta in modo coerente con la banda della dinamica che si vuole osservare o controllare.

Un secondo punto spesso sottovalutato riguarda le grandezze derivate. Se da una posizione campionata si ricava la velocita` con una differenza finita,
\[
\dot x[k]\approx \frac{x[k]-x[k-1]}{T_s},
\]
il rumore ad alta frequenza viene amplificato di un fattore proporzionale a $1/T_s$. Quindi diminuire $T_s$ migliora la risoluzione temporale, ma puo` peggiorare in modo netto la stima numerica delle derivate se la misura non e` sufficientemente pulita.

\subsection{Quantizzazione e conversione A/D}

\slidepair{images/page-16.png}{images/page-17.png}{Quantizzazione del segnale e definizione dell’LSB in un convertitore A/D.}

Dopo il campionamento, il segnale subisce una seconda discretizzazione: la quantizzazione in ampiezza. Un convertitore A/D a $N$ bit divide l’intervallo di misura in $2^N$ livelli. Se lo span è $R$, il passo di quantizzazione vale
\[
q = \frac{R}{2^N}.
\]
Questo valore coincide con l’LSB, cioè la minima variazione idealmente distinguibile dal codice digitale.

\slidefig{images/page-18.png}{Errore di quantizzazione e intervallo di incertezza associato al codice digitale.}

Se si usa l’arrotondamento al livello più vicino, l’errore di quantizzazione appartiene all’intervallo
\[
e_q \in \left[-\frac{q}{2},\frac{q}{2}\right].
\]
La quantizzazione è quindi una non linearità intrinseca del sistema di misura digitale. In molti casi è piccola rispetto al rumore analogico, ma in altri casi, soprattutto per segnali piccoli o lentamente variabili, può diventare significativa.

Per un ADC ideale, quando l’errore di quantizzazione puo` essere approssimato come rumore uniforme indipendente dal segnale, la varianza vale
\[
\sigma_q^2=\frac{q^2}{12}.
\]
Da qui discende anche la classica stima del rapporto segnale-rumore di quantizzazione per una sinusoide a fondo scala:
\[
\text{SNR}_{\text{ideal}} \approx 6.02\,N + 1.76\ \text{dB}.
\]
Questa formula e` molto utile per capire subito che aumentare di un bit la risoluzione nominale equivale, idealmente, a guadagnare circa \(6\) dB di SNR.

\slidepair{images/page-19.png}{images/page-20.png}{Esempio numerico di quantizzazione e interpretazione della risoluzione nominale dell’ADC.}

L’esempio numerico mostrato nelle ultime slide del blocco, con intervallo $0$-$10$ V e ADC a $12$ bit, dà
\[
q=\frac{10}{4096}\approx 2.44\ \text{mV},
\qquad
|e_q|_{\max}\approx 1.22\ \text{mV}.
\]
Questo mostra bene che la risoluzione numerica nominale dell’ADC non coincide automaticamente con l’accuratezza reale del sistema di misura.

Infatti la risoluzione nominale e` soltanto la granularita` del codice digitale. L’accuratezza effettiva dipende anche da offset, non linearita` integrale e differenziale, rumore d’ingresso, jitter di campionamento e stabilita` della tensione di riferimento. Nei convertitori reali, per questo, si usa spesso anche il concetto di \emph{ENOB} (effective number of bits), che traduce tutte le non idealita` in un numero di bit effettivamente utilizzabili.

% =========================================================
\subsection{Caratteristica statica ideale e non linearità}

\slidepair{images/page-21.png}{images/page-22.png}{Caratteristica statica ideale del sensore e prime deviazioni dal comportamento lineare.}

La caratteristica statica ideale di un sensore è la relazione tra grandezza vera e grandezza misurata:
\[
y_{\text{measured}} = f(y_{\text{true}}).
\]
Nel caso ideale questa relazione sarebbe perfettamente nota, lineare e invertibile. Nella realtà compaiono invece non linearità del trasduttore, offset, errori di scala ed effetti introdotti dall’elettronica di condizionamento.

\slidepair{images/page-23.png}{images/page-24.png}{Non linearità, errori di scala e scostamenti sistematici dalla caratteristica ideale.}

\slidepair{images/page-25.png}{images/page-26.png}{Origine fisica ed elettronica delle non idealità: esempio di sensore calibrato in fabbrica.}

Le slide centrali del blocco mostrano correttamente che queste non idealità possono nascere sia dal principio fisico del sensore sia dai circuiti che lo accompagnano. L’esempio del BMP085, riportato subito dopo, è utile perché mostra che molti sensori moderni non producono direttamente una misura finale, ma dati grezzi che vengono corretti con coefficienti di calibrazione memorizzati in fabbrica. In forma sintetica, un modello utile è
\[
y = f(x) + b + n,
\]
dove $f$ è la caratteristica nominale, $b$ un bias e $n$ il rumore residuo.

Un parametro molto utile per leggere localmente la caratteristica e` la sensibilita`
\[
S(x)=\frac{dy}{dx}.
\]
Nel caso ideale affine $S(x)$ e` costante; in presenza di non linearita`, invece, la sensibilita` cambia con il punto di lavoro. Questo ha un impatto diretto sul controllo e sulla stima, perche' la stessa variazione fisica puo` produrre variazioni di uscita diverse in regioni diverse del range operativo.

Nelle specifiche industriali la non linearita` viene spesso espressa come massimo scostamento dalla retta ideale o dalla migliore retta di fit, normalizzato al fondo scala:
\[
\varepsilon_{NL} = 100\,\frac{\max_x |y(x)-y_{\text{lin}}(x)|}{FS}.
\]
Questa forma e` comoda per confrontare sensori diversi, ma va sempre interpretata ricordando quale retta di riferimento e` stata scelta.

\subsection{Isteresi}

\slidepair{images/page-27.png}{images/page-28.png}{Isteresi: dipendenza dell’uscita dalla storia del segnale di ingresso.}

L’isteresi è una non idealità diversa dalla semplice non linearità. In presenza di isteresi, l’uscita del sensore dipende non solo dal valore attuale dell’ingresso, ma anche dalla storia precedente del segnale. Per uno stesso valore del misurando si possono quindi ottenere uscite diverse a seconda che l’ingresso stia crescendo o diminuendo.

\slidepair{images/page-29.png}{images/page-30.png}{Ciclo di andata e ritorno e implicazioni pratiche dell’isteresi in robotica.}

Le cause più comuni sono attrito, backlash, remanenza magnetica, deformazioni elastiche e tensioni interne nei materiali. In pratica questo significa che il sensore non è più descritto da una funzione statica univoca $y=f(x)$, ma da un ciclo di andata e ritorno. In robotica il fenomeno è particolarmente rilevante quando il sistema compie movimenti ripetitivi con inversioni di direzione, perché l’errore non dipende solo dalla posizione, ma dal percorso seguito per raggiungerla.

Una misura sintetica dell’isteresi e` spesso la massima distanza tra il ramo crescente e quello decrescente:
\[
H_{\max}=\max_x |y_{\uparrow}(x)-y_{\downarrow}(x)|.
\]
Dal punto di vista del modello, questa formula fa vedere bene che l’isteresi non e` riducibile a un semplice offset costante: e` una memoria del sistema, e quindi una non idealita` intrinsecamente piu` difficile da compensare.

\subsection{Dead-band e saturazione}

\slidefig{images/page-31.png}{Dead-band e saturazione come non linearità statiche del sensore.}

La dead-band è l’intervallo di ingresso attorno a zero in cui il sensore non produce una variazione rilevabile in uscita. La saturazione è invece la perdita di sensibilità vicino ai limiti del fondo scala, dove ingressi diversi generano praticamente la stessa lettura. Sono entrambe non linearità statiche molto comuni.

La dead-band è critica quando si vogliono rilevare piccole velocità, piccole forze o piccoli spostamenti. La saturazione è critica perché tronca l’informazione disponibile e può indurre il controllore a sottostimare la vera ampiezza della grandezza fisica. In entrambi i casi, si perde qualità della misura in una porzione del range operativo.

In forma estremamente compatta si possono modellare come
\[
y(x)\approx
\begin{cases}
0 & |x|<x_d \quad \text{(dead-band)}\\
\operatorname{sat}(Kx) & \text{(saturazione)}
\end{cases}
\]
dove $\operatorname{sat}(\cdot)$ rappresenta il classico troncamento ai limiti di uscita. Questo tipo di modello a tratti e` molto semplice, ma spesso sufficiente per capire perche' il sensore smette di essere utile in certe zone del proprio range.

\subsection{Schema sintetico delle non idealit\`a di misura}

Per fissare le idee, conviene raccogliere in uno schema unico le principali non idealit\`a incontrate nelle prime lezioni. Questa vista d’insieme \`e utile perch\'e aiuta a distinguere rapidamente errori correggibili in calibrazione, errori dipendenti dalla storia del moto ed errori stocastici che richiedono modellazione statistica.

{\renewcommand{\arraystretch}{1.18}
\begin{longtable}{@{}>{\raggedright\arraybackslash}p{0.16\linewidth}>{\raggedright\arraybackslash}p{0.24\linewidth}>{\raggedright\arraybackslash}p{0.22\linewidth}>{\raggedright\arraybackslash}p{0.28\linewidth}@{}}
\toprule
\textbf{Effetto} & \textbf{Significato fisico} & \textbf{Modello tipico} & \textbf{Impatto pratico} \\
\midrule
\endfirsthead
\toprule
\textbf{Effetto} & \textbf{Significato fisico} & \textbf{Modello tipico} & \textbf{Impatto pratico} \\
\midrule
\endhead
Bias o offset & Scostamento sistematico quasi costante anche a ingresso nullo & $y=ax+b$ con $b\neq 0$ & Produce errore medio persistente; si corregge con zeroing o calibrazione \\
Errore di scala & Guadagno diverso da quello ideale & $y=ax$ con $a\neq 1$ & La misura cresce troppo o troppo poco rispetto al valore vero \\
Non linearit\`a & Relazione ingresso-uscita non affine & $y=f(x)$ non lineare & Un solo fattore di scala non basta; servono modelli pi\`u ricchi o linearizzazioni locali \\
Isteresi & Uscita dipendente dalla storia del segnale & $y=f_{\uparrow}(x)$ o $y=f_{\downarrow}(x)$ & Lo stesso ingresso pu\`o produrre letture diverse in salita e in discesa \\
Dead-band & Piccole variazioni del misurando non vengono rilevate & $y\approx 0$ per $|x|<x_d$ & Perdita di sensibilit\`a vicino a zero; critica per micro-spostamenti e basse velocit\`a \\
Saturazione & L’uscita non cresce oltre una certa soglia & $y\approx y_{\max}$ oltre il fondo scala & Il controllore sottostima il valore reale quando il sensore esce dal range utile \\
Rumore & Componente aleatoria campione per campione & $y=h(x)+n$ & Peggiora stima, derivazione numerica e stabilit\`a della retroazione \\
Deriva lenta & Variazione lenta di bias o parametri nel tempo & $b_{k+1}=b_k+w_k$ & Richiede ricalibrazione, compensazione termica o stima online \\
\bottomrule
\end{longtable}
}

% =========================================================
\subsection{Rumore e sorgenti fisiche}

\slidepair{images/page-32.png}{images/page-33.png}{Rumore di misura e principali sorgenti fisiche che lo generano.}

Le slide introducono il rumore come componente indesiderata che si somma alla misura e ne limita l’utilizzabilità in controllo e stima. Un modello semplice è
\[
y_k = h(x_k) + b_k + n_k,
\]
dove $n_k$ rappresenta la parte rumorosa e $b_k$ un eventuale bias o deriva lenta.

Le sorgenti fisiche richiamate nelle slide sono quelle più comuni: rumore termico, shot noise, flicker noise, rumore ambientale e rumore di quantizzazione. Il punto importante è che non tutte hanno la stessa struttura spettrale. Alcune sono quasi bianche su una certa banda, altre dominano alle basse frequenze, altre dipendono soprattutto dall’ambiente operativo.

Per questa ragione il rumore non va letto solo come ``dispersione'' attorno al valore medio, ma come processo temporale con correlazione e contenuto spettrale propri. Se si indica con
\[
R_n[\ell]=\mathbb{E}\{n[k]\,n[k-\ell]\}
\]
l’autocorrelazione del rumore, allora rumori diversi possono avere la stessa varianza $R_n[0]$ ma effetti molto diversi nel tempo e in frequenza.

\subsection{Rumore di quantizzazione, bianco, gaussiano e colorato}

\slidepair{images/page-34.png}{images/page-35.png}{Rumore di quantizzazione e modello uniforme equivalente.}

Le prime due slide del blocco riguardano il rumore di quantizzazione, che può essere modellato, in prima approssimazione, come una variabile uniforme
\[
n_q \sim \mathcal U\left(-\frac{\Delta}{2},\frac{\Delta}{2}\right),
\qquad
\sigma_q^2=\frac{\Delta^2}{12}.
\]
Questa approssimazione è utile per stimare l’ordine di grandezza dell’errore introdotto dall’ADC.

\slidepair{images/page-36.png}{images/page-37.png}{Rumore bianco, rumore gaussiano, rumore colorato e deriva lenta del bias.}

Le slide successive chiariscono poi una distinzione importante: ``bianco'' si riferisce allo spettro, mentre ``gaussiano'' si riferisce alla distribuzione dei campioni. Un rumore può quindi essere gaussiano ma non bianco, o bianco ma non gaussiano. Questa distinzione è fondamentale nei problemi di filtraggio e stima.

Nel caso di rumore bianco ideale, la densita` spettrale di potenza e` piatta:
\[
S_n(\omega)=S_0.
\]
Nel caso gaussiano, invece, l’informazione riguarda la distribuzione dei campioni:
\[
n \sim \mathcal N(\mu,\sigma^2).
\]
Questa separazione concettuale e` cruciale, perche' molti algoritmi di stima assumono simultaneamente gaussianita` e bianchezza, ma sono ipotesi diverse e vanno controllate separatamente sui dati reali.

Infine vengono introdotti modelli di rumore colorato e di deriva del bias. Un modello semplice di rumore colorato è
\[
n_{k+1}=\alpha n_k + v_k,
\]
mentre per il bias drift si usa spesso il random walk:
\[
b_{k+1}=b_k+w_k.
\]
Questi modelli sono particolarmente rilevanti nelle IMU e in molti sensori inerziali.

\subsection{Istogramma e spettro del rumore}

\slidepair{images/page-38.png}{images/page-39.png}{Istogramma dei campioni e spettro del rumore come strumenti complementari di analisi.}

Le due slide finali del blocco mostrano due strumenti complementari per caratterizzare il rumore. L’istogramma consente di osservare la distribuzione statistica dei campioni e di verificare se un modello gaussiano o uniforme sia plausibile. Lo spettro o la densità spettrale di potenza mostrano invece come la potenza del rumore sia distribuita sulle frequenze.

Questa distinzione è importante perché due processi con la stessa varianza possono avere effetti molto diversi se uno è concentrato alle alte frequenze e l’altro alle basse. Per questo, in pratica, la sola deviazione standard non basta quasi mai a caratterizzare in modo completo la qualità di un sensore.

In termini matematici, istogramma e spettro sono collegati dall’idea che distribuzione e struttura temporale sono due aspetti distinti dello stesso processo. L’istogramma risponde alla domanda ``quali valori assume il rumore?'', mentre lo spettro risponde alla domanda ``con quale velocita` e su quali frequenze varia?''. Un’analisi seria dei dati sperimentali richiede quasi sempre entrambe.

% =========================================================
\subsection{Calibrazione affine di base}

\slidepair{images/page-40.png}{images/page-41.png}{Impostazione della calibrazione affine e significato dei parametri di scala e offset.}

Questo blocco introduce il problema della calibrazione nel suo caso più semplice. L’idea è stimare una relazione tra ingresso fisico noto $x$ e uscita misurata $y$. Il modello base è affine:
\[
y = ax + b,
\]
dove $a$ rappresenta il fattore di scala e $b$ il bias o offset.

\slidefig{images/page-42.png}{Stima dei parametri del modello affine tramite minimi quadrati.}

La soluzione proposta è quella classica dei minimi quadrati. Organizzando i dati in forma matriciale,
\[
X=
\begin{bmatrix}
x_1 & 1\\
x_2 & 1\\
\vdots & \vdots\\
x_N & 1
\end{bmatrix},
\qquad
y=
\begin{bmatrix}
y_1\\y_2\\ \vdots\\ y_N
\end{bmatrix},
\]
si ottiene
\[
\hat\theta = (X^T X)^{-1}X^T y,
\qquad
\theta=
\begin{bmatrix}
a\\b
\end{bmatrix}.
\]

Se poi l’obiettivo non e` soltanto prevedere l’uscita del sensore ma correggere la misura, spesso serve anche il modello inverso. Nel caso affine, assumendo $a\neq 0$, si ottiene
\[
\hat x = \frac{y-b}{a}.
\]
Questo passaggio e` banale solo in apparenza: in pratica significa decidere se la calibrazione venga usata nel verso diretto per simulare il sensore oppure nel verso inverso per compensarlo.

\slidetriple{images/page-43.png}{images/page-44.png}{images/page-45.png}{Dai dati sperimentali al modello calibrato: interpretazione grafica della stima.}

Il valore di questo blocco non sta nella complessità matematica, ma nel fatto che presenta la calibrazione come un problema di identificazione parametrica. Le ultime slide rendono proprio visibile questo passaggio: dai punti sperimentali grezzi si arriva a un modello stimato che poi potrà essere invertito o compensato.

Se si assume rumore bianco a varianza costante, la covarianza dei parametri affini stimati e` inoltre
\[
\operatorname{Cov}(\hat\theta)=\sigma^2 (X^TX)^{-1}.
\]
Anche nel caso semplice affine compare quindi gia` l’idea, molto importante per il seguito del corso, che il modello non sia descritto solo dai suoi parametri nominali, ma anche dall’incertezza con cui quei parametri sono stati identificati.

\subsection{Calibrazione di IMU e lettura della noise density}

\slidepair{images/page-46.png}{images/page-47.png}{Calibrazione di una IMU e modello matriciale della misura.}

Le ultime slide delle prime due lezioni collegano la calibrazione a un caso robotico concreto: le IMU. Il modello mostrato è
\[
y = Sx + b + n,
\]
dove $S$ è una matrice di scala e accoppiamento, $b$ è il bias e $n$ il rumore. È già un modello più realistico del caso scalare e prepara il terreno alla lezione successiva, in cui la calibrazione verrà trattata in modo più sistematico.

\slidefig{images/page-48.png}{Introduzione alla noise density nei datasheet dei sensori inerziali.}

La slide successiva apre invece il tema della \emph{noise density} nei datasheet, che è spesso una delle voci meno intuitive da leggere ma una delle più utili quando si vuole tradurre una specifica commerciale in un modello di rumore utilizzabile.

\slidepair{images/page-49.png}{images/page-50.png}{Dalla noise density alla varianza del rumore su una banda finita.}

Le ultime due slide esplicitano il passaggio quantitativo. Se un accelerometro dichiara, per esempio,
\[
100\ \mu g/\sqrt{\text{Hz}},
\]
questo valore va prima convertito in unità SI e poi interpretato come densità spettrale di ampiezza del rumore. Elevandolo al quadrato si ottiene la densità spettrale di potenza. Su una banda finita $B$, la varianza si ricava come
\[
\sigma_a^2 = S_a B.
\]
Questo passaggio è molto utile in pratica, perché permette di tradurre un dato di targa in una stima concreta della deviazione standard del rumore sul segnale.

Nel caso tipico di acquisizione digitale con filtro anti-alias e banda approssimativamente pari a $B\simeq f_s/2$, la deviazione standard discreta si puo` stimare direttamente come
\[
\sigma_a \approx N_a \sqrt{\frac{f_s}{2}},
\]
dove $N_a$ e` la noise density espressa in unita` fisiche per radice di hertz. Questa formula e` molto pratica in laboratorio, perche' consente di passare quasi immediatamente dal datasheet a un modello di rumore da inserire in simulazione, filtraggio o fusione sensoriale.

Nel caso delle IMU triassiali, inoltre, il modello matriciale $y=Sx+b+n$ va letto in modo ancora piu` ricco: la matrice $S$ non descrive solo i fattori di scala sui singoli assi, ma anche i disallineamenti geometrici e il cross-coupling tra canali. Proprio per questo la calibrazione di una IMU non e` una semplice taratura asse per asse, ma un problema realmente multivariabile.

\subsection*{Sintesi finale della Lezione 1-2}

Le prime due lezioni costruiscono il quadro di base necessario per tutto il corso. Viene chiarito che la misura non è mai un atto ideale: passa attraverso trasduzione, campionamento, quantizzazione, non idealità statiche, rumore e calibrazione. Dal punto di vista robotico, il messaggio è semplice: prima di progettare bene il controllo, bisogna capire bene come nasce e come si degrada la misura.

% =========================================================
% LEZIONE 3
% =========================================================

\clearpage
\section{Lezione 3}

\subsection{Visione generale della lezione}

Questa lezione sviluppa in modo più sistematico il tema della calibrazione dei sensori. L’idea centrale è che un sensore non va descritto solo con una curva nominale, ma con un modello completo che separi tre contributi distinti: la parte deterministica della caratteristica, un eventuale bias residuo e il rumore di misura. In questo modo la calibrazione diventa uno strumento utile non solo per correggere la lettura del sensore, ma anche per costruire modelli coerenti da usare in stima, filtraggio e fusione sensoriale.

\begin{focusbox}
\textbf{Idea chiave.} Calibrare un sensore significa identificare un modello di misura sufficientemente accurato da poter essere usato in modo affidabile dentro un sistema di controllo o di stima.
\end{focusbox}

\subsection{Modelli polinomiali di calibrazione}

Le prime slide introducono il caso di modello polinomiale. Il primo esempio è il modello quadratico
\[
y = a_2 x^2 + a_1 x + a_0,
\]
che è non lineare rispetto alla variabile fisica $x$, ma lineare nei parametri $a_2$, $a_1$, $a_0$. Questo consente di usare ancora i minimi quadrati classici.

\slidefig{images2/page-1.png}{Modello di calibrazione quadratico}

Per un insieme di campioni $(x_i,y_i)$ si può scrivere
\[
y_i = a_2 x_i^2 + a_1 x_i + a_0 + n_i,
\]
e raccogliere i dati nella forma matriciale
\[
\mathbf y = \Phi \theta + \mathbf n,
\qquad
\theta =
\begin{bmatrix}
a_2\\ a_1\\ a_0
\end{bmatrix},
\qquad
\Phi =
\begin{bmatrix}
x_1^2 & x_1 & 1\\
x_2^2 & x_2 & 1\\
\vdots & \vdots & \vdots\\
x_N^2 & x_N & 1
\end{bmatrix}.
\]

La slide successiva generalizza il ragionamento al polinomio di ordine $n$:
\[
y = \sum_{k=0}^{n} a_k x^k.
\]

\slidefig{images2/page-2.png}{Modello polinomiale generico}

Anche in questo caso il problema resta lineare nei parametri, quindi la stima si può ottenere con
\[
\hat \theta = (\Phi^T \Phi)^{-1}\Phi^T \mathbf y,
\]
purché i dati siano abbastanza informativi da rendere invertibile $\Phi^T\Phi$.

Qui compare pero` gia` un primo problema numerico serio: all’aumentare dell’ordine del polinomio, la matrice di regressione assume struttura di Vandermonde e puo` diventare mal condizionata. In pratica, campioni raccolti su un intervallo ampio o non ben distribuiti rendono i coefficienti di ordine alto molto sensibili al rumore. Per questo, anche quando il modello fisico suggerisce un polinomio, conviene spesso riscalare la variabile $x$ o usare basi numericamente piu` stabili.

\subsection{Scelta dell’ordine del modello}

Una delle slide più importanti è quella sulla scelta dell’ordine del polinomio. Il principio corretto è scegliere il modello più semplice possibile che lasci nei residui soltanto una componente compatibile con rumore casuale non strutturato.

\slidefig{images2/page-3.png}{Scelta dell’ordine del modello}

Se il modello è troppo semplice, parte della fisica del sensore resta nei residui. Se il modello è troppo ricco, si finisce per adattare anche il rumore sperimentale. In termini pratici, il residuo
\[
e_i = y_i - f(x_i,\hat\theta)
\]
dovrebbe avere media circa nulla e non mostrare andamento sistematico rispetto a $x$ o nel tempo.

Da un punto di vista piu` formale, la scelta dell’ordine va letta come compromesso tra \emph{bias di modello} e \emph{varianza della stima}. Un modello troppo povero introduce errore sistematico; un modello troppo ricco rende i parametri instabili e peggiora la capacita` di generalizzare su dati nuovi. E` la stessa logica che in identificazione di sistema si traduce nella distinzione tra adattamento ai dati di stima e capacita` predittiva sul set di validazione.

\begin{keybox}
\textbf{Regola pratica.} L’ordine del modello non va scelto per “far passare la curva dentro tutti i punti”, ma per descrivere bene la fisica del sensore senza inseguire il rumore.
\end{keybox}

\subsection{Incertezza dei parametri stimati}

Le slide successive introducono il tema dell’incertezza parametrica. Se il rumore di misura è modellato come gaussiano a media nulla, allora anche i parametri stimati sono variabili casuali, con covarianza
\[
\operatorname{Cov}(\hat\theta)=\sigma^2(\Phi^T\Phi)^{-1}.
\]

\slidefig{images2/page-4.png}{Intervalli di confidenza dei parametri}

La varianza del parametro $\hat\theta_i$ si legge sulla diagonale di questa matrice. Questo permette di costruire intervalli di confidenza e capire se i dati raccolti sono sufficientemente ricchi oppure no. Dal punto di vista ingegneristico, un parametro stimato con grande incertezza segnala che il dataset o il modello non sono ben bilanciati.

L’informazione parametrica si puo` poi propagare anche sull’uscita calibrata. Linearizzando il modello attorno ai parametri stimati, la varianza della previsione in un punto $x^\star$ si puo` approssimare con
\[
\operatorname{Var}\!\left(\hat y(x^\star)\right)\approx \phi(x^\star)^T \operatorname{Cov}(\hat\theta)\,\phi(x^\star),
\]
dove $\phi(x^\star)$ e` il vettore delle basi del modello. Questo passaggio e` molto utile quando si vuole sapere non solo quale sia la curva calibrata, ma con quanta confidenza si possa usare in un certo punto del range.

\subsection{Bias e rumore di misura}

Una parte essenziale della lezione è la separazione tra bias e rumore. Il modello completo diventa
\[
y = f(x,\theta) + b + n.
\]

\slidefig{images2/page-5.png}{Ruolo del bias nella misura}

Il bias rappresenta una deviazione sistematica o lentamente variabile. Il rumore rappresenta invece la parte aleatoria, che cambia rapidamente da campione a campione. Questa distinzione è fondamentale perché i due termini si trattano in modo diverso: il bias si corregge o si stima, il rumore si modella statisticamente e si attenua con medie o filtri.

Quando il bias non è costante, un modello semplice e molto usato è il random walk:
\[
b(k+1)=b(k)+w(k).
\]
Questo tipo di modello è centrale nelle IMU e nei giroscopi, dove la deriva lenta è spesso più critica del rumore bianco ad alta frequenza.

Dal punto di vista della stima, questa scrittura suggerisce subito un’interpretazione importante: il bias e` esso stesso uno stato dinamico del sistema di misura. Per questo, in molti filtri di Kalman e osservatori, il bias non viene trattato come una costante sconosciuta da compensare una volta per tutte, ma come una variabile addizionale da stimare online insieme alle grandezze fisiche di interesse.

\subsection{Stima del rumore residuo}

Una volta stimato il modello deterministico ed eventualmente il bias, si analizza il residuo:
\[
e_i = y_i - f(x_i,\hat\theta) - b.
\]

\slidefig{images2/page-6.png}{Modellazione del rumore residuo}

La varianza residua si può stimare con
\[
\hat\sigma^2 = \frac{1}{N-p}\sum_{i=1}^{N} e_i^2,
\]
dove $p$ è il numero di parametri stimati. Le verifiche utili sui residui sono poche ma molto importanti: media circa nulla, dispersione compatibile con il livello di rumore atteso, assenza di andamenti sistematici. Se i residui non soddisfano queste proprietà, significa che il modello è ancora incompleto oppure che il rumore non è ben descritto come bianco e indipendente.

Una verifica ancora piu` forte consiste nel controllare l’autocorrelazione residua:
\[
\hat R_e[\ell]=\frac{1}{N}\sum_{k=\ell+1}^{N} e_k e_{k-\ell}.
\]
Se il modello e` corretto e il rumore residuo e` davvero quasi bianco, allora $\hat R_e[\ell]$ dovrebbe essere piccola per tutti i lag $\ell\neq 0$. Quando invece compaiono strutture residue persistenti, significa che parte della dinamica o della non linearita` del sensore non e` ancora stata catturata dal modello scelto.

\subsection{Modello finale di calibrazione}

La lezione converge così su un modello sintetico ma molto potente:
\[
y = f(x,\theta) + b + n.
\]

\slidefig{images2/page-7.png}{Modello completo di calibrazione}

Questa forma è quella corretta da portare nelle applicazioni successive. La parte deterministica descrive la caratteristica del sensore, il bias rappresenta l’errore sistematico residuo, il rumore modella la variabilità stocastica.

Scritta in questo modo, la misura si presta anche a una scomposizione concettuale molto utile:
\[
\underbrace{f(x,\theta)}_{\text{fisica nominale}}
\;+\;
\underbrace{b}_{\text{errore sistematico}}
\;+\;
\underbrace{n}_{\text{errore casuale}}.
\]
La forza del modello sta proprio nel separare fenomeni che vanno trattati con strumenti diversi: identificazione parametrica per $f$, compensazione o stima lenta per $b$, modellazione statistica e filtraggio per $n$.

\subsection{Esempio implementativo}

L’ultima slide mostra un esempio Matlab/Simulink. Il punto non è tanto lo strumento software in sé, quanto il flusso concettuale: dati sperimentali, scelta del modello, stima dei parametri, analisi dei residui, uso finale del modello calibrato.

\slidefig{images2/page-8.png}{Esempio Matlab/Simulink}

Dal punto di vista implementativo, questa slide suggerisce anche un’idea metodologica importante: la calibrazione dovrebbe essere riproducibile e automatizzabile. Un buon script di calibrazione non serve solo a stimare i parametri una volta, ma a rendere tracciabile tutto il processo, inclusi i dataset usati, i criteri di validazione e le statistiche dei residui.

\subsection{Procedura pratica di calibrazione}

Per rendere operativa la lezione, conviene riassumere il flusso di calibrazione in pochi passi sempre uguali. Nella pratica di laboratorio, seguire una procedura fissa riduce il rischio di scegliere un modello troppo complesso o di interpretare male i residui.

\begin{focusbox}
\textbf{Flusso consigliato.}
\begin{enumerate}[leftmargin=1.5em]
\item Definire la grandezza fisica da calibrare, il range di lavoro e le condizioni operative rilevanti.
\item Acquisire coppie \((x_i,y_i)\) con un riferimento esterno affidabile, coprendo bene tutto il range utile.
\item Scegliere il modello pi\`u semplice compatibile con i dati: affine, polinomiale o modello fisico dedicato.
\item Stimare i parametri con minimi quadrati o con una procedura equivalente di identificazione.
\item Analizzare i residui \(e_i=y_i-f(x_i,\hat\theta)\): media quasi nulla, varianza plausibile, assenza di struttura sistematica.
\item Validare il modello su dati non usati in stima e solo dopo usarlo in controllo, stima o compensazione online.
\end{enumerate}
\end{focusbox}

In forma compatta, la logica del processo \`e
\[
\text{dati} \longrightarrow \text{modello} \longrightarrow \text{stima} \longrightarrow \text{residui} \longrightarrow \text{validazione}.
\]
Questa catena vale per sensori molto diversi tra loro: un potenziometro, un accelerometro MEMS o un sensore di forza cambiano per fisica, ma non per il metodo generale con cui la loro misura viene resa credibile.

\subsection*{Sintesi finale della Lezione 3}

La calibrazione non consiste solo nel trovare una curva che assomigli ai dati. Consiste nel costruire un modello di misura credibile, abbastanza semplice da essere robusto e abbastanza ricco da catturare il comportamento fisico del sensore. La distinzione tra caratteristica nominale, bias e rumore è il punto centrale da ricordare.

% =========================================================
% LEZIONE 4
% =========================================================

\clearpage
\section{Lezione 4}

\subsection{Visione generale della lezione}

Questa lezione si concentra su quattro temi collegati tra loro: caratteristiche dinamiche dei sensori, identificazione sperimentale della loro dinamica, distinzione tra accuratezza e precisione, e principali tecnologie per la misura di posizione. Dopo aver studiato la calibrazione statica, l’attenzione si sposta su come il sensore reagisce nel tempo e su come la sua tecnologia influenzi il dato usato dal controllore.

\begin{focusbox}
\textbf{Idea chiave.} Un sensore non va giudicato solo da quanto misura “vicino al valore vero”, ma anche da quanto rapidamente risponde e da quanto è affidabile la sua uscita nel tempo.
\end{focusbox}

\subsection{Caratteristiche dinamiche dei sensori}

La prima parte della lezione chiarisce che un sensore è un sistema dinamico. Tra grandezza fisica e uscita misurata esiste quindi una funzione di trasferimento:
\[
F(s)=\frac{\Delta y(s)}{\Delta u(s)}.
\]

\slidefig{images3/page-01.png}{Introduzione alle caratteristiche dinamiche dei sensori}

Questo significa che la misura non è istantanea. Se la dinamica del sensore è lenta rispetto alla banda del loop di controllo, il sensore introduce ritardo e perdita di fase nel ramo di retroazione. In alcuni casi questo è trascurabile. In altri casi modifica in modo sostanziale la risposta del sistema chiuso.

Per un sensore del primo ordine il legame tra ingresso e uscita puo` essere scritto anche nel dominio del tempo come
\[
\tau \dot y(t)+y(t)=K\,u(t),
\]
con risposta al gradino
\[
y(t)=K\,u_0\left(1-e^{-t/\tau}\right).
\]
Questa espressione rende immediatamente visibile il significato fisico di $\tau$: non e` un numero astratto, ma il tempo caratteristico con cui la misura rincorre il fenomeno reale.

L’esempio della termocoppia è particolarmente istruttivo. Il comportamento può essere approssimato con un primo ordine:
\[
\tau \frac{dT}{dt}+T=T_f.
\]

\slidefig{images3/page-02.png}{Esempio di termocoppia}

La costante di tempo $\tau$ dipende dalla capacità termica e dallo scambio con l’ambiente. Ne segue un compromesso chiaro: più il sensore è robusto e protetto, più tende a diventare lento.

\subsection{Dinamica del sensore dentro il loop di controllo}

Le slide successive mostrano l’effetto della dinamica del trasduttore su un sistema in retroazione. Se il sensore è ideale si può scrivere
\[
T(s)=k_T,
\]
mentre se il sensore ha dinamica propria si usa un modello come
\[
T(s)=\frac{k_T}{\tau_T s+1}.
\]

\slidepair{images3/page-03.png}{images3/page-04.png}{Definizione della caratteristica dinamica e confronto con il caso di sensore ideale.}

Le prime due slide servono a isolare il punto concettuale: il blocco di misura non è sempre un semplice guadagno, ma può introdurre una dinamica vera e propria tra variabile fisica e segnale restituito al controllore.

\slidepair{images3/page-05.png}{images3/page-06.png}{Sensore dinamico nel loop di controllo e conseguenze sulla risposta del sistema.}

Il messaggio è molto netto: anche se pianta e controllore restano invariati, un sensore più lento può peggiorare overshoot, oscillazioni e margine di stabilità. Questo accade perché la dinamica del sensore introduce ritardo di fase nel ramo di misura. La simulazione mostrata nell’ultima slide rende visibile proprio questa degradazione. Dal punto di vista progettuale, il sensore deve quindi essere incluso nel modello del loop già in fase di sintesi del controllo.

Se si indica con $P(s)$ la pianta, con $C(s)$ il controllore e con $H_s(s)$ il sensore, la funzione di trasferimento in anello chiuso diventa infatti
\[
T_{cl}(s)=\frac{P(s)C(s)}{1+P(s)C(s)H_s(s)}.
\]
Questa forma mostra bene che il sensore entra nel denominatore del sistema chiuso, e quindi influisce direttamente su poli, stabilita` e margini. Non e` quindi corretto ``aggiungere il sensore dopo'' come dettaglio di implementazione: va modellato gia` nel progetto.

\subsection{Identificazione sperimentale della dinamica}

Le slide propongono tre approcci classici per identificare sperimentalmente la funzione di trasferimento del sensore.

Il primo è il metodo della risposta al gradino. Per un modello del primo ordine
\[
G(s)=\frac{K}{\tau s+1},
\]
il guadagno statico $K$ si ricava dal valore finale, mentre la costante di tempo $\tau$ si ricava dal transitorio.

\slidefig{images3/page-07.png}{Risposta al gradino}

Il vantaggio del metodo è la semplicità: serve poca strumentazione e il comportamento del sensore è leggibile direttamente nel tempo. Il limite è che non sempre è facile applicare un vero gradino del misurando, e che la presenza di rumore o di dinamiche più complesse può rendere ambigua la stima.

Per un primo ordine puro, un fatto utile da ricordare è che il tempo in cui la risposta raggiunge il \(63.2\%\) del valore finale coincide con la costante di tempo \(\tau\). Questo rende il metodo del gradino estremamente rapido in laboratorio, purché il sensore sia ben approssimabile con un modello senza zeri o dinamiche multiple dominanti.

Il secondo è il metodo della risposta impulsiva.

\slidefig{images3/page-08.png}{Risposta impulsiva}

Dal punto di vista teorico è molto informativo, perché la risposta impulsiva contiene tutta la dinamica del sistema LTI. In pratica, però, è meno comodo del gradino: un impulso ideale non è fisicamente realizzabile e spesso bisogna usare eccitazioni brevi o colpi equivalenti.

In termini formali, la risposta impulsiva $h(t)$ e` il kernel di convoluzione che lega ingresso e uscita:
\[
y(t)=\int_{0}^{\infty} h(\tau)\,u(t-\tau)\,d\tau.
\]
Per questo e` la descrizione piu` completa della dinamica LTI. Tuttavia proprio la sua ricchezza teorica contrasta con la scarsa comodita` sperimentale, che spiega perche' in molti contesti industriali si preferiscano prove al gradino o sweep sinusoidali.

Il terzo è l’identificazione in frequenza. Si applicano ingressi sinusoidali di frequenza diversa e si ricavano modulo e fase della risposta:
\[
|G(j\omega)|=\frac{A_{out}}{A_{in}},
\qquad
\angle G(j\omega)=\phi.
\]

\slidefig{images3/page-09.png}{Identificazione in frequenza}

Questo approccio è spesso il più utile quando si vuole collegare direttamente la caratterizzazione del sensore al progetto del controllo, perché rende immediatamente visibili banda passante, attenuazione, ritardo di fase ed eventuali risonanze.

Per un primo ordine, ad esempio, si ha
\[
|G(j\omega)|=\frac{K}{\sqrt{1+(\omega\tau)^2}},
\qquad
\angle G(j\omega)=-\arctan(\omega\tau).
\]
La frequenza di taglio \(\omega_c=1/\tau\) e` il punto in cui il modulo scende di \(3\) dB rispetto al valore statico. Questa lettura frequenziale e` particolarmente utile quando si deve scegliere se il sensore e` abbastanza veloce rispetto alla banda desiderata del controllo.

\subsection{Accuratezza e precisione}

La parte successiva della lezione chiarisce la distinzione tra accuratezza e precisione. L’accuratezza misura quanto la lettura è vicina al valore vero. La precisione misura quanto le letture ripetute sono concentrate tra loro.

\slidefig{images3/page-10.png}{Accuratezza e precisione}

Questa distinzione è fondamentale perché i due problemi si affrontano in modo diverso. Un errore sistematico si corregge con calibrazione o compensazione. Una dispersione casuale si attenua con medie, filtri o strategie statistiche.

Le slide propongono anche espressioni quantitative dell’accuratezza, per esempio in percentuale di fondo scala oppure in percentuale del valore misurato:
\[
\varepsilon_f = 100\,\frac{x_m-x_v}{x_{FS}},
\qquad
\varepsilon_a = 100\,\frac{x_m-x_v}{x_v}.
\]

\slidepair{images3/page-11.png}{images3/page-12.png}{Accuratezza quantitativa e precisione come ripetibilit\`a della misura.}

Le ultime due slide completano il quadro con una lettura quantitativa. La precisione, cio\`e la ripetibilit\`a delle prove, \`e infatti legata alla dispersione delle misure ripetute. In pratica, un sensore pu\`o essere molto preciso ma poco accurato, oppure accurato in media ma poco preciso. Solo quando bias e dispersione sono entrambi piccoli si ha una misura veramente buona.

Se si raccolgono $N$ prove ripetute nelle stesse condizioni, la precisione si puo` quantificare con la deviazione standard campionaria
\[
s=\sqrt{\frac{1}{N-1}\sum_{i=1}^{N}(x_i-\bar x)^2}.
\]
L’accuratezza, invece, richiede il confronto con un riferimento o valore vero noto. Questa separazione e` decisiva in taratura: la precisione si valuta per ripetizione, l’accuratezza per confronto metrologico.

\subsection{Sensori di posizione}

L’ultima parte della lezione introduce una classificazione dei sensori di posizione basata su principio fisico, tipo di moto e metodo di codifica.

\slidefig{images3/page-13.png}{Introduzione ai sensori di posizione}

La prima slide ha il ruolo di aprire il panorama: la posizione può essere letta con tecnologie molto diverse, e ogni famiglia di sensori porta con sé compromessi specifici tra risoluzione, robustezza e facilità di interfacciamento.

\slidefig{images3/page-14.png}{Classificazione generale}

La seconda slide ordina queste famiglie in una classificazione generale. Le prime tecnologie approfondite in questa parte sono potenziometri ed encoder incrementali.

\subsubsection{Potenziometri}

Il potenziometro lineare resistivo realizza un partitore di tensione in cui la posizione si traduce direttamente in uscita analogica. Il principio ideale è
\[
e_o=e_i\frac{R_1}{R_1+R_2}\approx e_i\frac{x}{L}.
\]

\slidepair{images3/page-15.png}{images3/page-16.png}{Potenziometro lineare: realizzazione fisica e principio circuitale.}

Le due slide sul potenziometro lineare rendono bene il vantaggio principale di questa tecnologia: la posizione viene convertita subito in una tensione analogica leggibile senza logiche di conteggio. I vantaggi sono semplicità, basso costo e misura assoluta. I limiti sono usura, rumore di contatto e vita utile limitata.

In realta` anche il modello del partitore va letto con un minimo di attenzione circuitale: la relazione \(e_o\approx e_i x/L\) vale bene solo se il carico a valle ha impedenza molto alta rispetto alla resistenza del potenziometro. Se il circuito successivo assorbe corrente in modo non trascurabile, il partitore viene caricato e la caratteristica ingresso-uscita si deforma.

\slidetriple{images3/page-17.png}{images3/page-18.png}{images3/page-19.png}{Potenziometri rotativi: varianti costruttive e misura assoluta della posizione angolare.}

Le slide successive mostrano che le stesse idee valgono anche per il potenziometro rotativo, che misura una posizione angolare. Il vantaggio principale resta la misura assoluta immediata. Lo svantaggio principale resta la natura a contatto, che introduce usura e variabilità nel tempo.

Per questa ragione i potenziometri sono molto adatti a sistemi semplici, joint con escursione limitata, manopole o attuazioni poco gravose, mentre diventano meno attraenti quando servono lunga vita utile, elevata ciclicit\`a e bassa deriva nel tempo.

\subsubsection{Encoder incrementali}

Gli encoder incrementali usano invece una codifica digitale basata sul conteggio di impulsi. Un disco con tacche o fenditure, letto tramite sistema ottico, genera sequenze periodiche di livelli logici.

\slidefig{images3/page-20.png}{Principio di funzionamento degli encoder incrementali}

La prima slide introduce il principio fisico: la posizione non viene letta come tensione continua, ma trasformata in una sequenza di eventi discreti.

\slidefig{images3/page-21.png}{Canali e segnali dell’encoder incrementale}

La seconda slide rende esplicita la struttura dei segnali. I due canali A e B, sfasati di $90^\circ$, permettono di stabilire anche il verso del moto. Un eventuale canale Z fornisce un riferimento una volta per giro. La posizione viene ricostruita contando gli impulsi, quindi il sensore è incrementale e non assoluto. Se si perde il conteggio o si riavvia il sistema senza memoria della posizione, occorre una procedura di homing o una logica di riferimento iniziale.

I vantaggi principali sono alta risoluzione, robustezza e buona integrazione con sistemi digitali. Dal punto di vista robotico, sono tra i sensori di posizione più usati proprio perché combinano precisione, velocità e facilità di elaborazione.

Dal conteggio si ricava anche una stima di velocita`:
\[
\dot\theta[k]\approx \frac{\theta[k]-\theta[k-1]}{T_s}.
\]
Questa stima e` molto semplice da implementare, ma soffre di quantizzazione e rumore quando la velocita` e` bassa. In molti servoazionamenti si usano quindi finestre di conteggio variabili, filtraggio digitale o osservatori dedicati per ottenere una velocita` piu` pulita a partire dagli stessi canali dell’encoder.

\begin{keybox}
\textbf{Confronto sintetico.} Il potenziometro è semplice e assoluto ma soggetto a usura. L’encoder incrementale è preciso e veloce ma richiede conteggio e riferimento iniziale.
\end{keybox}

\subsection*{Sintesi finale della Lezione 4}

La quarta lezione introduce i principali criteri con cui si giudica un sensore di posizione: principio fisico, natura assoluta o incrementale della misura, qualità del dato nel tempo e facilità di integrazione nel controllo. Potenziometri ed encoder incrementali rappresentano due approcci molto diversi, ma entrambi mostrano bene che la misura di posizione non è mai un dato ``neutro'': dipende dalla tecnologia usata per produrlo.

% =========================================================
% LEZIONE 5
% =========================================================

\clearpage
\section{Lezione 5}

\subsection{Visione generale della lezione}

Questa lezione prosegue l’analisi dei sensori di posizione entrando più a fondo in tre famiglie molto importanti: encoder ottici avanzati, sensori differenziali induttivi e resolver. L’obiettivo è capire come aumentare la risoluzione, come ottenere una misura assoluta robusta e come trattare sensori che non forniscono direttamente una tensione continua, ma segnali modulati da demodulare e convertire.

\subsection{Encoder incrementali: struttura, decodifica e risoluzione}

\slidepair{images4/page-01.png}{images4/page-02.png}{Encoder incrementali: punti di forza, limiti e principali configurazioni costruttive.}

Le prime due slide del blocco servono soprattutto a mettere a fuoco il profilo pratico dell’encoder incrementale. I vantaggi principali sono alta risoluzione, robustezza e uscita nativamente digitale. Questo lo rende molto adatto ai sistemi di controllo numerico e ai servoazionamenti, perché il segnale entra quasi direttamente nella logica di conteggio del microcontrollore o dell’azionamento. I limiti, però, restano quelli tipici dei sensori incrementali: all’accensione non si conosce la posizione assoluta e, nel caso ottico, la qualità della lettura può degradare se polvere, sporco o disallineamenti riducono il contrasto tra zone trasparenti e opache.

Dal punto di vista costruttivo, la quadratura può essere ottenuta in modi diversi: con più piste sul disco oppure con più ricevitori spazialmente sfalsati. In entrambi i casi l’obiettivo è generare due segnali $A$ e $B$ sfasati di $90^\circ$ elettrici. Spesso è presente anche una maschera fissa che migliora il contrasto ottico e quindi la precisione della commutazione.

\slidefig{images4/page-03.png}{Decodifica dell’encoder incrementale.}

La slide di decodifica introduce il legame tra geometria del disco e conteggio digitale. Se il disco presenta $N_e$ tacche o finestre per giro, l’incremento meccanico elementare associato a un periodo del segnale è
\[
\Delta \theta_{\text{mecc}}=\frac{2\pi}{N_e}.
\]
Le piste $A$ e $B$ sono in quadratura proprio per permettere il riconoscimento del verso. Quando $A$ anticipa $B$ il conteggio evolve in un verso; quando $B$ anticipa $A$ il conteggio evolve nel verso opposto.

Da un punto di vista logico, la coppia \((A,B)\) attraversa una sequenza ciclica di quattro stati:
\[
00 \rightarrow 01 \rightarrow 11 \rightarrow 10 \rightarrow 00
\]
in un verso, e la sequenza inversa nel verso opposto. Questa struttura e` la base di tutta la decodifica in quadratura: non si contano solo fronti, ma si verifica anche che la transizione osservata sia coerente con il verso di moto.

\slidefig{images4/page-04.png}{Lettura dei segnali e logica di conteggio}

La slide successiva completa il quadro mostrando la logica concreta di lettura. L’eventuale canale $Z$ fornisce un impulso una volta per giro e viene usato come riferimento di zero o come evento di homing. La nozione di ``grado elettrico'' serve proprio a leggere questi diagrammi temporali: un periodo elettrico del segnale corrisponde a un singolo passo elementare del reticolo, quindi la variabile elettrica evolve molto più rapidamente dell’angolo meccanico. La logica di decodifica usa i fronti dei segnali e uno stato memorizzato in flip-flop o in un contatore up/down per capire se incrementare oppure decrementare il conteggio.

Questa lettura permette anche di rilevare condizioni anomale. Se, tra due campioni successivi, si osserva una transizione che salta due stati logici, il decoder puo` interpretarla come errore di campionamento, disturbo o velocita` troppo elevata rispetto alla frequenza di acquisizione. I decoder robusti, sia hardware sia software, usano proprio questo controllo di coerenza per segnalare eventuali perdita di impulsi.

\subsection{Codifiche X1, X2 e X4 degli encoder incrementali}

\slidetriple{images4/page-05.png}{images4/page-06.png}{images4/page-07.png}{Codifica X1: fronte utile, riconoscimento del verso ed evoluzione temporale dei segnali.}

Le prime tre slide di questo blocco mostrano in dettaglio come i fronti vengano trasformati in conteggi. In codifica X1 si utilizza un solo fronte per periodo, tipicamente il fronte di salita del canale $A$. Il livello istantaneo dell’altro canale viene campionato nello stesso istante per ricavare il verso di rotazione. In questo modo si ottiene una logica semplice e robusta, ma la risoluzione sfrutta solo una parte dell’informazione disponibile.

\slidepair{images4/page-08.png}{images4/page-09.png}{Implementazione della codifica X1 e comportamento della macchina a stati durante l’inversione del moto.}

Le sequenze temporali dedicate al cambio di verso chiariscono un punto importante per i sistemi reali: quando il moto si inverte, la sequenza degli stati logici percorre la stessa macchina a stati ma in ordine inverso. Per questo è naturale implementare la decodifica con automi a quattro stati oppure con contatori up/down dedicati.

Dal punto di vista dell’elettronica digitale, questa e` una macchina a stati finiti estremamente naturale: lo stato corrente memorizza l’ultima coppia \((A,B)\), l’ingresso e` la nuova coppia letta, e l’uscita e` l’azione \(\{-1,0,+1,\text{errore}\}\). Questa formulazione e` importante perche' separa bene il problema della misura del verso da quello del semplice conteggio di fronti.

\slidetriple{images4/page-10.png}{images4/page-11.png}{images4/page-12.png}{Confronto tra codifiche X1, X2 e X4 e loro effetto sulla risoluzione per giro.}

Le slide successive mettono a confronto i diversi livelli di sfruttamento dei fronti. Con la codifica X2 si usano salita e discesa di un singolo canale, mentre con la codifica X4 si sfruttano tutti i fronti di $A$ e $B$. Se il costruttore dichiara un encoder da $N_e$ PPR sul canale $A$, il numero di conteggi per giro diventa
\[
N_c=
\begin{cases}
N_e & \text{in codifica X1},\\
2N_e & \text{in codifica X2},\\
4N_e & \text{in codifica X4}.
\end{cases}
\]
Di conseguenza, la risoluzione angolare effettiva vale
\[
\Delta\theta=\frac{2\pi}{N_c}.
\]
La posizione ricostruita dal contatore si può quindi scrivere come
\[
\theta[k]=\frac{2\pi}{N_c}\,n[k],
\]
dove $n[k]$ è il conteggio accumulato al campione $k$.

In pratica, scegliere X1, X2 o X4 non e` solo una decisione di risoluzione, ma anche di robustezza. La codifica X4 sfrutta tutta l’informazione disponibile, ma diventa piu` sensibile a jitter, rimbalzi logici e limiti di banda dell’elettronica. Per questo gli azionamenti industriali permettono spesso di selezionare la modalita` di decodifica in funzione della velocita` massima attesa e della qualita` del segnale.

\slidefig{images4/page-13.png}{Decodifica con microcontrollore}

L’ultima slide del blocco sottolinea proprio l’aspetto implementativo. In un microcontrollore moderno, per esempio STM32, il timer può acquisire i canali $A$ e $B$, filtrare i fronti, riconoscere il verso e aggiornare il contatore senza caricare la CPU. Dal punto di vista del controllo robotico questo è molto utile, perché libera tempo di calcolo per il regolatore e riduce il rischio di perdere impulsi ad alta velocità.

Dal conteggio accumulato si possono poi derivare sia posizione sia velocita`. La posizione e` disponibile in modo diretto, mentre la velocita` richiede una differenziazione numerica oppure una misura del tempo tra impulsi. La prima soluzione e` migliore ad alta velocita`, la seconda a bassa velocita`; e` proprio da questo compromesso che nascono molte architetture di stima di velocita` usate negli azionamenti.

\subsection{Encoder ottici incrementali e encoder assoluti}

\slidepair{images4/page-14.png}{images4/page-15.png}{Encoder ottico incrementale ed encoder assoluto: principio fisico e differenza concettuale tra conteggio e codifica diretta.}

La prima slide mostra la realizzazione fisica dell’encoder ottico incrementale: disco codificato, emettitore, ricevitori e elettronica di shaping. Il passaggio concettuale verso l’encoder assoluto, mostrato subito dopo, consiste nel sostituire la misura per conteggio con una codifica diretta della posizione.

\slidepair{images4/page-16.png}{images4/page-17.png}{Lettura diretta del codice assoluto e risoluzione associata al numero di piste.}

Negli encoder assoluti il disco non ha più una o due sole piste periodiche, ma più corone concentriche. Ogni corona rappresenta un bit e ogni fotodetettore legge direttamente una cifra del codice. Se sono presenti $N$ piste indipendenti, il numero di posizioni distinguibili in un giro è $2^N$ e la risoluzione nominale vale
\[
\Delta\theta=\frac{2\pi}{2^N}.
\]
Per esempio, un encoder single-turn a $13$ bit distingue $8192$ posizioni per giro, con risoluzione di circa
\[
\Delta\theta=\frac{360^\circ}{8192}\approx 0.044^\circ.
\]

Il vantaggio decisivo è che il codice letto corrisponde direttamente alla posizione assoluta all’interno del giro. Non serve contare impulsi e non serve eseguire homing dopo l’accensione, almeno nella versione single-turn. Questo rende l’encoder assoluto molto comodo quando il robot deve conoscere subito la configurazione iniziale o quando un reset improvviso non può portare alla perdita dello stato.

Nel caso multi-turn il principio viene esteso aggiungendo una seconda informazione che conteggia i giri completi. Questa pu\`o essere ottenuta con ingranaggi codificati, memorie non volatili, sistemi magnetici ausiliari o piccoli accumulatori energetici che mantengono la traccia del numero di giri. Il vantaggio e` enorme in robotica articolata: dopo un power cycle il sistema puo` conoscere direttamente la configurazione dei giunti senza dover effettuare una nuova procedura di referenziazione.

\slidetriple{images4/page-18.png}{images4/page-19.png}{images4/page-20.png}{Problema delle transizioni simultanee nel codice binario e introduzione del codice Gray.}

Le slide dedicate al codice mostrano però anche un problema fondamentale del binario puro. Nel passaggio tra due settori adiacenti possono cambiare più bit simultaneamente; se la transizione fisica non è perfettamente simultanea, il lettore può osservare configurazioni spurie e quindi una posizione errata. Il caso classico è la transizione da $7$ a $8$, in cui molti bit cambiano insieme.

Per evitare questo problema si usa il codice Gray, in cui due parole adiacenti differiscono per un solo bit. In questo modo, mentre il sensore attraversa il bordo tra due settori, il massimo errore dovuto a transizioni non simultanee si riduce drasticamente. Il codice Gray non elimina i problemi costruttivi o di allineamento, ma rende la codifica molto più robusta dal punto di vista logico.

Dal punto di vista algoritmico, il codice Gray va poi riconvertito in binario se si vuole usare la parola letta per calcoli numerici ordinari. La conversione si ottiene con una somma XOR cumulativa tra bit piu` significativi e meno significativi. E` un dettaglio apparentemente secondario, ma fa capire bene che l’encoder assoluto non e` solo un dispositivo meccanico-ottico: e` anche una piccola architettura di codifica dell’informazione.

\slidepair{images4/page-21.png}{images4/page-22.png}{Esempio di encoder assoluto e riepilogo dei principali vantaggi e limiti applicativi.}

Il compromesso finale è chiaro: l’encoder assoluto fornisce direttamente la posizione, mantiene l’informazione anche dopo un’interruzione di alimentazione se l’architettura lo consente, ma costa di più e richiede una meccanica più precisa. Inoltre, quando si vuole conoscere anche il numero di giri completi, servono soluzioni multi-turn con elettronica aggiuntiva o memorie interne, che aumentano ulteriormente costo e complessità.

\subsection{Encoder ottici lineari con interpolazione}

\slidepair{images4/page-23.png}{images4/page-24.png}{Encoder lineare ottico: aumento della risoluzione sfruttando l’interpolazione analogica del segnale.}

Le prime due slide sull’interpolazione ottica introducono un’idea molto elegante: per aumentare la risoluzione non è necessario limitarsi a sogliare il segnale luminoso in modo binario. Si può invece sfruttare la modulazione continua dell’intensità luminosa dovuta al moto relativo tra reticoli ottici. In questo caso l’informazione di posizione non è contenuta solo nel numero di finestre attraversate, ma anche nella fase del segnale analogico generato dai fotodetettori.

\slidepair{images4/page-25.png}{images4/page-26.png}{Geometria del reticolo ottico e costruzione di segnali opposti per la lettura differenziale.}

La slide sui segnali sfasati di $180^\circ$ anticipa già il tema della lettura differenziale: si costruiscono coppie di segnali opposti per cancellare offset e disturbi comuni, migliorando il rapporto segnale-rumore.

Questo passaggio e` cruciale perche' l’interpolazione fine non lavora piu` su segnali digitali saturi, ma su forme analogiche quasi sinusoidali. In questo regime, piccoli offset DC, variazioni dell’intensita` luminosa o guadagni diversi tra canali introducono errori sistematici sulla fase stimata. La lettura differenziale nasce proprio per ridurre questi contributi prima ancora della stima dell’angolo.

\slidepair{images4/page-27.png}{images4/page-28.png}{Generazione di segnali in quadratura tramite sfasamenti di \(90^\circ\) e \(270^\circ\).}

\slidepair{images4/page-29.png}{images4/page-30.png}{Segnale differenziale nel periodo di scansione e costruzione della coppia seno-coseno.}

Con opportune geometrie dei reticoli si ottengono infatti due segnali in quadratura, approssimativamente sinusoidali:
\[
s_1(x)=A\sin\!\left(\frac{2\pi x}{p}\right),
\qquad
s_2(x)=A\cos\!\left(\frac{2\pi x}{p}\right),
\]
dove $p$ è il passo del reticolo. Le slide mostrano anche segnali sfasati di $180^\circ$, utili per costruire letture differenziali che cancellano offset e disturbi comuni. Indicando con $s_3=-s_1$ e $s_4=-s_2$ i segnali opposti, si possono formare
\[
d_s=s_1-s_3=2A\sin\!\left(\frac{2\pi x}{p}\right),
\qquad
d_c=s_2-s_4=2A\cos\!\left(\frac{2\pi x}{p}\right).
\]

\slidepair{images4/page-31.png}{images4/page-32.png}{Equazioni di interpolazione, ricostruzione della fase e riferimento assoluto di zero.}

Il vantaggio della coppia sinusoidale in quadratura, mostrato esplicitamente nelle ultime due slide del blocco, è che permette di ricavare la posizione frazionaria all’interno del singolo passo ottico usando la fase:
\[
x_{\text{frac}}=\frac{p}{2\pi}\operatorname{atan2}(d_s,d_c).
\]
In questo modo non ci si limita più a sapere in quale intervallo cade la posizione, ma si stima anche dove ci si trova all’interno di quell’intervallo. È questo il principio che consente di superare la risoluzione puramente geometrica del reticolo.

Per ottenere la posizione assoluta lungo tutta la corsa bisogna poi combinare la parte frazionaria con il conteggio dei periodi completi:
\[
x = kp + x_{\text{frac}},
\]
dove $k$ è il numero di passi interi attraversati. Le slide mostrano anche come un fotodetettore aggiuntivo o una griglia di riferimento possano fornire un punto di zero per resettare il conteggio di $k$.

Dal punto di vista del controllo, questa tecnologia è molto interessante perché combina alta risoluzione e natura continua del segnale. L’aspetto più delicato è il condizionamento analogico: ampiezze sbilanciate, offset, distorsioni o sfasamenti non esattamente ortogonali peggiorano la qualità dell’angolo calcolato con \(\operatorname{atan2}\) e quindi introducono errore metrico sulla posizione stimata.

In forma geometrica, questi difetti trasformano il locus ideale seno-coseno, che dovrebbe essere un cerchio nel piano \((d_c,d_s)\), in un’ellisse decentrata o ruotata. La correzione di guadagni, offset e non ortogonalita` puo` quindi essere letta come un problema di ``circolarizzazione'' della traiettoria misurata prima del calcolo della fase. Questo e` uno dei motivi per cui gli encoder interpolati di alta gamma includono elettroniche di compensazione molto sofisticate.

\subsection{Frange di Moir\'e e aumento della rilevabilit\`a}

\slidepair{images4/page-33.png}{images4/page-34.png}{Frange di Moir\'e: principio di battimento geometrico e aumento della rilevabilit\`a dello spostamento.}

Le prime due slide di questo blocco spiegano come migliorare la precisione anche a parità di passo nominale del reticolo. L’idea è sovrapporre due reticoli con passo quasi uguale. L’interferenza geometrica tra i due genera una figura di battimento, cioè le frange di Moir\'e, il cui periodo spaziale è molto più grande del passo originale.

Se i due passi sono $p_1$ e $p_2$, il periodo delle frange vale
\[
\Lambda=\frac{p_1p_2}{|p_1-p_2|}.
\]
Nel caso mostrato nelle slide, con $p_1=p$ e $p_2=p+\Delta p$, per $\Delta p$ piccolo si ottiene l’approssimazione
\[
\Lambda \approx \frac{p^2}{|\Delta p|},
\qquad
\Lambda \gg p.
\]
Questo significa che uno spostamento meccanico molto fine viene trasformato otticamente in una variazione più lenta e più facile da misurare con buon rapporto segnale-rumore.

\slidefig{images4/page-35.png}{Confronto tra interpolazione diretta e Moire}

Il punto importante è che il sensore non ``crea'' informazione dal nulla. La risoluzione resta legata alla qualità geometrica del reticolo e alla qualità dell’elettronica di lettura. Tuttavia il Moir\'e rende più osservabile la fase del moto e quindi migliora la robustezza dell’interpolazione, riducendo la sensibilità a imperfezioni locali, tolleranze di montaggio e rumore di misura.

Un modo utile per leggere quantitativamente il fenomeno e` introdurre il fattore di magnificazione geometrica
\[
M=\frac{\Lambda}{p}.
\]
Quando \(M\gg 1\), piccoli errori di allineamento o piccole traslazioni meccaniche producono variazioni piu` lente e piu` facilmente osservabili nel pattern Moire. E` per questo che queste strutture sono cosi` efficaci nei sistemi di misura di precisione: non aumentano la quantita` di informazione disponibile, ma la rendono molto piu` leggibile.

\subsection{LVDT: Linear Variable Differential Transformer}

\slidepair{images4/page-36.png}{images4/page-37.png}{LVDT: principio generale e struttura fisica del sensore differenziale lineare.}

Le prime due slide introducono l’LVDT come un sensore lineare basato su un principio da trasformatore: una bobina primaria viene eccitata in AC, due bobine secondarie raccolgono le tensioni indotte e la posizione del nucleo ferromagnetico mobile modifica l’accoppiamento magnetico. La traduzione più fedele dell’idea mostrata nelle figure è: nucleo centrato $\Rightarrow$ tensioni indotte uguali $\Rightarrow$ uscita differenziale nulla; nucleo spostato $\Rightarrow$ una tensione cresce, l’altra diminuisce $\Rightarrow$ l’uscita cambia in funzione dello spostamento.

\slidepair{images4/page-38.png}{images4/page-39.png}{Primario, secondari e costruzione dell’uscita differenziale nell’LVDT.}

Le due slide successive entrano nel cuore del principio di misura. Il primario viene eccitato con una sinusoide, indicata nella slide come
\[
V_e(t)=V_m\sin(\omega t).
\]
Questa scelta di notazione, con la pulsazione \(\omega\), viene mantenuta qui per restare coerenti con le equazioni riportate nelle slide stesse. Sui due secondari si inducono tensioni sinusoidali con la stessa pulsazione della portante, ma con ampiezza dipendente dalla posizione \(x\) del nucleo:
\[
V_1(t)=k_1(x)V_m\sin(\omega t),
\qquad
V_2(t)=k_2(x)V_m\sin(\omega t).
\]
Poiché i due secondari sono collegati in opposizione, il sensore espone direttamente la differenza
\[
V_o(t)=V_2(t)-V_1(t).
\]

\slidefig{images4/page-40.png}{Condizione di simmetria, zero elettrico e inversione di fase nell’LVDT.}

La slide sulla condizione di simmetria mostra il punto di lavoro fondamentale: per \(x=0\), in un circuito idealmente simmetrico, si ha \(V_1=V_2\) e quindi \(V_o=0\). Se il nucleo si sposta verso un secondario, l’accoppiamento magnetico di quel ramo aumenta mentre quello dell’altro diminuisce. In una linearizzazione locale attorno allo zero si può scrivere
\[
k_1(x)\approx k_0-k_x x,
\qquad
k_2(x)\approx k_0+k_x x,
\]
da cui segue
\[
V_o(t)\approx 2k_xV_m\,x\,\sin(\omega t).
\]
Il segno di \(x\) non è codificato solo nell’ampiezza, ma anche nella fase relativa rispetto alla portante: da un lato il segnale è in fase, dall’altro è in controfase.

Questo punto e` essenziale: l’LVDT non e` soltanto un sensore di ampiezza, ma un sensore a modulazione di ampiezza con informazione di segno trasportata dalla fase. Se si osservasse solo \(|V_o|\), si perderebbe il verso dello spostamento rispetto al centro; e` la demodulazione sincrona rispetto alla portante che consente di ricostruire una grandezza firmata proporzionale a \(x\).

\slidetriple{images4/page-41.png}{images4/page-42.png}{images4/page-43.png}{Forma del segnale di uscita dell’LVDT, esempio temporale e andamento dell’ampiezza rispetto alla posizione.}

Questo terzetto di slide collega il modello al segnale reale. L’LVDT non restituisce una tensione continua direttamente proporzionale a \(x\), ma una sinusoide modulata in ampiezza. La terza figura del gruppo è particolarmente utile perché mostra separatamente
\[
|V_o(x)| \approx S_x |x|,
\]
cioè un’ampiezza che cresce in modo pressoché lineare nel tratto centrale, e una fase che salta di \(\pi\) quando si attraversa lo zero. Questo è il motivo per cui l’elettronica di lettura deve recuperare sia modulo sia segno dello spostamento.

Fuori dalla regione centrale, pero`, l’andamento smette di essere ben approssimabile con una retta. Il campo magnetico disperso, le non uniformita` geometriche e la saturazione del materiale ferromagnetico rendono meno accurata l’ipotesi lineare. La ``zona utile'' dell’LVDT e` quindi il tratto in cui la differenza tra i due accoppiamenti resta quasi affine rispetto a \(x\).

\slidepair{images4/page-44.png}{images4/page-45.png}{Prestazioni tipiche dell’LVDT e schema di condizionamento per ottenere una misura digitale utilizzabile.}

La slide prestazionale riporta ordini di grandezza tipici utili da ricordare: risoluzione dell’ordine di \(2\)–\(20\,\mu\text{m}\) per sensori lineari, sensibilità circa \(50\)–\(100\ \text{mV}/(\text{mm}\cdot\text{V})\), errore di linearità attorno a \(0.1\%\)–\(0.5\%\) e campi di misura lineari tipicamente tra \(1\) e \(10\ \text{cm}\). Sono numeri coerenti con l’immagine applicativa dell’LVDT: sensore molto preciso e molto lineare, ma non pensato per corse estremamente grandi.

La slide sul condizionamento chiarisce bene un punto che spesso nei riassunti resta implicito: per determinare la posizione serve un circuito capace di stimare l’ampiezza della sinusoide e di rilevarne la fase rispetto alla portante. I circuiti di condizionamento per LVDT fanno esattamente questo, filtrando il segnale differenziale, estraendo modulo e segno e, nei convertitori LVDT-to-digital, campionando il risultato per portarlo in formato digitale.

Una catena tipica comprende: generatore della portante, amplificazione del differenziale, demodulatore sincrono, filtro passa-basso e stadio finale di scalatura. In molti casi il demodulatore viene implementato moltiplicando \(V_o(t)\) per una copia sincrona di \(\sin(\omega t)\) o \(\cos(\omega t)\); dopo il filtraggio resta una componente continua proporzionale a \(x\). E` esattamente lo stesso principio che riapparira` piu` avanti nel resolver-to-digital converter.

\slidepair{images4/page-46.png}{images4/page-47.png}{Vantaggi, limiti applicativi ed esempio costruttivo di sensore LVDT commerciale.}

L’ultima coppia di slide conclude con il classico compromesso applicativo: assenza di contatto, robustezza, ripetibilità elevata e buona immunità ambientale da un lato; necessità di eccitazione AC, demodulazione dedicata, corsa limitata e costo non trascurabile dall’altro. In pratica, l’LVDT si sceglie quando contano soprattutto linearità, affidabilità e stabilità della misura.

\subsection{RVDT: Rotary Variable Differential Transformer}

\slidetriple{images4/page-48.png}{images4/page-49.png}{images4/page-50.png}{RVDT: struttura rotativa, accoppiamento trasformatore ed equazioni di base della misura angolare.}

Le tre slide del blocco mostrano che l’RVDT è la controparte rotativa dell’LVDT. La struttura fisica cambia per adattarsi a una variabile angolare, ma il principio resta quello di un trasformatore differenziale: eccitazione AC sul primario, due uscite secondarie e lettura della differenza.

Seguendo esattamente la notazione riportata nelle figure, l’eccitazione primaria viene scritta come
\[
V_p(t)=V_0\sin(\omega t),
\]
mentre le due uscite secondarie valgono
\[
V_{s1}(t)=k_1(\theta)V_0\sin(\omega t),
\qquad
V_{s2}(t)=k_2(\theta)V_0\sin(\omega t),
\]
con uscita differenziale
\[
V_{out}(t)=V_{s1}(t)-V_{s2}(t).
\]

Se \(k_1(\theta)\) e \(k_2(\theta)\) sono quasi lineari e simmetriche nell’intorno della posizione centrale, allora \(V_{out}\) varia linearmente con \(\theta\) e si annulla per \(\theta=0\). Anche qui la fase di \(V_{out}\), pari a \(0^\circ\) o \(180^\circ\) rispetto a \(V_p(t)\), permette di capire il segno dello spostamento angolare. La slide riporta anche un ordine di grandezza della sensibilità, tipicamente \(1\)–\(5\ \text{mV}/(\text{deg}\cdot\text{V})\).

La differenza principale rispetto all’LVDT e` che la regione di buona linearita` angolare e` normalmente limitata a un intervallo relativamente ridotto attorno alla posizione centrale. L’RVDT non nasce quindi come sensore ``su tutto il giro'', ma come sensore robusto di piccoli spostamenti angolari. E` questa caratteristica che lo rende molto adatto a servomeccanismi, superfici mobili e attuazioni in cui conta soprattutto la robustezza in un intervallo limitato.

\slidefig{images4/page-51.png}{Vantaggi, applicazioni e limiti dell’RVDT}

La slide finale del blocco traduce in modo molto chiaro il profilo applicativo dell’RVDT. I vantaggi sono: assenza di contatto e quindi assenza di usura, robustezza meccanica, buona ripetibilità, tolleranza ad alte temperature e immunità a polvere e olio. Le limitazioni sono invece il range angolare limitato, la necessità di eccitazione AC e di elettronica di demodulazione, e una risoluzione intrinseca inferiore a quella dei migliori encoder ottici. Non a caso la slide conclude con un criterio semplice ma efficace: l’RVDT è la scelta giusta quando \emph{robustezza} conta più di \emph{risoluzione estrema}.

\subsection{Resolver}

\slidetriple{images4/page-52.png}{images4/page-53.png}{images4/page-54.png}{Resolver: principio di trasformatore rotante, accoppiamento magnetico e dipendenza seno-coseno dall’angolo.}

Le tre slide introduttive definiscono il resolver come \emph{rotating transformer}: un trasformatore rotante in cui un avvolgimento rotorico eccitato in AC induce sugli avvolgimenti statorici due tensioni sinusoidali e cosinusoidali. La differenza chiave rispetto all’RVDT è esplicitata chiaramente: la bobina di eccitazione è avvolta sul rotore, così da consentire la rotazione continua su \(360^\circ\).

Questa osservazione cambia completamente il profilo del sensore: il resolver non e` un trasduttore differenziale limitato a una piccola zona quasi lineare, ma una macchina elettromagnetica progettata per generare direttamente una codifica seno-coseno continua della posizione angolare. In questo senso e` molto piu` vicino a un ``encoder analogico assoluto'' che a un semplice sensore incrementale.

Nella forma più vicina alla slide introduttiva, le ampiezze utili sono
\[
V_s = kV_{ref}\sin\theta,
\qquad
V_c = kV_{ref}\cos\theta,
\]
e l’angolo può essere ricostruito idealmente da
\[
\theta=\tan^{-1}\!\left(\frac{V_s}{V_c}\right).
\]
Dal punto di vista implementativo, però, è più corretto usare \(\operatorname{atan2}(V_s,V_c)\), perché la funzione \(\tan^{-1}(V_s/V_c)\) da sola non distingue i quadranti.

Le figure sull’accoppiamento magnetico e sulla dipendenza sinusoidale dall’angolo spiegano il motivo fisico della coppia seno-coseno: la rotazione proietta il flusso magnetico su due assi ortogonali dello statore. È proprio questa ortogonalità a rendere il resolver molto robusto per la misura angolare continua, anche in ambienti dove polvere, vibrazioni o temperatura renderebbero più critica una tecnologia ottica.

Dal punto di vista geometrico, le due uscite ideali soddisfano
\[
V_s^2+V_c^2 = (kV_{ref})^2,
\]
cioe` descrivono un cerchio nel piano \((V_c,V_s)\). Anche qui, come negli encoder interpolati, molte procedure di correzione reale consistono nel riportare il segnale misurato il piu` possibile a questa struttura circolare ideale compensando offset, differenze di guadagno e sfasamenti non perfettamente ortogonali.

\slidepair{images4/page-55.png}{images4/page-56.png}{Resolver reale: segnali modulati, demodulazione e ricostruzione dell’angolo durante la rotazione.}

Le due slide successive chiariscono infatti che il resolver reale non restituisce direttamente \(V_s\) e \(V_c\) come grandezze continue già demodulate, ma due sinusoidi modulate dalla portante. Una delle convenzioni mostrate nelle slide è
\[
V_s(t)=kV_{ref}\sin\theta\,\cos(\omega t),
\qquad
V_c(t)=kV_{ref}\cos\theta\,\cos(\omega t).
\]
In altre slide dello stesso blocco compare invece \(\sin(\omega t)\) al posto di \(\cos(\omega t)\): non è una contraddizione fisica, ma solo una diversa convenzione di fase sulla portante di riferimento.

Dopo una demodulazione sincrona si recuperano i corrispondenti segnali in continua
\[
V_{s,DC}=K\sin\theta,
\qquad
V_{c,DC}=K\cos\theta,
\]
e a quel punto l’angolo si ricava in modo robusto con
\[
\theta=\operatorname{atan2}(V_{s,DC},V_{c,DC}).
\]
La slide 136 aiuta molto a visualizzare questo passaggio: durante la rotazione, il segnale di eccitazione viene modulato da \(\sin\theta\) e \(\cos\theta\), e il resolver converter ha il compito di estrarre proprio questi inviluppi.

In applicazioni di precisione si esegue spesso anche una normalizzazione dell’ampiezza prima del calcolo dell’angolo, per ridurre la sensibilita` a variazioni di guadagno della catena analogica:
\[
\tilde S = \frac{V_{s,DC}}{\sqrt{V_{s,DC}^2+V_{c,DC}^2}},
\qquad
\tilde C = \frac{V_{c,DC}}{\sqrt{V_{s,DC}^2+V_{c,DC}^2}}.
\]
In questo modo l’angolo dipende soprattutto dal rapporto tra i due canali e molto meno dall’ampiezza assoluta della portante o dai guadagni elettronici.

\subsection{Resolver-to-Digital Converter}

\slidepair{images4/page-57.png}{images4/page-58.png}{Resolver-to-Digital Converter: demodulazione sincrona e calcolo diretto dell’angolo con \(\operatorname{atan2}\).}

Il blocco finale della lezione spiega come i segnali del resolver vengano convertiti in una posizione digitale. La prima coppia di slide scrive esplicitamente
\[
V_s(t)=KV_{ref}\sin\theta\,\sin(\omega t),
\qquad
V_c(t)=KV_{ref}\cos\theta\,\sin(\omega t),
\]
e mostra che l’RDC moltiplica queste uscite per il segnale di riferimento \(\sin(\omega t)\). Usando l’identità
\[
\sin^2(\omega t)=\frac{1}{2}\bigl(1-\cos(2\omega t)\bigr),
\]
dopo filtraggio passa-basso spariscono i termini a \(2\omega\) e restano due componenti continue proporzionali a seno e coseno dell’angolo:
\[
V'_s \propto \sin\theta,
\qquad
V'_c \propto \cos\theta.
\]
Qui mantengo deliberatamente la notazione con \(\omega\), e non con \(\omega_c\), perché è quella usata nelle slide del blocco RDC.

La slide 138 formalizza allora il problema in modo molto pulito: dati
\[
S=A\sin\theta,
\qquad
C=A\cos\theta,
\]
bisogna calcolare \(\theta\). Il primo metodo è il calcolo diretto
\[
\theta=\operatorname{atan2}(S,C).
\]
La slide sottolinea bene vantaggi e limiti: soluzione concettualmente semplice, ma bisognosa di ADC accurati e sensibile a rumore, offset e squilibri di ampiezza.

Se si vuole ragionare in termini architetturali, il percorso del segnale e` il seguente:
\[
\text{resolver} \rightarrow \text{amplificazione} \rightarrow \text{demodulazione sincrona} \rightarrow \text{filtro} \rightarrow \text{ADC} \rightarrow \text{calcolo dell’angolo}.
\]
Ogni blocco puo` introdurre errore: la demodulazione puo` lasciare residui di portante, il filtro puo` introdurre ritardo, l’ADC puo` aggiungere quantizzazione, il calcolo digitale puo` essere sensibile a squilibri dei canali. L’RDC va quindi letto come una catena completa di misura, non come una singola formula \(\operatorname{atan2}\).

\slidepair{images4/page-59.png}{images4/page-60.png}{Errore di inseguimento, tracking servo loop e interpretazione del convertitore come PLL.}

Il secondo metodo, mostrato nella coppia di slide successiva, è il \emph{tracking servo loop}. Invece di calcolare subito \(\theta\), il convertitore genera una stima interna \(\hat\theta\), sintetizza i corrispondenti seno e coseno e confronta segnali misurati e segnali previsti. L’errore di inseguimento è definito nella slide come
\[
\varepsilon=S\cos(\hat\theta)-C\sin(\hat\theta).
\]
Sostituendo \(S\) e \(C\) si ottiene una quantità proporzionale a
\[
\varepsilon \propto \sin(\theta-\hat\theta)\approx \theta-\hat\theta
\]
per piccoli errori. La slide usa la proporzionalità proprio per evidenziare che l’ampiezza complessiva può essere assorbita nei guadagni del loop o normalizzata a parte.

La stima viene poi aggiornata con un controllore di tipo PI:
\[
\dot{\hat\theta}=K_p\varepsilon + K_i\int \varepsilon\,dt.
\]
La slide lo interpreta correttamente come un PLL: quando \(\varepsilon=0\), la stima diventa costante e si ha \(\hat\theta=\theta\). I vantaggi messi in evidenza sono rumore meglio rigettato, alta risoluzione, inseguimento continuo e disponibilità quasi diretta di una stima di velocità dall’anello interno.

Questa e` probabilmente l’idea piu` importante dell’intero schema RDC. Invece di stimare ogni campione indipendentemente, il convertitore costruisce un modello interno dell’angolo e corregge solo l’errore di inseguimento. In termini di teoria dei sistemi, il tracking loop agisce come un osservatore dinamico non lineare che sfrutta la continuita` temporale della traiettoria meccanica. E` proprio per questo che puo` essere molto piu` robusto al rumore rispetto al calcolo diretto istante per istante.

Se si linearizza per piccoli errori, il loop vede un errore quasi proporzionale a \(\theta-\hat\theta\). La scelta di \(K_p\) e \(K_i\) determina allora banda, rapidita` di aggancio e sensibilita` al rumore. Un loop molto veloce segue bene accelerazioni e velocita` elevate, ma lascia passare piu` rumore; un loop piu` lento filtra meglio ma puo` accumulare errore dinamico quando il moto varia rapidamente.

\slidetriple{images4/page-61.png}{images4/page-62.png}{images4/page-63.png}{Architettura concreta dell’RDC, campionamento sincrono, VCO e bilancio finale della tecnologia resolver.}

Le ultime tre slide rendono concreta l’architettura dell’RDC. La prima mostra un esempio di estrazione dell’inviluppo con campionamento sincrono: generatore di eccitazione, condizionamento analogico, circuito sample-and-hold sincronizzato, ADC e infine CPU o blocco RDC che ricostruisce posizione e velocità. La seconda entra ancora più nel dettaglio e indica la variabile digitale interna con \(\varphi\): il detector produce un errore proporzionale a
\[
K\sin(\theta-\varphi),
\]
che viene integrato e usato per pilotare un VCO, cioè un \emph{voltage-controlled oscillator}, il quale aggiorna il contatore up/down. In aggancio, la slide riporta la condizione
\[
\varphi = \theta \pm 1\,\text{LSB}.
\]
Questa è la traduzione circuitale del principio di tracking visto nella slide precedente.

Letto blocco per blocco, lo schema dell’RDC funziona cosi`:
\begin{enumerate}[leftmargin=1.5em]
\item il resolver genera due segnali AM portati dalla stessa eccitazione;
\item il demodulatore sincrono estrae le componenti lente proporzionali a seno e coseno dell’angolo;
\item il phase detector confronta tali componenti con seno e coseno sintetizzati dalla stima interna \(\varphi\);
\item il loop filter integra l’errore e ne ricava un comando lento e robusto;
\item il VCO o numerically controlled oscillator traduce questo comando in un incremento di fase;
\item il contatore up/down aggiorna la rappresentazione digitale dell’angolo.
\end{enumerate}

Questa lettura e` importante perche' spiega finalmente perche' in molti datasheet compaiono contemporaneamente termini come demodulator, phase detector, loop filter, VCO e counter: non sono blocchi scollegati, ma i componenti standard di una catena di tracking di fase applicata alla misura angolare.

La disponibilita` di una fase interna \(\varphi\) ha anche un altro vantaggio decisivo: la velocita` angolare si ottiene quasi gratuitamente come derivata della fase stimata o, nel linguaggio del loop, come frequenza del VCO in aggancio. E` per questo che molti RDC forniscono non solo posizione ma anche una stima di velocita` gia` ben filtrata, estremamente utile nei servoazionamenti.

La terza slide del gruppo riassume in modo molto netto il compromesso tecnologico: grande robustezza, affidabilità, tolleranza a temperature elevate e buona immunità a sporco e olio da un lato; necessità di elaborazione analogica e risoluzione nativa inferiore a quella dei migliori encoder ottici dall’altro. È per questo che il resolver resta fortissimo in ambienti industriali severi, applicazioni aerospaziali e attuazione di potenza, dove la robustezza vale più della massima risoluzione ottica.

\begin{keybox}
\textbf{Mappa rapida delle tecnologie.} L’encoder incrementale eccelle quando servono conteggio veloce e integrazione digitale. L’encoder assoluto privilegia disponibilità immediata della posizione. LVDT e RVDT offrono grande linearità e robustezza locale, ma richiedono eccitazione AC e demodulazione. Il resolver estende la stessa filosofia a una misura angolare continua su tutto il giro: meno immediato da interfacciare di un encoder, ma molto più tollerante agli ambienti difficili.
\end{keybox}

\subsection*{Sintesi finale della Lezione 5}

Questa quinta lezione amplia in modo deciso il panorama dei sensori di posizione. Lo stesso stato meccanico del robot può essere rappresentato come conteggio di impulsi, parola binaria assoluta, fase di segnali sinusoidali oppure differenza di accoppiamenti magnetici. Ogni scelta porta con sé un compromesso diverso tra risoluzione, robustezza, elettronica necessaria, costo e facilità di integrazione nel controllo.

Il punto più importante da ricordare è che il sensore non è mai solo un componente da catalogo: è parte del modello complessivo del sistema. Un encoder incrementale richiede logiche di conteggio affidabili; un encoder assoluto richiede una codifica robusta; un LVDT, un RVDT o un resolver richiedono demodulazione e condizionamento coerenti con la dinamica del loop. Capire questi dettagli significa progettare una retroazione più credibile, più stabile e più adatta all’applicazione robotica reale.

% =========================================================
% LEZIONE 6
% =========================================================

\clearpage
\section{Lezione 6}

\subsection{Visione generale della lezione}

Questa lezione riprende il tema dei sensori elettromagnetici di posizione e lo porta su un piano piu` costruttivo e applicativo. Dopo aver introdotto il resolver nella lezione precedente, qui l’attenzione si concentra su un problema molto concreto: come alimentare il rotore mantenendo una misura affidabile nel tempo. Da questo punto di vista nasce il confronto tra resolver con slip rings e resolver brushless.

La seconda parte della lezione allarga poi il quadro ai \emph{synchro}, dispositivi elettromeccanici molto vicini ai resolver ma pensati anche per trasmettere l’informazione angolare a distanza. Il blocco finale introduce invece gli encoder magnetici, in particolare la variante rotativa assoluta basata su due componenti ortogonali del campo magnetico. Nel complesso, la lezione mostra tre strategie diverse per misurare o trasmettere un angolo in ambienti reali: contatto elettrico strisciante, accoppiamento magnetico contactless e codifica magnetica diretta.

\subsection{Resolver con e senza slip rings}

\slidepair{images5/page-01.png}{images5/page-02.png}{Resolver: problema pratico dell’eccitazione rotorica e principali modi per trasferirla al rotore.}

Le prime due slide chiariscono bene un punto che spesso, quando si studia il resolver solo a livello di formula, resta implicito. Per funzionare, il resolver ha bisogno che il rotore sia eccitato con un segnale AC di riferimento, tipicamente scritto come
\[
V_r(t)=V_0\sin(\omega t).
\]
Poiche' il rotore gira, questa eccitazione deve essere trasferita a una parte in movimento. Tutta la differenza tra resolver con slip rings e resolver brushless nasce da qui.

Il caso con slip rings usa un collegamento elettrico diretto tra parte fissa e parte rotante mediante anelli e spazzole. Il caso brushless, invece, evita il contatto elettrico e trasferisce la portante con un accoppiamento magnetico, cioe` tramite un piccolo trasformatore rotante. Il vantaggio della seconda soluzione non e` nella formula finale seno-coseno, che resta concettualmente la stessa, ma nell’affidabilita` meccanica e nella qualita` del segnale lungo la vita del dispositivo.

\slidefig{images5/page-03.png}{Resolver con slip rings: collegamento diretto, punti di forza e principali limiti pratici.}

La slide dedicata ai resolver con slip rings evidenzia un compromesso classico. Dal lato positivo, la soluzione e` semplice, economica e capace di fornire un’eccitazione rotorica di buona ampiezza, perche' non introduce un ulteriore stadio magnetico tra sorgente e rotore. Dal lato negativo, pero`, compaiono tutti i limiti dei contatti striscianti: usura delle spazzole, manutenzione, variazioni della resistenza di contatto, rumore elettrico e vita utile limitata.

Dal punto di vista della misura, questo significa che il segnale di eccitazione non e` perfettamente costante. Piccole fluttuazioni di contatto possono tradursi in variazioni della portante, disturbi impulsivi o addirittura brevi dropout. Poiche' il resolver codifica l’angolo come rapporto coerente tra segnali modulati, ogni degradazione della portante si riflette poi sulla qualita` della stima di posizione.

\subsection{Resolver brushless e accoppiamento magnetico}

\slidepair{images5/page-04.png}{images5/page-06.png}{Resolver brushless: struttura con rotary transformer e realizzazione fisica del trasferimento di eccitazione.}

Le slide successive mostrano la soluzione brushless. Qui non esiste alcun contatto elettrico diretto con il rotore: l’eccitazione viene trasferita da un avvolgimento primario stazionario a un avvolgimento secondario montato sul rotore tramite accoppiamento magnetico. La struttura puo` essere letta come una cascata di due dispositivi:
\[
\text{trasformatore di eccitazione} \longrightarrow \text{resolver vero e proprio}.
\]
Il primo stadio porta energia alla parte rotante; il secondo usa tale eccitazione per generare le due uscite seno e coseno in funzione dell’angolo.

Se si indica con \(V_{\text{rotor}}\) la tensione effettivamente presente sul rotore dopo il trasferimento contactless, il comportamento ideale del resolver resta quello gia` visto nella lezione precedente:
\[
V_c \propto V_{\text{rotor}}\cos\theta,
\qquad
V_s \propto V_{\text{rotor}}\sin\theta.
\]
La differenza, quindi, non sta nella codifica geometrica dell’angolo, ma nel modo in cui si ottiene e si mantiene la portante rotorica.

\slidefig{images5/page-07.png}{Accoppiamento magnetico nel resolver brushless: modello a due stadi, sfasamento di fase e attenuazione della portante.}

La slide sull’accoppiamento magnetico e` la piu` importante del blocco, perche' rende esplicito il modello equivalente. La tensione sul rotore puo` essere descritta in prima approssimazione come
\[
V_{\text{rotor}} = k\,V_{\text{primary}},
\]
oppure, in forma temporale piu` realistica,
\[
V_{\text{rotor}}(t)=V_0\sin(\omega t+\phi_e),
\]
dove \(k\) rappresenta il trasferimento di ampiezza del trasformatore rotante e \(\phi_e\) il piccolo sfasamento introdotto dal trasferimento magnetico.

Questo significa che il resolver brushless si comporta come un sistema a due trasformatori in cascata: uno trasferisce l’eccitazione alla parte rotante, l’altro realizza la modulazione seno-coseno legata all’angolo. Il prezzo da pagare e` un leggero calo di ampiezza e un piccolo sfasamento di fase della portante; il beneficio, pero`, e` molto grande in termini di affidabilita`, perche' vengono eliminati del tutto i contatti soggetti a usura.

Se si scrivono le uscite in forma esplicita, si ottiene quindi una struttura del tipo
\[
V_s(t)=K_sV_0\sin\theta\,\sin(\omega t+\phi_e),
\qquad
V_c(t)=K_cV_0\cos\theta\,\sin(\omega t+\phi_e).
\]
Rispetto al resolver ideale, l’informazione angolare continua a vivere nei fattori \(\sin\theta\) e \(\cos\theta\), ma la catena di demodulazione deve tenere conto della fase reale della portante. E` proprio per questo che i resolver brushless richiedono spesso una compensazione di fase nell’elettronica di lettura.

\slidepair{images5/page-08.png}{images5/page-09.png}{Qualita` del segnale nei resolver e confronto prestazionale tra architettura con slip rings e brushless.}

Le due slide finali del blocco riassumono bene il confronto applicativo. Nei resolver con slip rings, il problema principale e` la qualita` instabile del contatto elettrico: variazioni di resistenza, rumore e possibili interruzioni degradano integrita` del segnale e durata del sensore. Nei resolver brushless, invece, l’accoppiamento e` molto piu` stabile e la qualita` del segnale migliora, pur a costo di una leggera attenuazione e di un piccolo sfasamento da compensare.

La tabella di confronto rende esplicito il compromesso: gli slip rings costano meno e hanno struttura piu` semplice, ma soffrono di usura e manutenzione; il brushless costa di piu`, ma offre rumore piu` basso, affidabilita` molto maggiore e vita utile piu` lunga. In applicazioni industriali severe, il resolver brushless e` quindi spesso la scelta naturale.

\begin{keybox}
\textbf{Idea pratica.} I resolver con slip rings e quelli brushless codificano l’angolo allo stesso modo. La vera differenza sta in come si trasferisce la portante al rotore: contatto diretto piu` semplice ma usurabile, oppure accoppiamento magnetico piu` robusto ma leggermente piu` complesso da compensare.
\end{keybox}

\subsection{Synchro}

\slidepair{images5/page-10.png}{images5/page-11.png}{Synchro: struttura a trasformatore rotante con statore trifase e relazione tra tensioni indotte e angolo.}

Le slide sul synchro introducono un dispositivo molto vicino al resolver dal punto di vista elettromagnetico, ma con un’architettura diversa: il rotore e` eccitato in AC, mentre lo statore ha tre avvolgimenti sfasati di \(120^\circ\). Indicando con
\[
V_r(t)=V_0\sin(\omega t)
\]
l’eccitazione rotorica, le tre tensioni statoriche possono essere scritte nella forma
\[
V_{s1}=K\sin(\omega t)\cos(\theta-0^\circ),
\]
\[
V_{s2}=K\sin(\omega t)\cos(\theta-120^\circ),
\]
\[
V_{s3}=K\sin(\omega t)\cos(\theta+120^\circ).
\]
La posizione angolare e` quindi distribuita su una terna di segnali anziche' sulla classica coppia seno-coseno del resolver.

Dal punto di vista concettuale, il synchro puo` essere letto come una codifica elettromagnetica trifase dell’angolo. Questa rappresentazione e` particolarmente utile quando non si vuole soltanto misurare l’angolo localmente, ma anche trasmetterlo elettricamente a distanza o usarlo come riferimento per un altro asse.

\slidefig{images5/page-12.png}{Sistema synchro TX-RX-CT: trasmissione master-slave e generazione di un errore angolare elettrico.}

La slide successiva mostra proprio questa vocazione alla trasmissione. In un sistema classico compaiono tre blocchi: \emph{Synchro Transmitter} (TX), \emph{Synchro Receiver} (RX) e \emph{Control Transformer} (CT). Se lo statore del trasmettitore viene collegato elettricamente a quello del ricevitore, il rotore del ricevitore tende ad allinearsi all’angolo del trasmettitore:
\[
\theta_{RX}\rightarrow \theta_{TX}.
\]
In questo modo il synchro non si limita a ``misurare'' un angolo: puo` anche replicarlo o trasmetterlo meccanicamente a distanza.

Nel control transformer, invece, l’uscita viene costruita come errore angolare elettrico. La slide la scrive come
\[
V_{out}=K\sin(\theta_{TX}-\theta_{CT}),
\]
che si annulla quando master e slave sono allineati. Questa e` un’idea molto importante per il controllo: il sensore non restituisce necessariamente la posizione assoluta gia` pronta, ma puo` fornire direttamente una grandezza di errore utile a pilotare un servoanello.

\slidepair{images5/page-13.png}{images5/page-14.png}{Principali applicazioni dei synchro ed esempio costruttivo con statore trifase e rotore con slip rings.}

Le slide applicative collocano il synchro nel suo contesto storico e industriale: avionica, sistemi di navigazione navale, indicatori remoti, posizionamento di antenne e, in generale, contesti in cui l’informazione angolare deve essere trasmessa in modo robusto in ambienti soggetti a vibrazioni o disturbi elettromagnetici. La presenza di tre avvolgimenti statorici e, spesso, di slip rings sul rotore rende il dispositivo meno compatto e meno ``digitale'' di un encoder moderno, ma molto interessante quando contano robustezza e trasmissione remota analogica dell’angolo.

Per questo, nella robotica contemporanea, il synchro compare meno spesso nelle architetture compatte e low-cost, ma resta concettualmente importante: mostra che una misura di posizione puo` essere trattata come grandezza elettromagnetica distribuita, trasmessa e confrontata direttamente all’interno di un sistema servoassistito.

\subsection{Encoder magnetici}

\slidefig{images5/page-15.png}{Encoder magnetici: classificazione generale, componenti fondamentali e vantaggi nei contesti gravosi.}

L’ultima parte della lezione cambia tecnologia ma conserva la stessa logica di fondo: usare un fenomeno fisico robusto per ottenere informazione di posizione. Gli encoder magnetici possono essere lineari o rotativi, assoluti o relativi. Gli elementi base indicati nelle slide sono tre: un magnete permanente montato sul rotore o su una strip lineare, un sensore magnetico e l’elettronica di condizionamento.

La famiglia dei sensori magnetici citata nella slide e` molto ampia: Hall, AMR (\emph{anisotropic magnetoresistance}), GMR (\emph{giant magnetoresistance}) e TMR (\emph{tunneling magnetoresistance}). Il principio applicativo, pero`, resta comune: misurare localmente il campo magnetico generato dal magnete e convertirlo in una stima di posizione. Il vantaggio pratico e` evidente: buona robustezza a polvere, olio e vibrazioni, dimensioni contenute e costo relativamente basso.

\slidepair{images5/page-16.png}{images5/page-17.png}{Encoder magnetico rotativo assoluto: ricostruzione dell’angolo da due componenti ortogonali del campo magnetico.}

Le ultime due slide introducono la variante rotativa assoluta piu` semplice e diffusa. Un magnete cilindrico magnetizzato trasversalmente genera un campo che, visto dal sensore, ruota con l’albero. Due elementi sensibili ortogonali misurano le due componenti del flusso magnetico, tipicamente indicate con \(B_x\) e \(B_y\). In prima approssimazione si puo` scrivere
\[
B_x = B_0\cos\theta,
\qquad
B_y = B_0\sin\theta.
\]
L’angolo assoluto si ricava allora in modo diretto come
\[
\theta=\operatorname{atan2}(B_y,B_x).
\]

Questo schema e` concettualmente molto elegante, perche' realizza una versione magnetica della coppia seno-coseno gia` vista negli encoder ottici interpolati e nei resolver demodulati. Anche qui la funzione \(\operatorname{atan2}\) e` la scelta corretta, perche' consente di ricostruire il quadrante senza ambiguita`.

Le caratteristiche riportate nella slide finale spiegano bene perche' questi encoder sono cosi` diffusi: posizione vera disponibile all’accensione, nessuna procedura di homing e risoluzioni tipiche dell’ordine di \(14\)–\(18\) bit. In molte applicazioni meccatroniche questa combinazione li rende un ottimo compromesso tra robustezza, costo e facilita` di integrazione.

Naturalmente non si tratta di sensori ``perfetti''. Eccentricita` del magnete, non ortogonalita` dei canali, offset dei sensori e distorsioni del campo magnetico possono trasformare la traiettoria ideale nel piano \((B_x,B_y)\) da cerchio a ellisse decentrata. Per questo, nei dispositivi di qualita` piu` alta, si eseguono compensazioni molto simili a quelle viste per resolver ed encoder interpolati.

\begin{keybox}
\textbf{Mappa rapida della lezione.} Il resolver con slip rings privilegia semplicit\`a e costo ma soffre di usura. Il resolver brushless sacrifica qualcosa in ampiezza e richiede compensazione di fase, ma guadagna tantissimo in affidabilita`. Il synchro aggiunge la possibilita` di trasmettere l’angolo elettricamente in configurazione master-slave. L’encoder magnetico assoluto, infine, fornisce direttamente la posizione con una struttura compatta e molto robusta agli ambienti sporchi.
\end{keybox}

\subsection*{Sintesi finale della Lezione 6}

Questa sesta lezione mostra che, nei sensori di posizione elettromagnetici, la geometria della codifica e la qualita` dell’interfacciamento sono inseparabili. Il resolver brushless migliora il resolver tradizionale non cambiando la legge seno-coseno, ma cambiando il modo in cui la portante arriva al rotore. Il synchro mostra come la stessa fisica possa essere usata non solo per misurare un angolo, ma anche per trasmetterlo e confrontarlo a distanza. L’encoder magnetico assoluto, infine, porta la stessa idea di codifica ortogonale in una forma piu` compatta, economica e nativamente adatta all’elettronica moderna.

Il filo conduttore da ricordare e` quindi questo: una misura affidabile di posizione non dipende solo dalla formula finale con cui si ricava \(\theta\), ma dall’intera catena fisica che genera, trasferisce, condiziona e interpreta il segnale. E` precisamente questa catena che determina robustezza, rumore, manutenzione necessaria e qualita` della retroazione nel sistema robotico reale.

% =========================================================
% LEZIONE 7
% =========================================================

\clearpage
\section{Lezione 7}

\subsection{Visione generale della lezione}

Questa lezione chiude il blocco dedicato ai sensori di posizione e apre in modo esplicito il tema della misura di distanza e di velocita`. Il punto di partenza e` ancora una volta la tecnologia magnetica, ma in una forma piu` semplice e robusta: encoder rotativi relativi e sensori lineari magnetici. Da qui la trattazione passa ai sensori di distanza, cioe` dispositivi che non misurano direttamente una coordinata interna del robot, ma la separazione da un bersaglio o da un ostacolo.

La seconda meta` della lezione entra poi nel tema della velocita`. Prima vengono introdotti i sensori che forniscono una misura diretta della velocita`, come i tachogeneratori, e poi quelli che la ricavano in modo indiretto da misure di posizione o da impulsi incrementali. L’ultima parte mostra infine un fatto molto importante per il controllo: in molti casi la velocita` migliore non e` quella misurata ``grezzamente'', ma quella ottenuta con un osservatore o con un filtro di stima costruito sul modello dinamico del sistema.

\subsection{Encoder magnetici relativi e sensori lineari magnetici}

\slidepair{images6/page-01.png}{images6/page-02.png}{Encoder magnetici relativi: realizzazione con magneti permanenti o con struttura dentata ferromagnetica.}

Le prime due slide mostrano due modi diversi di costruire un encoder magnetico relativo. Nel primo caso si usa un magnete cilindrico magnetizzato trasversalmente, oppure un disco con poli nord-sud alternati; nel secondo si usa invece una ruota dentata ferromagnetica che modifica il campo generato da un magnete permanente. In entrambi i casi il sensore non legge una posizione assoluta gia` codificata, ma rileva il passaggio periodico di poli o denti.

Dal punto di vista funzionale, il comportamento e` molto vicino a quello di un encoder incrementale ottico: il moto genera una sequenza di eventi, la posizione viene ricostruita integrando gli impulsi e la direzione richiede un secondo canale o una seconda sonda opportunamente sfalsata. Se il passo magnetico corrisponde a \(N_p\) eventi per giro, la stima della posizione puo` essere scritta nella forma
\[
\theta[k]=\frac{2\pi}{N_p}\,n[k],
\]
dove \(n[k]\) e` il conteggio accumulato. Il prezzo da pagare, come ricordano le slide, e` che occorre homing e la risoluzione tipica e` spesso inferiore a quella dei migliori encoder ottici.

La variante con ruota dentata e sensore di flusso magnetico e` particolarmente interessante dal punto di vista industriale, perche' e` robusta, semplice e ben adattata ad ambienti con sporco, olio o vibrazioni. Qui il passaggio di un dente devia il campo magnetico e il sensore ne rileva la variazione locale. E` una logica molto usata nei sensori di velocita` ruota in ambito automotive e, piu` in generale, in tutte le applicazioni dove la robustezza meccanica conta piu` della risoluzione estrema.

\slidetriple{images6/page-03.png}{images6/page-04.png}{images6/page-05.png}{Sensori lineari magnetici: principio, limiti pratici e interpolazione sub-pitch con segnali in quadratura.}

Le tre slide successive traslano la stessa idea dal moto rotativo a quello lineare. Al posto del disco magnetico compare una strip magnetizzata o una scala magnetica; al posto del sensore puntuale rotativo compare una testina di lettura con elementi Hall o magnetoresistivi. Il principio di base e` sempre lo stesso: il campo generato dalla scala viene campionato localmente e convertito in una stima di posizione.

La slide con la formula dell’effetto Hall riporta
\[
V_H=\frac{I\,B}{q\,n\,t},
\]
cioe` una tensione proporzionale alla densita` di flusso magnetico \(B\) quando una corrente \(I\) attraversa un materiale di spessore \(t\) con densita` di portatori \(n\). Questo spiega perche' i sensori Hall o MR siano cosi` utili in queste architetture: consentono di misurare il campo magnetico senza contatto e con elettroniche relativamente compatte.

Nel caso incrementale, la geometria della testina permette di ottenere due segnali quasi sinusoidali in quadratura:
\[
V_s=A\sin\!\left(\frac{2\pi x}{p}\right),
\qquad
V_c=A\cos\!\left(\frac{2\pi x}{p}\right),
\]
dove \(p\) e` il passo magnetico. Da qui la posizione fine si ricava con
\[
x=\frac{p}{2\pi}\operatorname{atan2}(V_s,V_c)+kp,
\]
esattamente come negli encoder ottici interpolati. La velocita` puo` poi essere stimata per differenziazione numerica, ad esempio con \(\hat v \approx \Delta \hat x/\Delta t\).

Le slide sottolineano bene il compromesso applicativo: questi sensori lavorano bene in presenza di polvere o olio e hanno meccanica semplice, ma soffrono di distorsioni del campo, deriva termica e invecchiamento del magnete. In pratica, sono ottimi quando si cerca robustezza ambientale con una buona risoluzione, ma richiedono compensazione se si vuole precisione elevata su lunghe corse.

\subsection{Sensori di distanza}

\slidepair{images6/page-06.png}{images6/page-07.png}{Sensori di distanza: quadro generale e principio del range sensing ultrasonico.}

La parte successiva della lezione cambia prospettiva: invece di misurare posizione di un albero o di una slitta, si misura la distanza da un bersaglio o da un ostacolo. Le slide ricordano che questa informazione puo` essere usata direttamente come dato di prossimita` oppure indirettamente per stimare la posizione del robot rispetto all’ambiente, per esempio combinando una mappa con piu` misure di distanza.

Le due famiglie principali introdotte sono ultrasonica e ottica. Il sensore ultrasonico a tempo di volo (\emph{Time-of-Flight}, ToF) emette un impulso acustico e misura il tempo che passa prima di ricevere l’eco riflessa dal bersaglio. E` una misura assoluta di distanza, robusta e poco costosa, adatta a obstacle detection, level sensing e semplici misure di range.

\slidepair{images6/page-08.png}{images6/page-09.png}{Tempo di volo ultrasonico: schema di eco, dipendenza dalla velocita` del suono e ricostruzione della distanza.}

Le slide quantitative sul ToF ultrasonico esplicitano la relazione fondamentale:
\[
d=\frac{c\,t_{\text{echo}}}{2},
\]
dove \(t_{\text{echo}}\) e` il tempo di andata e ritorno dell’onda e \(c\) la velocita` del suono nell’aria. Le slide riportano anche l’approssimazione pratica
\[
c \approx 331 + 0.6\,T(^\circ\text{C})\ \text{m/s},
\]
che rende immediato un punto importante: la misura dipende dall’ambiente, in particolare da temperatura e, in misura minore, da umidita`.

Se si vuole una stima grezza della velocita` lungo la direzione di misura, si puo` usare
\[
\hat v \approx \frac{\Delta d}{\Delta t},
\]
ma la slide sottolinea correttamente che si tratta di una stima a bassa banda. L’ultrasonico e` molto utile quando conta sapere ``quanto lontano'' si trova un ostacolo, meno quando servono aggiornamenti rapidissimi o precisione fine sulla dinamica.

\slidepair{images6/page-10.png}{images6/page-11.png}{Triangolazione laser: principio fisico, impieghi tipici e geometria elementare della misura.}

Nel sensore ottico a triangolazione, invece, un punto laser viene proiettato sul bersaglio e la sua immagine si sposta su un sensore PSD o CMOS al variare della distanza. La geometria del sistema trasforma questo spostamento d’immagine in una stima della distanza. Per piccole angolazioni e in configurazioni semplici, la slide riporta l’approssimazione
\[
z \approx \frac{b\,f}{u},
\]
dove \(b\) e` la baseline ottica, \(f\) la lunghezza focale e \(u\) lo spostamento dell’immagine sul sensore.

Il punto forte di questa tecnologia e` la rapidita`: il sensore e` non-contact, veloce e adatto a brevi o medie distanze. Le applicazioni citate nelle slide vanno dall’evitamento ostacoli in robotica mobile alla misura di grandi spostamenti in macchine industriali. I limiti principali sono invece la sensibilita` a riflettivita` superficiale e luce ambiente, cioe` proprio a fattori che alterano la qualita` del punto ottico osservato.

\begin{keybox}
\textbf{Confronto rapido.} L’ultrasonico misura il tempo di volo di un’onda acustica: economico e robusto, ma lento e sensibile alle condizioni dell’aria. La triangolazione laser usa invece una geometria ottica: e` piu` rapida e precisa a corto-medio raggio, ma piu` sensibile alle proprieta` della superficie e alla luce ambiente.
\end{keybox}

\subsection{Sensori di velocita`: classificazione e principio elettromagnetico di base}

\slidepair{images6/page-12.png}{images6/page-13.png}{Passaggio ai sensori di velocita` e distinzione tra misura diretta e indiretta.}

Le slide di apertura del blocco velocita` separano subito due idee molto diverse. La velocita` puo` essere misurata \emph{direttamente}, quando il sensore ha un’uscita fisicamente proporzionale a \(\omega\) o a \(v\), oppure \emph{indirettamente}, quando viene ricavata da una misura di posizione tramite differenze finite o conteggio impulsi:
\[
\hat v \approx \frac{\Delta x}{\Delta t}.
\]
La distinzione e` importante perche' cambia sia il rumore sia il ritardo della misura disponibile al controllore.

Le applicazioni citate dalle slide coprono due scenari diversi: la misura della velocita` relativa di parti di una macchina, come albero motore o carrello lineare, e la misura della propria velocita` in robotica mobile, come velocita` angolare di un drone o velocita` lineare di un velivolo. Sono problemi simili dal punto di vista matematico, ma molto diversi per dinamica, rumore e tecnologia impiegata.

\slidefig{images6/page-14.png}{Tassonomia sintetica dei sensori di velocita`: misure dirette e stime indirette da posizione.}

La tassonomia riportata nelle slide aiuta a mettere ordine: tra i sensori diretti compaiono tachogeneratori DC e AC, sensori Doppler, sensori di flusso e giroscopi MEMS; tra quelli indiretti compaiono encoder, resolver, LVDT/RVDT, encoder lineari e perfino sistemi di visione con tracking. La lezione si concentra soprattutto sui tachogeneratori e sulla velocita` stimata da encoder, ma il messaggio piu` generale e` che ``sensore di velocita`'' non identifica una singola tecnologia: identifica un problema di misura che puo` essere risolto in modi molto diversi.

\slidetriple{images6/page-15.png}{images6/page-16.png}{images6/page-17.png}{Legge di Faraday applicata a una spira rotante in campo magnetico uniforme.}

Il blocco successivo ricostruisce il principio fisico di base che rende possibili i tachogeneratori. Si considera una spira conduttrice di area \(A\) che ruota a velocita` angolare \(\omega\) in un campo magnetico uniforme \(B\). Il flusso magnetico concatenato vale
\[
\Phi(t)=BA\cos\theta(t),
\qquad
\theta(t)=\omega t + \theta_0.
\]
Applicando la legge di Faraday,
\[
e(t)=-\frac{d\Phi}{dt},
\]
si ottiene
\[
e(t)=BA\omega\sin(\omega t),
\]
e, per una bobina con \(N\) spire,
\[
e(t)=NBA\omega\sin(\omega t).
\]

Il punto chiave, che la slide finale del blocco rende molto chiaro, e` che l’ampiezza della tensione indotta e` proporzionale alla velocita` angolare \(\omega\). E` esattamente da qui che nasce l’idea del tachogeneratore: convertire una velocita` meccanica in una tensione la cui ampiezza cresca linearmente con essa.

\subsection{Tachogeneratore DC}

\slidepair{images6/page-18.png}{images6/page-19.png}{Tachogeneratore DC: principio di misura e ruolo di collettore e spazzole.}

Le prime due slide sul tachogeneratore DC riassumono il principio in modo estremamente compatto:
\[
V_{\text{out}} \simeq K\omega.
\]
Il rotore gira nel campo magnetico, la tensione indotta cresce con la velocita` e il dispositivo fornisce quindi una misura analogica diretta della velocita` senza richiedere differenziazione numerica. Questo e` il grande vantaggio rispetto alla stima ottenuta da un encoder di posizione.

La seconda slide, pero`, chiarisce subito un dettaglio essenziale: senza collettore la singola bobina produrrebbe una tensione alternata del tipo
\[
e(t)=E_0\sin(\omega t).
\]
Per ottenere un’uscita unidirezionale proporzionale alla velocita` occorrono molte bobine, un collettore segmentato e spazzole stazionarie. In pratica, il tachogeneratore DC usa una rettifica meccanica direttamente all’interno della macchina elettrica.

\slidetriple{images6/page-20.png}{images6/page-21.png}{images6/page-22.png}{Tachogeneratore DC: commutazione tra piu` bobine, ripple residuo e rettifica meccanica.}

Le slide centrali del blocco spiegano come la tensione quasi continua venga costruita sommando i contributi di molte bobine sfasate. Se la \(i\)-esima bobina ha fase
\[
\theta_i=\theta+\frac{2\pi i}{N},
\]
la tensione associata puo` essere schematizzata come
\[
e_i(t)=E_0\sin\theta_i.
\]
Il collettore seleziona di volta in volta le bobine la cui tensione e` vicina al massimo positivo, in modo da produrre in uscita una grandezza quasi continua.

La slide piu` importante del gruppo scrive infatti
\[
V_{\text{out}}(t)\approx K\omega + \varepsilon(t),
\]
dove \(\varepsilon(t)\) e` il ripple. Aumentando il numero di bobine e di segmenti di collettore, il ripple si riduce. Le spazzole e il collettore agiscono quindi come un raddrizzatore meccanico che trasforma la tensione AC indotta nell’armatura in una tensione media proporzionale alla velocita`.

\slidefig{images6/page-23.png}{Prestazioni e limiti del tachogeneratore DC.}

L’ultima slide del blocco ne riassume molto bene il profilo applicativo. I punti forti sono buona linearita`, uscita analogica diretta e polarita` che indica anche il verso di rotazione. I punti deboli sono invece usura delle spazzole, ripple residuo, rumore ad alta frequenza dovuto alla commutazione e prestazioni non brillanti a bassa velocita`, dove la tensione utile diventa piccola e i disturbi relativi pesano di piu`.

Le slide citano anche un ordine di grandezza tipico del ripple, spesso nell’intervallo \(1.5\%\)–\(5\%\) dell’uscita. Per questo, nei sistemi di controllo di precisione, un tachogeneratore DC richiede quasi sempre uno stadio di filtraggio o smoothing prima di usare il segnale all’interno del loop.

\subsection{Tachogeneratore AC}

\slidepair{images6/page-24.png}{images6/page-25.png}{Tachogeneratore AC: uscita sinusoidale proporzionale alla velocita` e modalita` di lettura.}

Il tachogeneratore AC nasce dalla stessa fisica elettromagnetica, ma evita collettore e spazzole. Le slide lo descrivono come una macchina con armatura statica e campo magnetico rotante, capace di generare
\[
V_{\text{out}}(t)=K\omega\sin(\omega t).
\]
Anche qui l’ampiezza e` proporzionale alla velocita`, ma l’uscita resta alternata.

La velocita` puo` allora essere ricavata in due modi: misurando l’ampiezza del segnale, tipicamente dopo rettifica e filtraggio, oppure misurandone la frequenza con un contatore. Il vantaggio principale rispetto al tachogeneratore DC e` l’assenza di spazzole e collettore, quindi minore usura e minore rumore di commutazione. Lo svantaggio e` che l’uscita non e` direttamente una tensione continua gia` pronta per il controllore, ma richiede comunque condizionamento.

\subsection{Stima di velocita` da encoder: metodo M}

\slidefig{images6/page-26.png}{Idea generale della stima di velocita` da sensori incrementali tramite frequenza degli impulsi.}

La lezione passa poi ai metodi indiretti, cioe` ai casi in cui la velocita` viene stimata a partire da un encoder o da un sensore di posizione impulsivo. L’idea di base e` semplice: se gli impulsi arrivano piu` frequentemente, la velocita` e` maggiore. Il problema, pero`, e` trasformare questa intuizione in una stima numerica con buon compromesso tra precisione, banda e ritardo.

\slidepair{images6/page-27.png}{images6/page-28.png}{Metodo M: conteggio impulsi in una finestra temporale fissa e calcolo diretto della velocita`.}

Nel metodo M (\emph{Fixed Time Method}) si conta il numero \(m\) di impulsi osservati in una finestra temporale fissa \(T_c\). Lo spostamento angolare accumulato nella finestra vale
\[
X=\frac{m}{PPR}\,2\pi,
\]
dove \(PPR\) e` il numero di impulsi per giro del sensore. La velocita` angolare stimata diventa allora
\[
\omega_m=\frac{X}{T_c}=\frac{2\pi}{T_c}\frac{m}{PPR},
\]
mentre in giri al minuto si ha
\[
N=\frac{60}{2\pi}\omega_m=\frac{60}{T_c}\frac{m}{PPR}.
\]

Il pregio fondamentale di questo metodo e` che il tempo di campionamento della velocita` e` costante. Dal punto di vista del controllo, questo e` comodissimo, perche' il regolatore riceve nuovi campioni a intervalli regolari e il progetto discreto del loop diventa piu` semplice.

\slidepair{images6/page-29.png}{images6/page-30.png}{Metodo M: errore di quantizzazione degli impulsi e criticita` a bassa velocita`.}

Le slide successive spiegano bene il limite strutturale del metodo M: l’errore massimo sul conteggio e` di un impulso, quindi l’errore massimo sulla velocita` in rpm e` dell’ordine di
\[
\Delta N_{\max}=\frac{60}{T_c\,PPR}.
\]
Se \(T_c\) e` piccolo, la banda cresce ma l’errore di quantizzazione puo` diventare molto grande. Se invece \(T_c\) viene aumentato, la misura si stabilizza ma il ritardo cresce e la banda si riduce.

Il problema peggiore si presenta a bassa velocita`, quando nella finestra \(T_c\) arrivano pochi impulsi. Le slide mostrano proprio questo caso: velocita` reale costante ma misure discrete diverse da una finestra all’altra. Si puo` migliorare la situazione aumentando il \(PPR\) oppure allungando \(T_c\), ma entrambe le soluzioni hanno un costo: sensore piu` caro nel primo caso, minore reattivita` del controllo nel secondo.

\subsection{Metodo T e metodo M/T}

\slidepair{images6/page-31.png}{images6/page-32.png}{Metodo T: misura del tempo tra impulsi consecutivi e ricostruzione della velocita` dal periodo.}

Il metodo T (\emph{Pulse distance measurement}) fa il contrario del metodo M. Invece di contare quanti impulsi cadono in un tempo fissato, misura quanto tempo passa tra due impulsi consecutivi. Lo spostamento angolare associato a due impulsi adiacenti e` fisso:
\[
X=\frac{2\pi}{PPR}.
\]
Se il periodo tra impulsi vale \(T_c\), la velocita` e`
\[
\omega_m=\frac{X}{T_c}.
\]
Nella realizzazione digitale mostrata dalla slide, \(T_c\) si ottiene contando \(m_c\) impulsi di un clock di riferimento di periodo \(T_{\text{clock}}\):
\[
T_c=m_cT_{\text{clock}}.
\]
Quindi
\[
\omega_m=\frac{2\pi}{PPR\,m_cT_{\text{clock}}},
\qquad
N=\frac{60}{PPR\,m_cT_{\text{clock}}}.
\]

Questo metodo elimina l’errore di ``impulso perso'' del metodo M ed e` molto piu` accurato alle basse velocita`, dove il tempo tra impulsi e` lungo e quindi facilmente misurabile. Il rovescio della medaglia, pero`, e` che i nuovi campioni di velocita` arrivano solo quando arriva un nuovo impulso. Il campionamento della velocita` non e` quindi uniforme e, per un controllore di velocita`, questo e` un problema serio.

\slidepair{images6/page-33.png}{images6/page-34.png}{Limiti del metodo T alle alte velocita` e introduzione del metodo ibrido M/T.}

Le slide successive mostrano bene il bilanciamento opposto del metodo T: alle alte velocita` il periodo tra impulsi diventa molto piccolo e il clock di riferimento deve essere molto veloce per mantenere una stima accurata. Per superare i limiti complementari dei metodi M e T, la lezione introduce il metodo ibrido M/T.

Nel metodo M/T si contano gli impulsi \(m_1\) in una finestra temporale nominale \(T_c\), come nel metodo M, ma si misura anche il tempo residuo \(\Delta T\) tra la fine della finestra e l’ultimo impulso effettivamente utile, come nel metodo T. In questo modo si mantiene un aggiornamento quasi regolare senza perdere del tutto l’informazione temporale fine.

\slidefig{images6/page-35.png}{Metodo M/T: formula completa con conteggio impulsi e correzione temporale residua.}

La slide formula il metodo in modo molto pulito. Lo spostamento associato ai \(m_1\) impulsi osservati vale
\[
X=\frac{2\pi}{PPR}m_1.
\]
Se \(\Delta T\) viene misurato contando \(m_2\) impulsi di un clock di periodo \(T_2\), il tempo totale associato ai \(m_1\) impulsi e`
\[
T=T_c+\Delta T=T_c+m_2T_2.
\]
La velocita` si stima allora come
\[
\omega_m=\frac{X}{T}=\frac{2\pi}{PPR}\frac{m_1}{T_c+m_2T_2},
\]
e in rpm
\[
N=\frac{60}{PPR}\frac{m_1}{T_c+m_2T_2}.
\]

\slidepair{images6/page-36.png}{images6/page-37.png}{Bilancio finale dei metodi M, T, M/T e passaggio alla stima di stato di posizione e velocita`.}

Le ultime slide del blocco chiariscono bene il messaggio finale: nessuno dei metodi M, T o M/T fornisce la velocita` istantanea vera. Tutti restituiscono una velocita` media discreta sul proprio intervallo di osservazione. Il metodo M e` comodo per il controllo perche' ha campionamento fisso, ma soffre alle basse velocita`. Il metodo T e` ottimo alle basse velocita`, ma produce campionamento variabile. Il metodo M/T e` il compromesso migliore dei tre, ma non elimina del tutto ritardi e irregolarita`.

E` qui che compare la quarta famiglia di metodi, cioe` la stima di stato. L’idea non e` piu` limitarsi a convertire impulsi in velocita`, ma costruire un estimatore che usi insieme il modello dinamico e la misura disponibile per ricostruire posizione e velocita` in modo piu` regolare e piu` robusto al rumore.

\begin{keybox}
\textbf{Lettura pratica dei metodi.} Il metodo M e` naturale quando il controllore lavora a campionamento fisso. Il metodo T e` vantaggioso alle basse velocita`. Il metodo M/T migliora il compromesso generale. Quando pero` servono insieme regolarita`, filtraggio del rumore e migliore coerenza dinamica, conviene passare a un osservatore di stato.
\end{keybox}

\subsection{Stima filtrata di posizione e velocita`}

\slidepair{images6/page-38.png}{images6/page-39.png}{Modello a velocita` costante e osservatore di Luenberger per angolo e velocita`.}

La parte finale della lezione costruisce un osservatore molto semplice ma molto istruttivo. In assenza di altre informazioni, si assume un modello a velocita` costante:
\[
\dot\theta=\omega,
\qquad
\dot\omega=0.
\]
Questo e` il modello minimo ragionevole quando il moto e` liscio e non si vogliono introdurre ipotesi troppo complesse sull’accelerazione.

La slide successiva scrive il problema in forma di stato:
\[
x=\begin{bmatrix}\theta\\ \omega\end{bmatrix},
\qquad
A=\begin{bmatrix}0 & 1\\ 0 & 0\end{bmatrix},
\qquad
C=\begin{bmatrix}1 & 0\end{bmatrix}.
\]
L’osservatore di Luenberger risultante e`
\[
\dot{\hat\theta}=\hat\omega + l_1(\theta-\hat\theta),
\qquad
\dot{\hat\omega}=l_2(\theta-\hat\theta).
\]
Definendo l’errore di stima \(e=x-\hat x\), la dinamica dell’errore diventa
\[
\dot e=(A-LC)e,
\]
e la convergenza si ottiene scegliendo \(L\) in modo che \(A-LC\) sia stabile.

\slidepair{images6/page-40.png}{images6/page-41.png}{Scelta dei guadagni dell’osservatore e altre possibili strutture di filtraggio della velocita`.}

La slide successiva propone una scelta molto comune dei guadagni:
\[
L=
\begin{bmatrix}
2/\varepsilon\\
1/\varepsilon^2
\end{bmatrix}.
\]
In questa parametrizzazione, \(\varepsilon\) fissa la banda dell’osservatore e quindi il compromesso tra rapidita` di inseguimento, ritardo di fase e sensibilita` al rumore. E` proprio qui che l’osservatore smette di essere una semplice formula e diventa uno strumento di progetto per il loop di controllo.

L’ultima slide utile del blocco ricorda poi che l’osservatore di Luenberger e` solo una possibilita`. Vengono citati anche il \emph{fading filter} di Zarchan, il filtro di Kalman e predittori non lineari che sfruttano informazioni aggiuntive, per esempio la coppia applicata a un motore, l’inerzia di carico nota e una stima della coppia disturbante. Il messaggio conclusivo e` molto forte: per ottenere una buona velocita` non basta contare impulsi, spesso bisogna introdurre modello, filtraggio e informazione fisica extra.

\subsection*{Sintesi finale della Lezione 7}

Questa settima lezione allarga in modo netto il corso oltre i soli sensori di posizione ``classici''. I sensori magnetici relativi e lineari mostrano come si possa ottenere una misura robusta con tecnologie semplici e ben adatte ad ambienti sporchi. I sensori di distanza mostrano invece che la posizione puo` essere ricostruita anche in modo indiretto, misurando separazioni geometriche dal mondo esterno. Infine, il blocco sui sensori di velocita` fa vedere che la misura della velocita` puo` essere sia un problema fisico diretto, come nei tachogeneratori, sia un problema di stima numerica a partire da impulsi o modelli.

Il punto piu` importante da portare a casa e` che la velocita` e` spesso la grandezza piu` delicata da usare nel controllo. Differenziare il rumore e` pericoloso, contare pochi impulsi introduce quantizzazione, campionamenti non uniformi complicano il progetto del loop e filtri troppo aggressivi introducono ritardo. Per questo la lezione converge in modo naturale verso gli osservatori: non come artifici matematici scollegati dall’hardware, ma come parte integrante della catena di misura realmente usata nel controllo dei robot.

% =========================================================
% LEZIONE 8
% =========================================================

\clearpage
\section{Lezione 8}

\subsection{Visione generale della lezione}

Questa lezione prosegue il blocco sui sensori di velocita`, ma sposta l'attenzione da una velocita` misurata rispetto a un organo meccanico, come nel tachogeneratore, a una velocita` \emph{propria} del veicolo o del corpo sensorizzato. Il caso piu` importante e` la velocita` angolare: sapere quanto rapidamente un drone ruota attorno ai propri assi, quanto sta imbardando un aereo, oppure come varia l'assetto di un satellite e` una condizione essenziale per stabilizzazione, navigazione e controllo.

Il protagonista della lezione e` quindi il giroscopio. Prima viene introdotto in modo fisico, attraverso inerzia, precessione e nutazione. Poi viene collegato alle applicazioni classiche di navigazione, come indicatore di virata, orizzonte artificiale, directional gyro e girobussola. La parte finale formalizza il modello dinamico del giroscopio a due gradi di liberta` e mostra come, bloccando un grado di liberta` e aggiungendo molla e smorzatore, si ottenga un \emph{rate gyro}: un giroscopio usato non per mantenere una direzione, ma per misurare direttamente una velocita` angolare.

\subsection{Misura della velocita` propria}

\slidefig{images7/page-01.png}{Misura della velocita` propria lineare o angolare, soprattutto in robotica mobile e veicoli.}

La prima slide chiarisce subito il cambio di scenario. Nella lezione precedente la velocita` era spesso una grandezza relativa: velocita` di un albero, di un carrello, di una ruota dentata, oppure velocita` stimata da impulsi di un encoder. Qui invece la domanda e`: quanto si sta muovendo il sistema stesso?

Per un robot mobile o un veicolo, la risposta puo` riguardare sia la velocita` lineare sia la velocita` angolare. La velocita` angolare di un drone, per esempio, entra direttamente nei loop di stabilizzazione di rollio, beccheggio e imbardata. La velocita` angolare di un satellite serve per il controllo d'assetto. La velocita` angolare di un'automobile e` fondamentale per sistemi di stabilita` e dinamica veicolo. La velocita` lineare di un aereo, invece, puo` essere riferita al terreno oppure all'aria: due grandezze diverse, entrambe importanti.

Il punto concettuale e` che, quando si misura la velocita` propria, non basta guardare un organo meccanico interno. Serve un sensore capace di fornire informazione inerziale, cioe` collegata al moto del corpo nello spazio. Il giroscopio nasce proprio per questo.

\subsection{Origine e ruolo dei giroscopi}

\slidepair{images7/page-02.png}{images7/page-03.png}{Cenni storici: dal giroscopio di Foucault alla girobussola navale di Anschuetz-Kaempfe e Sperry.}

Le slide storiche servono a capire perche' il giroscopio sia diventato cosi` importante nella navigazione. Foucault costruisce un primo giroscopio funzionale nel 1851, mostrando in modo sperimentale che un corpo in rapida rotazione tende a conservare la direzione del proprio asse. Questa proprieta` diventa immediatamente preziosa quando si vuole avere un riferimento direzionale indipendente da disturbi esterni.

All'inizio del Novecento il problema pratico diventa la navigazione navale, in particolare per sommergibili e grandi navi. Un ago magnetico puo` essere disturbato dalle masse ferromagnetiche, dalle correnti e dall'ambiente elettromagnetico; un giroscopio, invece, puo` fornire un riferimento legato alla rotazione terrestre e quindi al nord geografico. Anschuetz-Kaempfe e Sperry sviluppano soluzioni sempre piu` robuste, alimentate in modo continuo e montate su supporti capaci di sopportare il moto della nave.

Questa evoluzione storica mostra un aspetto che tornera` anche nella robotica moderna: il sensore non e` solo il principio fisico. E` anche meccanica di supporto, alimentazione, sospensione, smorzamento, compensazione degli errori e interfaccia con il sistema di controllo.

\slidefig{images7/page-04.png}{Definizione operativa di giroscopio e due principi fisici fondamentali: inerzia e precessione.}

La definizione essenziale e` semplice: un giroscopio contiene un rotore che gira rapidamente e il cui asse puo` orientarsi, almeno in parte, rispetto al supporto. In navigazione puo` essere usato in due modi complementari. Da un lato misura direttamente variabili di moto, soprattutto velocita` angolari. Dall'altro mantiene un riferimento inerziale rispetto al quale interpretare forze, assetto e direzioni del veicolo.

Le due parole chiave sono \emph{inerzia} e \emph{precessione}. Per inerzia giroscopica si intende la tendenza del rotore a mantenere la direzione del proprio asse di spin. Per precessione si intende invece la risposta apparentemente ``laterale'' del giroscopio: una coppia applicata non produce una rotazione direttamente nella direzione intuitiva della forza, ma cambia il momento angolare e genera una rotazione attorno a un asse perpendicolare.

\subsection{Inerzia giroscopica}

\slidepair{images7/page-05.png}{images7/page-06.png}{Principio di inerzia giroscopica: il rotore tende a mantenere il proprio piano di rotazione mentre il supporto si muove.}

Le due slide sull'inerzia mostrano il caso piu` intuitivo. Se la ruota gira rapidamente ed e` montata su giunti cardanici, il telaio esterno puo` essere ruotato senza trascinare immediatamente il piano della ruota. In assenza di coppie esterne significative, l'asse di spin tende a conservare la propria direzione nello spazio inerziale.

In termini di momento angolare, il rotore possiede
\[
\mathbf{L}=I\boldsymbol{\omega},
\]
dove \(I\) e` il momento d'inerzia rispetto all'asse di spin e \(\boldsymbol{\omega}\) e` la velocita` angolare del rotore. Se non agisce una coppia esterna, vale
\[
\boldsymbol{\tau}=\frac{d\mathbf{L}}{dt}\simeq 0,
\]
quindi \(\mathbf{L}\) resta quasi costante. Questa e` l'idea fisica che permette a un giroscopio meccanico di funzionare come riferimento di direzione.

Naturalmente, nei dispositivi reali non esiste il caso perfettamente ideale: attrito nei cuscinetti, sbilanciamenti, coppie elastiche dei supporti e vibrazioni introducono piccole variazioni di \(\mathbf{L}\). Pero` il principio rimane: maggiore e` il momento angolare del rotore, piu` il giroscopio e` resistente alle perturbazioni lente.

\subsection{Precessione: ruota appesa e trottola}

\slidetriple{images7/page-07.png}{images7/page-08.png}{images7/page-09.png}{Precessione di una ruota appesa: la coppia gravitazionale modifica la direzione del momento angolare.}

Il primo esempio di precessione e` la ruota di bicicletta appesa. Se la ruota non gira, il peso genera una coppia che la fa cadere. Se invece la ruota gira rapidamente, il peso genera comunque una coppia, ma il suo effetto principale non e` abbassare immediatamente l'asse: e` cambiare la direzione del momento angolare.

Indicando con \(l\) il braccio della forza peso rispetto al punto di sospensione, la coppia gravitazionale vale
\[
\tau=mgl.
\]
Questa coppia produce una variazione di momento angolare
\[
\Delta L=\tau\,\Delta t.
\]
Poiche' \(\Delta L\) e` perpendicolare a \(L\), il vettore momento angolare non cambia tanto modulo quanto direzione. Per piccoli angoli:
\[
\Delta\theta\simeq\frac{\Delta L}{L}.
\]
Dividendo per \(\Delta t\), si ottiene la velocita` angolare di precessione:
\[
\omega_p=\frac{\Delta\theta}{\Delta t}=\frac{\tau}{L}.
\]

Questa formula e` molto istruttiva. A parita` di coppia applicata, un rotore con momento angolare piu` grande precessa piu` lentamente. In pratica, se aumento la velocita` di spin o il momento d'inerzia, rendo il giroscopio piu` ``rigido'' rispetto alle coppie esterne.

\slidetriple{images7/page-10.png}{images7/page-11.png}{images7/page-12.png}{Precessione e nutazione della trottola: ruolo dell'angolo di inclinazione e oscillazioni attorno all'asse di simmetria.}

La trottola inclinata e` lo stesso fenomeno in una geometria leggermente diversa. La forza peso applicata al baricentro produce una coppia
\[
\tau=mgr\sin\phi,
\]
dove \(r\) e` la distanza del baricentro dal punto di contatto e \(\phi\) e` l'angolo di inclinazione rispetto alla verticale. In questo caso la rotazione di precessione interessa la componente \(L\sin\phi\), quindi
\[
\Delta\theta=\frac{\Delta L}{L\sin\phi}.
\]
La velocita` di precessione diventa
\[
\omega_p=\frac{\tau}{L\sin\phi}
       =\frac{mgr\sin\phi}{L\sin\phi}
       =\frac{mgr}{L}.
\]
Se \(L=I\omega\), allora
\[
\omega_p=\frac{mgr}{I\omega}.
\]

La dipendenza inversa da \(\omega\) spiega l'osservazione quotidiana: una trottola che gira molto rapidamente precessa lentamente e resta ``su'' in modo stabile; quando rallenta, la precessione diventa piu` evidente e il moto si degrada.

La slide sulla nutazione aggiunge un dettaglio importante. Durante la precessione, il momento angolare reale non e` sempre perfettamente allineato con l'asse di simmetria del corpo. Ne nasce un'oscillazione rapida attorno all'asse, chiamata nutazione. Nei giroscopi usati come sensori, la nutazione e` spesso un disturbo da smorzare; nei sistemi fisici liberi, invece, e` parte naturale della dinamica del corpo rigido.

\begin{keybox}
\textbf{Lettura fisica.} La precessione non e` una magia del giroscopio: e` semplicemente la legge \(\boldsymbol{\tau}=d\mathbf{L}/dt\) applicata a un vettore momento angolare molto grande. La coppia esterna cambia la direzione di \(\mathbf{L}\); la rotazione apparente dell'asse e` la conseguenza geometrica di questa variazione.
\end{keybox}

\subsection{Applicazioni alla navigazione aerea}

\slidepair{images7/page-13.png}{images7/page-14.png}{Indicatore di virata: la precessione del giroscopio viene trasformata in deflessione leggibile dal pilota.}

La navigazione aerea mostra molto bene come un principio fisico astratto diventi strumentazione concreta. L'indicatore di virata usa un giroscopio con asse e sospensioni elastiche. Quando l'aereo imbarda, il supporto del giroscopio ruota con il velivolo; per effetto della precessione, il giroscopio reagisce inclinando il proprio telaio. Questa inclinazione viene contrastata da molle e trasformata in una lettura sul quadrante.

La slide sottolinea anche la definizione pratica di \emph{standard rate turn}: una virata di \(3^\circ/\mathrm{s}\), che completa \(360^\circ\) in due minuti. Il punto di controllo e` importante: non basta sapere se l'aereo sta girando a sinistra o a destra, serve anche sapere con quale intensita`. Il rate gyro meccanico fornisce esattamente una grandezza proporzionale alla velocita` di virata.

\slidefig{images7/page-15.png}{Orizzonte artificiale e directional gyro: il giroscopio come riferimento di assetto e di direzione.}

L'orizzonte artificiale usa un giroscopio con asse orizzontale per mantenere un riferimento stabile rispetto al quale il velivolo puo` inclinarsi. Nella visualizzazione sembra che sia il giroscopio a inclinarsi, ma fisicamente e` il frame fissato all'aereo che ruota attorno al rotore, mentre il rotore tende a conservare la propria direzione.

Il directional gyro affronta invece il problema della direzione di prua. Una girobussola navale classica non e` adatta a un aereo, perche' il velivolo cambia rapidamente quota, rollio e beccheggio, accumulando errori se il riferimento non viene stabilizzato. Per questo il directional gyro viene montato su piattaforma stabilizzata e riallineato periodicamente o continuamente tramite bussola magnetica e motori di correzione. Il risultato e` uno strumento che fornisce al pilota una lettura rapida della direzione corrente.

\subsection{Girobussola e ricerca del nord}

\slidetriple{images7/page-16.png}{images7/page-17.png}{images7/page-18.png}{Giroscopio libero rispetto alla Terra: moto apparente diverso a equatore, poli e latitudini intermedie.}

Le slide sulla girobussola introducono un'idea sottile: un giroscopio libero conserva la propria direzione inerziale, ma l'osservatore e` solidale con la Terra, che ruota. Di conseguenza, il giroscopio sembra ruotare rispetto al riferimento terrestre.

All'equatore, un giroscopio libero appare compiere una rotazione completa in 24 ore attorno all'asse orizzontale nord-sud. Ai poli, invece, la rotazione apparente e` attorno all'asse verticale. A latitudini intermedie compaiono entrambe le componenti: una orizzontale e una verticale. Il comportamento osservato dipende quindi dalla decomposizione del vettore velocita` angolare terrestre nel sistema locale.

Questa osservazione e` la base della girobussola: se si riesce a vincolare e smorzare opportunamente il giroscopio, la sua dinamica puo` essere fatta convergere verso l'allineamento con l'asse di rotazione terrestre, cioe` verso il nord geografico.

\slidetriple{images7/page-19.png}{images7/page-20.png}{images7/page-21.png}{Principio della girobussola: una coppia di richiamo genera precessione fino all'allineamento con la rotazione terrestre.}

La soluzione proposta nelle slide e` introdurre una sospensione capace di generare coppie di richiamo quando il giroscopio si allontana dalla verticale. Storicamente si potevano usare, ad esempio, tubi di mercurio: lo spostamento del liquido produceva una coppia gravitazionale di richiamo.

La coppia \(\tau\) non riallinea il giroscopio in modo diretto, ma lo fa precessare con una velocita` \(\Omega\). La precessione continua finche' l'asse di spin \(\omega\) del giroscopio non si allinea con il vettore velocita` angolare terrestre. Quando l'allineamento e` raggiunto, la coppia media necessaria alla correzione si annulla e il sistema si stabilizza.

Questo e` il motivo per cui una girobussola indica il nord vero, non semplicemente il nord magnetico. Naturalmente nella pratica servono smorzamento, compensazioni e correzioni dinamiche: senza smorzamento, il sistema tenderebbe a oscillare attorno alla direzione corretta; con smorzamento opportuno, invece, converge.

\subsection{Modello matematico del giroscopio}

\slidepair{images7/page-22.png}{images7/page-23.png}{Configurazione a due gradi di liberta`: rotore, gimbal interno, gimbal esterno e velocita` angolari relative.}

La parte matematica formalizza il dispositivo come corpo rigido rotante montato su due giunti. Il rotore \(r\) ruota a velocita` costante \(n\) rispetto al gimbal interno \(g\). L'angolo \(\psi\) descrive la posizione del gimbal esterno rispetto al riferimento inerziale o al veicolo, mentre \(\theta\) descrive la posizione del gimbal interno rispetto al gimbal esterno.

Le velocita` angolari delle terne relative possono essere scritte come
\[
\Omega^\beta=\mathbf{1}_Z\dot\psi,
\]
\[
\Omega^\alpha
=\mathbf{1}_y\dot\theta+\mathbf{1}_Z\dot\psi
=\mathbf{1}_x(-\dot\psi\sin\theta)
 +\mathbf{1}_y\dot\theta
 +\mathbf{1}_z(\dot\psi\cos\theta),
\]
mentre per il rotore:
\[
\Omega^r
=\mathbf{1}_x n+\mathbf{1}_y\dot\theta+\mathbf{1}_Z\dot\psi
=\mathbf{1}_x(n-\dot\psi\sin\theta)
 +\mathbf{1}_y\dot\theta
 +\mathbf{1}_z(\dot\psi\cos\theta).
\]

Queste formule dicono una cosa molto concreta: anche se il rotore gira principalmente con velocita` \(n\), le rotazioni dei gimbal aggiungono componenti di velocita` angolare sugli altri assi. Sono proprio queste componenti accoppiate a generare le coppie giroscopiche.

\slidepair{images7/page-24.png}{images7/page-25.png}{Equazioni dinamiche tramite seconda legge di Newton per corpi rotanti e bilanci di momento.}

Le slide successive applicano la seconda legge di Newton in forma rotazionale. Per un corpo espresso in una terna rotante, la derivata inerziale del momento angolare contiene sia la derivata vista nella terna sia il termine di trasporto:
\[
\left(\frac{d\mathbf{H}}{dt}\right)_{\!i}
=
\left(\frac{d\mathbf{H}}{dt}\right)_{\!\alpha}
+\Omega^\alpha\times\mathbf{H}.
\]
Questo termine di prodotto vettoriale e` il cuore della dinamica giroscopica: quando la terna ruota, un momento angolare grande lungo un asse produce coppie apparenti sugli altri assi.

La notazione delle slide introduce \(M_b\) come coppia attorno all'asse \(Y\) tra i gimbal e \(M_c\) come coppia attorno all'asse \(Z\) esercitata dalla base sul gimbal esterno. Inoltre, ponendo
\[
h=J_x^r n,
\]
si identifica il momento angolare di spin del rotore rispetto al gimbal interno. Se \(n\) e` costante e molto maggiore delle altre velocita`, \(h\) diventa il parametro dominante del sensore.

\slidepair{images7/page-26.png}{images7/page-27.png}{Equazioni non lineari finali e linearizzazione per piccoli moti angolari.}

Combinando i bilanci di momento, le slide arrivano a un modello non lineare del tipo
\[
\begin{cases}
J_y\ddot\theta+h\dot\psi\cos\theta
+(J_z-J_x)\dot\psi^2\sin\theta\cos\theta=M_b,\\[2mm]
\left(J_z^o+J_z\cos^2\theta+J_x\sin^2\theta\right)\ddot\psi
-h\dot\theta\cos\theta
+2(J_x-J_z)\dot\psi\dot\theta\sin\theta\cos\theta=M_c.
\end{cases}
\]

La forma completa e` utile per capire da dove arrivano gli accoppiamenti, ma in molte applicazioni il giroscopio lavora con piccoli moti angolari: milliradianti o secondi d'arco. Inoltre la velocita` di spin \(n\) e` molto piu` grande delle velocita` dei gimbal. In queste condizioni si linearizza attorno a \(\theta\simeq 0\), trascurando termini quadratici:
\[
\begin{cases}
J_y\ddot\theta+h\dot\psi=M_b,\\
J_Z\ddot\psi-h\dot\theta=M_c.
\end{cases}
\]

Queste due equazioni sono molto piu` leggibili. La prima dice che una velocita` \(\dot\psi\) del case genera una coppia equivalente sull'asse \(\theta\), proporzionale a \(h\). La seconda dice il duale: una velocita` \(\dot\theta\) genera una coppia sull'asse \(\psi\). Il giroscopio e` quindi, per natura, un dispositivo di accoppiamento tra assi.

\subsection{Rate gyro: giroscopio per misurare velocita` angolare}

\slidepair{images7/page-28.png}{images7/page-29.png}{Dal modello a un grado di liberta` al rate gyro passivo con molla, smorzatore e pick-off angolare.}

Per trasformare il giroscopio in un sensore di velocita` angolare si blocca il gimbal esterno al veicolo e si lascia libero soltanto il gimbal interno. In questo caso \(\dot\psi\) non e` piu` una variabile da stimare internamente: e` l'ingresso del sistema, cioe` la velocita` angolare del case. Il modello a un grado di liberta` diventa
\[
J_y\ddot\theta=M_b-h\dot\psi.
\]

Nel rate gyro meccanico passivo si aggiungono una molla di richiamo e uno smorzatore viscoso. Con una coppia incerta o disturbante \(M_u\), la coppia sul gimbal si modella come
\[
M_b=-k\theta-b\dot\theta+M_u.
\]
Sostituendo:
\[
J_y\ddot\theta+b\dot\theta+k\theta=M_u-h\dot\psi.
\]

Se consideriamo il regime statico, trascuriamo accelerazione e velocita` del gimbal e assumiamo \(M_u\simeq 0\), resta
\[
k\theta=-h\dot\psi,
\qquad
\theta=-\frac{h}{k}\dot\psi.
\]

La lettura e` potentissima: l'uscita \(\theta\), cioe` la rotazione del gimbal misurata da un pick-off angolare, e` proporzionale alla velocita` angolare del case. La costante di sensibilita` e` \(h/k\): aumentando il momento angolare di spin \(h\) aumenta la sensibilita`; aumentando la rigidezza \(k\) si riduce la deflessione ma si rende il sensore piu` robusto e rapido.

Lo smorzatore \(b\) non cambia la relazione statica ideale, ma e` fondamentale nella dinamica: limita oscillazioni, nutazione e risonanze. In un sensore reale, quindi, il progetto di \(k\), \(b\), \(J_y\) e \(h\) determina banda, sensibilita`, rumore meccanico e tempo di assestamento.

\begin{keybox}
\textbf{Idea pratica del rate gyro.} Un giroscopio libero mantiene una direzione; un rate gyro, invece, trasforma la coppia giroscopica prodotta dalla velocita` del case in una piccola deflessione elastica misurabile. Non misura direttamente un angolo assoluto di assetto: misura una velocita` angolare.
\end{keybox}

\subsection*{Sintesi finale della Lezione 8}

Questa ottava lezione collega la misura di velocita` alla navigazione inerziale. Il giroscopio nasce dalla conservazione del momento angolare: un rotore veloce tende a mantenere la direzione del proprio asse e risponde alle coppie esterne con precessione. La stessa fisica spiega la ruota appesa, la trottola, la nutazione e gli strumenti classici di navigazione aerea.

Dal punto di vista del controllo, la parte piu` importante e` il passaggio finale: il giroscopio non e` soltanto un riferimento direzionale, ma puo` diventare un sensore di velocita` angolare. Nel rate gyro, la velocita` del case \(\dot\psi\) genera una coppia giroscopica proporzionale al momento angolare di spin \(h\); molla e smorzatore la trasformano in una deflessione \(\theta\) leggibile. In questo modo si ottiene una misura diretta, continua e fisicamente fondata della velocita` angolare propria del veicolo.

% =========================================================
% LEZIONE 9
% =========================================================

\clearpage
\section{Lezione 9}

\subsection{Visione generale della lezione}

Questa lezione chiude il discorso sui giroscopi meccanici e apre quello sui giroscopi ottici. La prima parte riprende il modello linearizzato della Lezione 8 e lo usa per capire due fenomeni dinamici che non si vedono nella formula statica del rate gyro: il \emph{coning}, cioe` il moto conico dell'asse di spin, e la distinzione tra precessione media e nutazione sovrapposta. Poi si mostra come il rate gyro possa essere trasformato da sensore passivo con molla a sensore a retroazione elettronica, in cui una coppia applicata da un motore mantiene quasi nullo l'angolo di precessione.

La seconda parte cambia tecnologia: invece di misurare la velocita` angolare tramite un rotore meccanico, si usa la propagazione della luce in due direzioni opposte. Il principio comune e` l'effetto Sagnac: se un cammino ottico chiuso ruota, i due fasci contro-propaganti impiegano tempi leggermente diversi per completare il giro. Questa differenza puo` essere letta come differenza di fase nei \emph{Fiber-Optic Gyro} (FOG) oppure come differenza di frequenza nei \emph{Ring Laser Gyro} (RLG).

\subsection{Dinamica linearizzata: poli e coning}

\slidefig{images8/page-01.png}{Modello linearizzato del giroscopio a due assi, poli del sistema e frequenza di oscillazione non smorzata.}

La slide riprende il modello linearizzato del giroscopio, con una notazione leggermente diversa:
\[
\begin{cases}
J_y\ddot\theta+h\dot\psi=M_y,\\
J_z\ddot\psi-h\dot\theta=M_z.
\end{cases}
\]
I termini \(J_y\ddot\theta\) e \(J_z\ddot\psi\) sono le inerzie dei due assi di gimbal, mentre i termini \(h\dot\psi\) e \(-h\dot\theta\) sono gli accoppiamenti giroscopici. Il parametro \(h\) resta il momento angolare di spin del rotore:
\[
h=I_{rx}n,
\]
dove \(I_{rx}\) e` il momento d'inerzia del rotore attorno all'asse di spin e \(n\) e` la velocita` di spin.

Il modello non smorzato presenta due poli nell'origine e una coppia di poli immaginari:
\[
s=0,0,\pm j\frac{h}{\sqrt{J_yJ_z}}.
\]
La frequenza naturale delle oscillazioni giroscopiche e`
\[
\omega_0=\frac{h}{\sqrt{J_yJ_z}}
=\frac{I_{rx}n}{\sqrt{J_yJ_z}}.
\]
Nel caso ideale di rotore a disco sottile e gimbal senza massa, la slide assume
\[
J_y=J_z=\frac{I_{rx}}{2},
\]
e quindi
\[
\omega_0=2n.
\]

Il messaggio fisico e` molto chiaro: oltre al moto lento di precessione, il giroscopio puo` avere un'oscillazione propria molto rapida, legata alla velocita` del rotore. In un sensore reale questa oscillazione va smorzata o tenuta fuori dalla banda utile.

\slidepair{images8/page-02.png}{images8/page-03.png}{Risposte a impulso e a coppia costante: coning, precessione media, nutazione e drift giroscopico.}

Se si applica un impulso di coppia \(M_y(t)=m\delta(t)\), la risposta del giroscopio e` oscillatoria:
\[
\theta(t)=\frac{\Omega_{y0}}{\omega_0}\sin(\omega_0t)\,u(t),
\qquad
\Omega_{y0}=\frac{m}{J_y},
\]
\[
\psi(t)=
\frac{\Omega_{y0}}{\omega_0}
\sqrt{\frac{J_y}{J_z}}
\left(1-\cos(\omega_0t)\right)u(t).
\]
Il moto risultante dell'asse di spin e` chiamato \emph{coning}, perche' l'asse descrive una piccola traiettoria conica. Se i momenti d'inerzia sui due assi sono uguali, il coning diventa circolare. Se invece c'e` smorzamento nei gimbal, il moto conico decade nel tempo.

La coppia costante \(M_{y0}\) mette in evidenza un altro aspetto. Le slide riportano:
\[
\theta(t)=\frac{M_{y0}}{J_y\omega_0^2}
\left(1-\cos(\omega_0t)\right)u(t),
\]
\[
\psi(t)=\frac{M_{y0}}{h}
\left(t-\frac{\sin(\omega_0t)}{\omega_0}\right)u(t).
\]
La parte lineare in \(t\) e` la precessione media, mentre il termine sinusoidale e` la nutazione sovrapposta. La velocita` di precessione stazionaria vale
\[
\dot\psi_{\text{ss}}=\frac{M_{y0}}{h}.
\]

Questa relazione e` una delle piu` importanti dell'intero blocco: una coppia costante non compensata produce un drift giroscopico. Aumentare \(h\) riduce il drift, ma non lo elimina se esistono coppie parassite permanenti.

\begin{keybox}
\textbf{Idea dinamica.} La formula statica del rate gyro nasconde una dinamica ricca: coning, nutazione e drift. Per usare bene un giroscopio nel controllo non basta conoscere la sensibilita` statica; bisogna anche sapere dove cadono le oscillazioni e quanto rapidamente vengono smorzate.
\end{keybox}

\subsection{Rate gyro con controllo elettronico}

\slidepair{images8/page-04.png}{images8/page-05.png}{Rate gyro a controllo elettronico: il torquer sostituisce la molla e mantiene quasi nullo l'angolo di precessione.}

Nel rate gyro passivo della Lezione 8, la coppia di richiamo era generata da una molla. Qui l'idea cambia: si usa un motore, o \emph{torquer}, per applicare la coppia necessaria a mantenere l'angolo \(\theta\) piccolo, idealmente nullo. Si parte ancora dall'equazione
\[
J_y\ddot\theta=M_b-h\dot\psi,
\]
che in Laplace puo` essere letta come una pianta doppio integratore:
\[
\theta(s)=G(s)\left(M_b(s)-h\dot\psi(s)\right),
\qquad
G(s)=\frac{1}{J_y s^2}.
\]

Il controllore calcola la coppia
\[
M_b=K(s)\left(\theta_{\text{REF}}-\theta\right).
\]
Gli obiettivi sono due: regolare \(\theta\) verso zero e, nello stesso tempo, stimare la perturbazione che sta cercando di far muovere il gimbal, cioe` la velocita` angolare del case \(\dot\psi\).

La funzione di trasferimento dalla reference all'angolo e`
\[
\frac{\theta}{\theta_{\text{REF}}}
=\frac{G(s)K(s)}{1+G(s)K(s)}.
\]
Ponendo \(\theta_{\text{REF}}=0\), la coppia richiesta al torquer diventa
\[
\frac{M_b}{\dot\psi}
=-h\,\frac{G(s)K(s)}{1+G(s)K(s)}.
\]
Se il loop ha guadagno elevato nella banda di interesse,
\[
G(s)K(s)\gg 1
\quad\Rightarrow\quad
M_b\simeq -h\dot\psi.
\]

Quindi il segnale utile non e` piu` la deflessione \(\theta\), che viene tenuta quasi a zero, ma la coppia elettrica necessaria ad annullarla. Questo principio si chiama spesso \emph{force rebalance}: si misura una grandezza osservando quanta azione di controllo serve per tenerla in equilibrio.

\subsection{Esempi di girobussola e tecnologie meccaniche}

\slidepair{images8/page-06.png}{images8/page-07.png}{Esempi di gyrocompass e dimostrazioni fisiche del comportamento giroscopico.}

Le slide di esempio sulla gyrocompass servono come ponte tra teoria e strumentazione reale. Da un lato si vedono strumenti storici di navigazione, in cui la meccanica giroscopica viene trasformata in indicazione di direzione; dall'altro si vede una dimostrazione fisica del comportamento interno, con un elemento che reagisce alla rotazione del supporto.

Il punto da ricordare e` che la girobussola e il rate gyro usano la stessa fisica, ma con scopi diversi. La girobussola cerca una direzione di equilibrio legata alla rotazione terrestre; il rate gyro cerca una misura istantanea della velocita` angolare del case. Nel primo caso conta la convergenza direzionale, nel secondo conta la banda di misura.

\slidefig{images8/page-08.png}{Tecnologie meccaniche per giroscopi: single-axis gyro flottante e giroscopio sferico sospeso elettricamente.}

La slide sulle tecnologie meccaniche mostra due soluzioni classiche per ridurre attriti e coppie parassite. Nel \emph{floating single-axis gyro}, il gimbal viene immerso o supportato in modo da alleggerire i carichi sui cuscinetti e ridurre l'attrito. Nel giroscopio sferico sospeso elettricamente, invece, il rotore e` mantenuto in sospensione senza contatto meccanico diretto, spesso in vuoto, e la lettura puo` essere ottica o capacitiva.

Queste soluzioni nascono dalla stessa esigenza: se una coppia parassita produce drift, allora bisogna progettare il sensore in modo da minimizzare attrito, sbilanciamento e contatti. La complessita` meccanica dei giroscopi ad alte prestazioni e` quindi una conseguenza diretta della loro sensibilita` alle coppie indesiderate.

\subsection{Effetto Sagnac: base dei giroscopi ottici}

\slidetriple{images8/page-09.png}{images8/page-10.png}{images8/page-11.png}{Effetto Sagnac su cammino circolare: due fasci contro-propaganti percorrono lunghezze efficaci diverse se il sistema ruota.}

La tecnologia ottica elimina il rotore meccanico e misura la velocita` angolare attraverso la luce. Due fasci laser generati dalla stessa sorgente percorrono lo stesso cammino chiuso in direzioni opposte: uno orario e uno antiorario. Se il sistema e` fermo, i due tempi di percorrenza sono uguali. Se il sistema ruota con velocita` angolare \(\omega\), il cammino efficace di un fascio si allunga e quello dell'altro si accorcia.

Per un cammino circolare di raggio \(R\), con area racchiusa \(A=\pi R^2\), le slide ricavano i tempi
\[
t^+=\frac{2\pi R}{c(1-\omega R/c)},
\qquad
t^-=\frac{2\pi R}{c(1+\omega R/c)}.
\]
Poiche' nei casi reali \(\omega R/c\ll 1\), si usa l'approssimazione
\[
\frac{1}{1-x}\simeq 1+x.
\]
La differenza di tempo diventa
\[
\Delta t=t^+-t^-
\simeq \frac{4\pi R^2}{c^2}\omega
=\frac{4A}{c^2}\omega.
\]
Moltiplicando per \(c\), si ottiene la differenza di cammino ottico:
\[
\Delta L=c\Delta t=\frac{4A}{c}\omega.
\]

Questa e` la formula centrale dell'effetto Sagnac. La misura cresce con l'area racchiusa dal cammino ottico e con la velocita` angolare. Il segnale e` piccolissimo, ma ha un vantaggio enorme: non richiede parti meccaniche in rotazione.

\slidefig{images8/page-12.png}{Risultato compatto dell'effetto Sagnac: differenza di tempo e differenza di cammino ottico proporzionali all'area racchiusa.}

La slide riassume il risultato nella forma piu` utile:
\[
\Delta t=\frac{4A}{c^2}\omega,
\qquad
\Delta L=\frac{4A}{c}\omega.
\]
Queste relazioni sono valide, al primo ordine, anche per cammini non perfettamente circolari: cio` che conta e` l'area orientata racchiusa dal percorso, non il fatto che il percorso sia una circonferenza ideale.

\subsection{Fiber-Optic Gyro: misura di fase}

\slidefig{images8/page-13.png}{Fiber-Optic Gyro: la differenza di cammino diventa differenza di fase interferometrica.}

Nel Fiber-Optic Gyro (FOG), i due fasci percorrono una bobina di fibra ottica in direzioni opposte e poi vengono ricombinati. La differenza di cammino \(\Delta L\) produce una differenza di fase:
\[
\Delta\phi=2\pi\frac{\Delta L}{\lambda}.
\]
Sostituendo il risultato Sagnac:
\[
\Delta\phi
=2\pi\frac{1}{\lambda}\frac{4A}{c}\omega
=\frac{8\pi A}{c\lambda}\omega.
\]

Una singola spira avrebbe un segnale molto piccolo. Per aumentare la sensibilita`, si avvolge la fibra in \(N\) spire:
\[
\Delta\phi=\frac{8\pi A N}{c\lambda}\omega.
\]
Per una bobina di raggio \(R\) e lunghezza totale di fibra \(L\), la stessa formula viene scritta nelle slide come
\[
\Delta\phi=\frac{4\pi R L}{c\lambda}\omega.
\]
Il significato e` immediato: piu` fibra e piu` area ottica racchiusa aumentano la sensibilita` del sensore.

\slidepair{images8/page-14.png}{images8/page-15.png}{Interferenza: la differenza di fase viene convertita in variazione di intensita` luminosa e letta in quadratura.}

Per misurare \(\Delta\phi\), il FOG sfrutta l'interferenza. Due onde coerenti che arrivano in fase producono interferenza costruttiva e quindi una zona piu` luminosa; due onde sfasate di \(180^\circ\) producono interferenza distruttiva e quindi una zona scura. L'intensita` luminosa osservata contiene quindi informazione sulla differenza di fase.

La slide finale del blocco aggiunge un'analogia importante con gli encoder: per capire anche il segno della velocita`, non basta una sola misura di intensita`, perche' attorno a certi punti la risposta puo` essere ambigua. Si usano allora due letture sfasate di un quarto di frangia, ottenendo segnali in quadratura, simili a una coppia seno-coseno. In questo modo si ricava sia il modulo sia il verso della velocita` angolare.

\slidefig{images8/page-16.png}{Effetto Sagnac su cammino non circolare: lo stesso risultato dipende dall'area racchiusa dal percorso ottico.}

La slide sul cammino quadrato mostra che la formula Sagnac non e` limitata all'anello circolare. Tre specchi e un beam splitter possono formare un cammino chiuso non circolare; per piccoli angoli e al primo ordine, la differenza di tempo rimane
\[
\Delta t=\frac{4A}{c^2}\omega,
\]
e quindi
\[
\Delta L=\frac{4A}{c}\omega.
\]
La geometria cambia, ma l'informazione fisica resta la stessa: la rotazione modifica diversamente i tempi di percorrenza dei due fasci contro-propaganti.

\subsection{Ring Laser Gyro: misura di frequenza}

\slidetriple{images8/page-17.png}{images8/page-18.png}{images8/page-19.png}{Ring Laser Gyro: cavita` laser, onde stazionarie e conversione dell'effetto Sagnac in battimento di frequenza.}

Nel Ring Laser Gyro (RLG), il cammino ottico non e` solo una fibra: e` una cavita` laser chiusa, tipicamente formata da specchi. I due fasci contro-propaganti sono modi della cavita`; quando il sistema ruota, l'effetto Sagnac modifica le lunghezze ottiche efficaci e quindi le frequenze dei due modi.

Le slide elencano i principali vantaggi dei laser gyro rispetto ai giroscopi meccanici: assenza di parti mobili, progetto relativamente robusto, insensibilita` a vibrazioni \(g\) e \(g^2\), ampio range dinamico, uscita naturalmente digitale, aggiornamento rapido e lunga vita operativa.

Il passaggio matematico parte da
\[
\nu\lambda=c.
\]
Per piccole variazioni:
\[
(\nu+\Delta\nu)(\lambda+\Delta\lambda)\simeq \nu\lambda,
\]
da cui
\[
\Delta\nu=-\frac{\Delta\lambda}{\lambda}\nu.
\]

La condizione di risonanza della cavita` impone che la lunghezza ottica contenga un numero intero \(p\) di lunghezze d'onda:
\[
L=p\lambda.
\]
Se la rotazione produce una variazione di lunghezza \(\Delta L\), allora
\[
L+\Delta L=p(\lambda+\Delta\lambda),
\qquad
\Delta L=p\Delta\lambda=\frac{4A}{c}\omega.
\]
Sostituendo nella relazione di frequenza, le slide arrivano a
\[
\Delta\nu=-4\frac{A\nu}{cL}\omega.
\]

Questa e` la misura del RLG: la velocita` angolare viene trasformata in una differenza di frequenza tra i due fasci.

\slidepair{images8/page-20.png}{images8/page-21.png}{Nel RLG fermo le frange sono stazionarie; in rotazione si muovono alla frequenza di battimento.}

Quando il RLG e` fermo, i due fasci formano al rivelatore un pattern di frange stazionario. Quando il sistema ruota, le frequenze dei fasci diventano leggermente diverse e nasce un battimento. Le frange si muovono sul detector con una frequenza proporzionale a
\[
\Delta\nu=-4\frac{A\nu}{cL}\omega.
\]
Contare i passaggi delle frange per unita` di tempo significa quindi misurare direttamente \(\omega\).

Questo spiega perche' l'uscita del RLG venga spesso descritta come intrinsecamente digitale: l'informazione puo` essere letta come conteggio di frange o frequenza di battimento, non come piccolissima deflessione meccanica.

\slidepair{images8/page-22.png}{images8/page-23.png}{Esempi reali di RLG triassiale e confronto applicativo tra Ring Laser Gyro e Fiber-Optic Gyro.}

Le ultime slide mostrano esempi reali. Un RLG triassiale contiene tre cavita` orientate lungo assi diversi, in modo da misurare le tre componenti della velocita` angolare. Un sistema FOG puo` essere usato in girobussole e sistemi di riferimento d'assetto, combinando piu` assi ottici con elettronica di stima e filtraggio.

FOG e RLG risolvono lo stesso problema con due letture diverse dell'effetto Sagnac. Nel FOG si misura una differenza di fase interferometrica; nel RLG si misura una differenza di frequenza. Entrambi eliminano il rotore meccanico, riducendo usura e attrito, ma richiedono ottica, sorgenti luminose, rivelatori e compensazioni sofisticate.

\begin{keybox}
\textbf{Confronto finale.} Il giroscopio meccanico misura la rotazione attraverso momento angolare e coppie giroscopiche. Il FOG misura una fase ottica proporzionale a \(\omega\). Il RLG misura una frequenza di battimento proporzionale a \(\omega\). Cambia la tecnologia, ma l'obiettivo resta identico: ottenere una misura robusta della velocita` angolare propria.
\end{keybox}

\subsection*{Sintesi finale della Lezione 9}

La nona lezione completa il percorso dai giroscopi meccanici ai giroscopi ottici. Dal lato meccanico, il modello linearizzato mostra che un giroscopio non e` solo un trasduttore statico: ha poli, oscillazioni, coning, nutazione e drift. Il controllo elettronico permette di migliorare il rate gyro tenendo quasi nullo l'angolo di precessione e usando la coppia del torquer come misura della velocita` angolare.

Dal lato ottico, l'effetto Sagnac sostituisce completamente la meccanica rotante. In un cammino chiuso in rotazione, due fasci contro-propaganti accumulano tempi di percorrenza diversi:
\[
\Delta t=\frac{4A}{c^2}\omega.
\]
Il FOG legge questa differenza come fase, il RLG come frequenza. Questa transizione e` molto importante per la robotica e la navigazione moderne: la grandezza fisica da misurare e` sempre \(\omega\), ma la tecnologia passa da masse rotanti e cuscinetti a luce, interferenza e conteggio di frange.

% =========================================================
% LEZIONE 10
% =========================================================

\clearpage
\section{Lezione 10}

\subsection{Visione generale della lezione}

Questa lezione completa il blocco sui giroscopi introducendo i \emph{vibrating element gyroscopes}, cioe` i giroscopi in cui non c'e` piu` un rotore macroscopico in spin, ma una microstruttura tenuta in vibrazione. La rotazione del corpo genera una forza di Coriolis proporzionale alla velocita` di vibrazione e alla velocita` angolare da misurare; leggendo lo spostamento laterale della massa si ottiene quindi una misura di \(\Omega\). Questo e` il principio fisico alla base di molti giroscopi MEMS.

La lezione poi confronta le principali tecnologie di giroscopi: meccanici classici, DTG, MEMS vibranti, FOG, RLG e HRG. Il confronto non e` solo una classifica di precisione: ogni tecnologia occupa una zona diversa del compromesso tra costo, dimensioni, consumo, rumore, deriva e robustezza.

Nella parte finale si passa dalla velocita` angolare alla velocita` lineare propria del veicolo. Per un veicolo terrestre si puo` stimare la velocita` dalle ruote, con il limite fondamentale dello slittamento. Per un aereo, invece, la grandezza piu` utile per il controllo e` spesso la velocita` rispetto all'aria, misurata con il tubo di Pitot tramite il principio di Bernoulli.

\subsection{Giroscopi a elemento vibrante}

\slidepair{images9/part1-01.png}{images9/part1-02.png}{Giroscopio vibrante elementare e giroscopio al quarzo: una vibrazione nota viene deviata dalla forza di Coriolis.}

L'idea dei giroscopi vibranti e` sostituire il rotore in spin con una massa che oscilla in modo controllato. La massa viene forzata a vibrare lungo l'asse \(x\):
\[
x(t)=X\sin(\omega_d t),
\qquad
\dot x(t)=X\omega_d\cos(\omega_d t),
\]
dove \(X\) e` l'ampiezza di vibrazione e \(\omega_d\) e` la frequenza di drive.

Se il sensore ruota con velocita` angolare \(\Omega_z\), la massa, che ha velocita` lungo \(x\), subisce una accelerazione di Coriolis lungo \(y\). Con la convenzione di segno usata nella slide:
\[
a_y=-2\Omega_z\dot x
=-2\Omega_zX\omega_d\cos(\omega_d t).
\]
In forma vettoriale il concetto e`
\[
\mathbf a_C=-2\boldsymbol{\Omega}\times\mathbf v_{\text{rel}},
\]
oppure, per la forza sulla massa,
\[
\mathbf F_C=m\mathbf a_C.
\]

Il sensore quindi lavora in due direzioni distinte. L'asse di \emph{drive} serve a creare una velocita` nota e periodica; l'asse di \emph{sense} misura lo spostamento laterale indotto dalla rotazione. Se non c'e` rotazione, la massa vibra solo lungo \(x\). Se c'e` rotazione attorno a \(z\), compare una componente di moto lungo \(y\), alla stessa frequenza del drive e con ampiezza proporzionale a \(\Omega_z\).

Nel giroscopio al quarzo la stessa logica viene realizzata con un cristallo prismatico mantenuto in flessione da trasduttori piezoelettrici. La rotazione attorno all'asse longitudinale altera la direzione effettiva della vibrazione: i piezoelementi laterali ricevono segnali differenti e questa differenza viene usata per stimare la velocita` angolare.

\begin{keybox}
\textbf{Idea chiave.} Un giroscopio vibrante non misura direttamente la rotazione come spostamento statico: misura una vibrazione secondaria generata dalla forza di Coriolis. Per questo e` fondamentale conoscere fase, frequenza e ampiezza della vibrazione primaria.
\end{keybox}

\slidepair{images9/part1-03.png}{images9/part1-04.png}{Esempi di giroscopi vibranti e tuning fork gyroscope.}

Il \emph{tuning fork gyroscope} usa due masse che vibrano in opposizione, come i rebbi di un diapason. Questa scelta non e` estetica: facendo muovere due masse in controfase si riducono le vibrazioni trasmesse al supporto e si migliora il rigetto dei disturbi comuni. Quando il dispositivo ruota attorno all'asse sensibile, le forze di Coriolis sulle due masse hanno verso coerente rispetto al modo di sense, e quindi producono un segnale differenziale leggibile.

La separazione tra modo di drive e modo di sense e` un tema ricorrente nei MEMS. Un buon progetto deve eccitare molto bene il modo utile, ma deve evitare che accelerazioni lineari, vibrazioni esterne o errori di fabbricazione generino un falso segnale di rotazione.

\subsection{Attuazione e lettura capacitiva nei MEMS}

\slidepair{images9/part1-05.png}{images9/part1-06.png}{Comb drive e richiamo elettrostatico: le strutture capacitive applicano forze alle micro-masse.}

Nei MEMS la massa vibrante viene spesso mossa e letta con strutture capacitive. La slide sul tuning fork con \emph{comb drive} mostra due funzioni diverse:
\[
\text{drive: creare la vibrazione primaria},\qquad
\text{sense: misurare lo spostamento da Coriolis}.
\]
I pettini elettrostatici applicano una forza controllata alla massa; altri elettrodi leggono lo spostamento tramite variazioni di capacita`.

La forza elettrostatica si capisce partendo dall'energia immagazzinata in un condensatore:
\[
U_c=\frac{1}{2}QV=\frac{1}{2}CV^2.
\]
Se la geometria cambia con la coordinata \(x\), anche la capacita` cambia con \(x\), e quindi l'energia dipende dalla posizione. La forza generalizzata lungo \(x\) e`
\[
F_x=-\frac{\partial U_c}{\partial x}.
\]
Il segno dipende dalla convenzione scelta per \(x\); fisicamente il sistema tende a muoversi nella direzione che aumenta la capacita` e riduce l'energia potenziale elettrostatica equivalente.

\slidepair{images9/part1-07.png}{images9/part1-08.png}{Condensatore piano e comb drive: dalla capacita` alla forza lungo la direzione cedevole.}

Per un condensatore piano, trascurando gli effetti di bordo, vale
\[
C=k\frac{\varepsilon_0 A}{z}
=k\frac{\varepsilon_0 WX}{z},
\]
dove \(W X\) e` l'area sovrapposta, \(z\) e` la distanza tra le armature, \(\varepsilon_0\) e` la permittivita` del vuoto e \(k\) tiene conto del dielettrico. L'energia e`
\[
U_c=\frac{1}{2}CV^2
=\frac{1}{2}\frac{k\varepsilon_0WX}{z}V^2.
\]
Quindi, se \(X\) cambia, la forza lungo la direzione cedevole e`
\[
F_x=-\frac{\partial U_c}{\partial X}
=-\frac{1}{2}\frac{k\varepsilon_0W}{z}V^2.
\]

Nel \emph{comb drive} si mettono molti condensatori in parallelo, ottenendo una capacita` totale approssimabile come
\[
C_{\text{tot}}=N\varepsilon_0\frac{WX}{z}.
\]
La forza diventa
\[
F_x=-\frac{\partial U_c}{\partial X}
=-N\frac{1}{2}\frac{k\varepsilon_0W}{z}V^2.
\]
Il vantaggio e` pratico: con molte dita interdigitate si ottiene una forza significativa anche con spostamenti micrometrici. Inoltre, nel tratto utile, la forza e` quasi indipendente dalla sovrapposizione \(X\), perche' \(\partial C/\partial X\) e` quasi costante.

\slidepair{images9/part1-09.png}{images9/part1-10.png}{Dettaglio del comb drive e modello dinamico del giroscopio vibrante.}

La slide del modello dinamico descrive una singola massa sospesa in modo da potersi muovere sia lungo \(x\) sia lungo \(y\). Lungo \(x\) si impone la vibrazione di drive:
\[
m\ddot x+c\dot x+kx=F\sin(\omega_d t).
\]
Lungo \(y\) compare invece la forza di Coriolis:
\[
m\ddot y+c\dot y+ky=-2m\Omega_z\dot x.
\]

Risolvendo la prima equazione si ottiene l'ampiezza di oscillazione lungo \(x\):
\[
X=
\frac{F}{m}\frac{1}{\omega_x^2}
\frac{1}{\sqrt{(1-r_x^2)^2+(2\xi_xr_x)^2}},
\]
dove
\[
\omega_x=\sqrt{\frac{k_x}{m}},\qquad
\xi_x=\frac{c_x}{2m\omega_x},\qquad
r_x=\frac{\omega_d}{\omega_x}.
\]
Sostituendo \(\dot x\) nella seconda equazione, l'ampiezza del moto di sense e`
\[
Y=
\Omega_z\frac{2m\omega_dX}{k_y}
\frac{1}{\sqrt{(1-r_y^2)^2+(2\xi_yr_y)^2}},
\]
con
\[
\omega_y=\sqrt{\frac{k_y}{m}},\qquad
\xi_y=\frac{c_y}{2m\omega_y},\qquad
r_y=\frac{\omega_d}{\omega_y}.
\]

Questa formula mostra tre leve progettuali: aumentare l'ampiezza di drive \(X\), scegliere bene la frequenza di drive \(\omega_d\), e sfruttare la risonanza del modo di sense. Avvicinarsi alla risonanza aumenta la sensibilita`, ma rende anche il sensore piu` delicato rispetto a variazioni termiche, dispersioni di fabbricazione e instabilita` del fattore di qualita`.

\slidepair{images9/part1-11.png}{images9/part1-12.png}{Strutture MEMS reali e giroscopi vibranti 2D.}

Le strutture reali usano quasi sempre pettini capacitivi per il drive e per il pick-off, perche' sono compatibili con la microfabbricazione e permettono integrazione con l'elettronica. Nella pratica il sensore non e` solo una massa con due molle: e` una struttura elastica progettata per avere modi propri ben separati, buona simmetria e bassa sensibilita` alle accelerazioni lineari.

I giroscopi vibranti 2D estendono il principio a piu` assi. Un esempio e` il \emph{vibrating wheel gyroscope}: una ruota o anello vibrante viene eccitato in un modo piano, e rotazioni attorno a due assi ortogonali nel piano generano inclinazioni o deformazioni leggibili con sensori capacitivi posti sopra e sotto la struttura.

\slidefig{images9/part1-13.png}{Esempi di giroscopi vibranti 2D e strutture a ruota/anello.}

Il punto didattico e` che un MEMS non ``vede'' magicamente tre assi: ogni asse richiede un modo meccanico, una simmetria e una catena di lettura. La miniaturizzazione rende il sensore economico e compatto, ma non elimina i problemi classici della sensoristica: bias, rumore, cross-axis sensitivity, sensibilita` termica e calibrazione.

\subsection{Confronto tra tecnologie di giroscopi}

\slidefig{images9/part2-01.png}{Classi di accuratezza dei giroscopi: consumer, industrial, tactical, navigation e strategic.}

Il parametro usato piu` spesso per confrontare i giroscopi e` la stabilita` del bias, spesso indicata come \emph{bias stability} o \emph{in-run bias stability}. Il bias e` l'uscita del sensore quando la velocita` angolare vera e` zero; se questo valore cambia durante l'uso, l'integrazione nel tempo produce errore di assetto o di posizione.

Le slide riportano una classificazione indicativa:
\[
\begin{array}{ll}
\text{consumer/automotive} & >20\ \text{deg/h},\\
\text{industrial} & 5\text{--}20\ \text{deg/h},\\
\text{tactical} & 0.5\text{--}5\ \text{deg/h},\\
\text{navigation} & 0.01\text{--}0.1\ \text{deg/h},\\
\text{strategic} & 0.0001\text{--}0.01\ \text{deg/h}.
\end{array}
\]
I confini non sono universali: produttori diversi possono usare nomi leggermente diversi. Pero` l'ordine di grandezza e` utile, perche' collega immediatamente una tecnologia al tipo di navigazione che puo` sostenere.

\slidetriple{images9/part2-02.png}{images9/part2-03.png}{images9/part2-04.png}{Confronti grafici tra tecnologie: bias, rumore, costo, dimensioni e campo applicativo.}

I grafici di confronto vanno letti come mappe di compromesso. I MEMS dominano quando contano costo, dimensioni, consumo e integrazione; FOG, RLG e HRG dominano quando servono stabilita` e rumore molto basso; i giroscopi meccanici classici restano importanti storicamente e in alcune architetture ad alte prestazioni, ma pagano parti mobili, complessita` e manutenzione.

\slidefig{images9/part2-05.png}{Tabella di confronto tra giroscopi meccanici, MEMS, FOG, RLG e HRG.}

La tabella della slide puo` essere riassunta cosi`:

{\renewcommand{\arraystretch}{1.15}
\begin{longtable}{@{}>{\raggedright\arraybackslash}p{0.19\linewidth}>{\raggedright\arraybackslash}p{0.21\linewidth}>{\raggedright\arraybackslash}p{0.25\linewidth}>{\raggedright\arraybackslash}p{0.25\linewidth}@{}}
\toprule
\textbf{Tecnologia} & \textbf{Principio} & \textbf{Punti forti} & \textbf{Limiti} \\
\midrule
\endfirsthead
\toprule
\textbf{Tecnologia} & \textbf{Principio} & \textbf{Punti forti} & \textbf{Limiti} \\
\midrule
\endhead
Meccanico classico & Momento angolare e precessione di un rotore & Fisica matura, buona stabilita` a breve termine, lunga tradizione in stabilizzazione e navigazione & Parti mobili, usura, warm-up, sensibilita` agli urti, ingombro e costo elevati \\
DTG / floated gyro & Rotore preciso con supporto accordato o flottante & Bias migliore dei MEMS comuni, prestazioni tattiche o di navigazione & Costoso, meccanicamente complesso, soggetto ad aging e manutenzione \\
MEMS vibrante & Forza di Coriolis su microstruttura vibrante & Piccolo, economico, basso consumo, integrabile, ottimo SWaP & Drift del bias, sensibilita` a temperatura e vibrazioni, prestazioni di lungo periodo inferiori agli ottici \\
FOG & Fase Sagnac in bobina di fibra ottica & Nessuna parte mobile, alta affidabilita`, basso rumore, buona stabilita` del bias & Costo e dimensioni maggiori dei MEMS, elettronica/ottica piu` complessa \\
RLG & Differenza di frequenza Sagnac in cavita` laser & Prestazioni navigation/strategic, rumore molto basso, uscita a frequenza & Costoso, grande, complesso, richiede mitigazione del lock-in \\
HRG / resonator gyro & Effetto Coriolis in guscio risonante tipo wineglass & Lunga vita, assenza di rotore in spin, ottima robustezza e stabilita` & Produzione specializzata, costo elevato, meno diffuso dei MEMS \\
\bottomrule
\end{longtable}
}

La scelta non e` quindi ``qual e` il migliore in assoluto'', ma ``qual e` il migliore per il vincolo del sistema''. Un drone economico accetta un MEMS con correzione software; una girobussola o una navigazione inerziale di qualita` richiedono FOG, RLG o HRG; un sistema legacy puo` ancora usare meccanica di precisione.

\begin{keybox}
\textbf{Criterio pratico.} Se il sensore deve solo stabilizzare rapidamente un assetto, banda e rumore a breve termine possono pesare piu` della deriva a lungo termine. Se invece la velocita` angolare viene integrata per fare navigazione, la stabilita` del bias diventa dominante.
\end{keybox}

\subsection{Velocita` lineare propria: veicoli terrestri}

\slidepair{images9/part2-06.png}{images9/part2-07.png}{Velocita` lineare propria e stima tramite velocita` angolare delle ruote.}

Misurare la velocita` lineare propria significa stimare la velocita` del veicolo usando sensori montati sul veicolo stesso. Nel caso di un veicolo con ruote, la misura piu` semplice deriva dalla velocita` angolare della ruota:
\[
v=R\omega,
\]
dove \(R\) e` il raggio effettivo della ruota e \(\omega\) e` la velocita` angolare, tipicamente ricavata da encoder incrementali.

Con piu` ruote si puo` mediare:
\[
v=\frac{1}{N}\sum_i v_i,
\qquad
v_i=R_i\omega_i.
\]
Questa misura e` economica, veloce, ad alta banda e funziona indoor senza GPS. Per questo e` molto usata nel controllo di robot mobili, AGV e veicoli terrestri.

Il limite fondamentale e` lo \emph{wheel slip}. In accelerazione, frenata, curva, terreno irregolare o con pneumatici deformati, la velocita` periferica della ruota non coincide piu` con la velocita` del veicolo:
\[
R\omega \neq v_{\text{veicolo}}.
\]
L'encoder continua a misurare benissimo la ruota, ma la ruota non rappresenta piu` fedelmente il moto del corpo. Questo e` un punto importante: il problema non e` il sensore, e` il modello fisico che collega ruota e veicolo.

\subsection{Velocita` lineare in aeronautica: aria, Bernoulli e Pitot}

\slidepair{images9/part2-08.png}{images9/part2-09.png}{Airspeed e principio di Bernoulli: per il controllo dell'aereo conta la velocita` rispetto all'aria.}

Per un aereo la velocita` piu` importante per il controllo non e` sempre la velocita` rispetto al suolo, ma la velocita` rispetto alla massa d'aria. Portanza, resistenza e momenti aerodinamici dipendono dall'aria che investe le superfici, quindi un controllore di volo ha bisogno dell'\emph{airspeed}. La velocita` al suolo, invece, dipende anche dal vento ed e` piu` adatta alla navigazione geografica.

La tecnologia classica per misurare l'airspeed e` il tubo di Pitot. Il principio fisico e` Bernoulli, valido nel caso ideale di flusso incomprimibile, inviscido e irrotazionale:
\[
p+\frac{1}{2}\rho v^2+\rho gh=\text{costante}.
\]
Se si confrontano due punti circa alla stessa quota, il termine gravitazionale si cancella e rimane lo scambio tra pressione statica e pressione dinamica:
\[
p+\frac{1}{2}\rho v^2=\text{costante}.
\]
Quando un flusso viene rallentato fino a velocita` nulla, l'energia cinetica per unita` di volume si trasforma in aumento di pressione.

\slidepair{images9/part2-10.png}{images9/part2-11.png}{Tubo di Pitot: pressione statica, pressione totale e calcolo della velocita`.}

Il tubo di Pitot misura due pressioni:
\[
p_0=\text{pressione statica},\qquad
p_T=\text{pressione totale o di ristagno}.
\]
La pressione statica e` quella ambiente, misurata da fori laterali non investiti direttamente dal flusso. La pressione totale viene misurata nel foro frontale, dove il fluido viene idealmente arrestato.

Applicando Bernoulli alla porta di ristagno:
\[
p_0+\frac{\rho v_0^2}{2}=p_T.
\]
Da cui:
\[
v_0=\sqrt{\frac{2(p_T-p_0)}{\rho}}.
\]
La differenza
\[
q=p_T-p_0
\]
e` la pressione dinamica. In prima approssimazione, quindi, il Pitot non misura direttamente una velocita`: misura una differenza di pressione e la converte in velocita` tramite densita` dell'aria e modello fluidodinamico.

\slidefig{images9/part2-12.png}{Esempi di sonde e prese di pressione su un aereo; caso Air France 447 richiamato come esempio di criticita` della misura di airspeed.}

L'ultima slide mostra che su un aereo reale non c'e` solo un tubicino isolato: ci sono sonde Pitot, porte statiche, sonde di temperatura, sensori di angolo d'attacco, vane di sideslip e ridondanze. Questo perche' la misura di airspeed e` critica. Se il Pitot si ostruisce, ghiaccia o fornisce dati incoerenti, il sistema di controllo e il pilota possono perdere una delle grandezze piu` importanti per capire lo stato aerodinamico dell'aereo.

Il richiamo al volo Air France 447 serve proprio come avvertimento ingegneristico: una misura apparentemente semplice, se usata dentro un sistema complesso, deve essere ridondata, diagnosticata e confrontata con altre informazioni. In robotica e controllo questa idea torna continuamente: non basta avere un sensore accurato in condizioni nominali, bisogna sapere riconoscere quando non e` piu` affidabile.

\begin{keybox}
\textbf{Da ricordare.} L'odometria su ruote e il Pitot misurano velocita` proprie, ma rispetto a riferimenti fisici diversi. Le ruote stimano il moto rispetto al terreno tramite contatto; il Pitot stima il moto rispetto all'aria tramite pressione dinamica. Entrambe le misure diventano fragili quando il legame fisico si rompe: slittamento per le ruote, vento/ghiaccio/ostruzione o modello aerodinamico non valido per il Pitot.
\end{keybox}

\subsection*{Sintesi finale della Lezione 10}

La decima lezione chiude il panorama dei giroscopi mostrando come la misura di velocita` angolare possa essere realizzata anche senza rotori macroscopici e senza ottica: basta una microstruttura vibrante, un modo di drive noto e una lettura del moto laterale prodotto dalla forza di Coriolis. Il giroscopio MEMS e` quindi una soluzione elegantissima, ma non gratuita: miniaturizzazione e basso costo vengono pagati con bias drift, sensibilita` termica, rumore e calibrazione accurata.

Il confronto tra tecnologie chiarisce che la scelta del giroscopio e` un problema di sistema. MEMS, FOG, RLG, HRG e giroscopi meccanici non competono tutti nello stesso spazio: cambiano costo, prestazioni, dimensioni e affidabilita`.

Infine, la lezione apre la misura della velocita` lineare propria. Nei veicoli terrestri si parte da
\[
v=R\omega,
\]
ma il risultato vale solo se la ruota rotola senza slittare. Negli aerei si usa invece Bernoulli:
\[
p_T=p_0+\frac{1}{2}\rho v^2,
\qquad
v=\sqrt{\frac{2(p_T-p_0)}{\rho}}.
\]
In entrambi i casi il sensore misura una grandezza fisica locale e la trasforma in velocita` tramite un modello. Quando il modello non vale piu`, anche una misura tecnicamente corretta puo` diventare una stima sbagliata dello stato del veicolo.

% =========================================================
% LEZIONE 11
% =========================================================

\clearpage
\section{Lezione 11}

\subsection{Visione generale della lezione}

Questa lezione prosegue il discorso sulla velocita` lineare propria introducendo la misura basata su effetto Doppler. L'idea e` elegante: invece di dedurre la velocita` dal moto di una ruota o da una differenza di pressione, si osserva lo spostamento in frequenza di un'onda riflessa da un bersaglio in moto. Questo principio vale per onde acustiche, onde elettromagnetiche, laser e sonar.

La seconda parte cambia famiglia di sensori e introduce gli \emph{accelerometri}, cioe` i sensori inerziali che misurano l'accelerazione tramite il moto relativo di una massa di prova dentro un case. Qui il passaggio importante non e` solo meccanico ma concettuale: l'uscita dell'accelerometro dipende dalla forza specifica, quindi risente anche della gravita` quando il sensore e` inclinato.

Infine, la lezione rassegna alcune delle tecnologie principali: accelerometri a pendolo con force rebalance, accelerometri al quarzo, MEMS capacitivi e accelerometri termici. Si chiude con un confronto di accuratezza, che mette in evidenza il compromesso tra prestazioni, costo, robustezza e integrazione.

\subsection{Velocita` da effetto Doppler}

\slidepair{images10/page-01.png}{images10/page-02.png}{Principio Doppler e principali implementazioni: laser Doppler e radar/microwave Doppler.}

Se un'onda viene riflessa da un bersaglio in moto, la frequenza del segnale ricevuto cambia rispetto a quella emessa. La slide riporta, per il caso di misura monostatica con riflessione, la relazione:
\[
f_d=\frac{2v\cos\theta}{\lambda}
=\frac{2v f_{\text{emit}}\cos\theta}{c},
\]
dove:
\[
\begin{aligned}
f_d&=\text{Doppler shift}, & v&=\text{velocita` relativa},\\
\lambda&=\text{lunghezza d'onda}, & \theta&=\text{angolo tra fascio e velocita`}.
\end{aligned}
\]

Il fattore \(2\) compare perche' c'e` un'andata e ritorno del segnale riflesso. Il punto fisico piu` importante e` l'ultimo della slide: il sensore non misura tutta la velocita`, ma solo la componente lungo il fascio:
\[
v_{\parallel}=v\cos\theta.
\]
Se il fascio e` ortogonale al moto, lo shift si annulla e il sensore ``non vede'' quella componente.

Questa formula si puo` anche invertire:
\[
v=\frac{f_d\lambda}{2\cos\theta}.
\]
In pratica, quindi, la misura di velocita` e` una misura di frequenza. Questo e` molto utile perche' le frequenze si misurano bene, con alta risoluzione e spesso con elettronica relativamente robusta.

La slide successiva mostra due implementazioni tipiche. Nel \emph{Laser Doppler Velocimetry} (LDV) si usa un laser coerente e si misura l'interferenza della luce riflessa. Il vantaggio e` una risoluzione altissima, adatta a vibrazioni superficiali, misure di flusso e processi molto veloci. Nel \emph{radar Doppler}, invece, si lavora a radiofrequenza: la precisione metrica e` in genere inferiore a quella di un LDV, ma si ottengono grande portata e robustezza a polvere, pioggia e nebbia.

\slidefig{images10/page-03.png}{Doppler Velocity Logger per veicoli subacquei: piu` fasci acustici per ricostruire la velocita` rispetto al fondale.}

Per i veicoli subacquei il principio Doppler viene implementato con onde acustiche. Il \emph{Doppler Velocity Log} (DVL) invia piu` fasci verso il fondale e misura lo shift su ciascun fascio. Se \(\mathbf n_i\) e` il versore del fascio \(i\)-esimo, si puo` scrivere in forma compatta:
\[
f_{d,i}\propto \mathbf n_i^T\mathbf v.
\]
Con quattro fasci si ottiene un sistema ridondante che permette di stimare il vettore velocita` \(\mathbf v\) rispetto al fondale, tipicamente tramite minimi quadrati.

Il DVL funziona bene finche' il fondale e` ``visibile'' e considerabile fermo. Se la riflessione peggiora o il riferimento fisico manca, la stima si indebolisce. La logica, pero`, e` la stessa gia` vista per ruote e Pitot: il sensore ricava la velocita` attraverso un modello di interazione con un riferimento esterno, che qui e` il fondale marino.

\begin{keybox}
\textbf{Idea comune.} Ruote, Pitot e Doppler sono tre modi diversi di misurare una velocita` propria: contatto con il terreno, pressione dinamica rispetto all'aria, shift di frequenza rispetto a un bersaglio. In tutti e tre i casi il sensore vede solo la componente fisicamente accoppiata al principio di misura.
\end{keybox}

\subsection{Accelerometri: principio fisico}

\slidepair{images10/page-04.png}{images10/page-05.png}{Ingresso nel tema degli accelerometri e schema meccanico dell'elemento sensibile inerziale.}

Con questa slide il corso passa ai sensori inerziali di accelerazione. L'idea base e` semplice: se il case del sensore accelera, una massa interna tende a rimanere indietro per inerzia. Misurando lo spostamento relativo tra case e massa si ricava l'accelerazione applicata lungo l'asse sensibile.

\slidefig{images10/page-06.png}{Principio di un accelerometro massa-molla-smorzatore con pick-off di spostamento relativo.}

La slide formalizza il modello piu` classico: una massa di prova \(m\), una molla, uno smorzatore e un sensore di spostamento relativo. Se \(x\) e` la posizione del case, \(y\) la posizione assoluta della proof mass e
\[
\xi \triangleq y-x
\]
lo spostamento relativo, allora l'accelerazione del case e`
\[
a=\ddot x.
\]

Quando il case accelera, la massa tende a opporsi al moto. In regime statico o quasi statico, la forza elastica bilancia la forza inerziale:
\[
k\xi \simeq -ma,
\qquad
\xi \simeq -\frac{m}{k}a.
\]
Questa relazione spiega perche' l'uscita sia proporzionale all'accelerazione: il pick-off non misura direttamente \(a\), ma una deformazione \(\xi\) causata da \(a\).

\slidepair{images10/page-07.png}{images10/page-08.png}{Equazioni del moto dell'accelerometro e contributo della gravita` in caso di inclinazione.}

Applicando la seconda legge di Newton alla massa mobile si ottiene:
\[
m\ddot y = -k\xi - b\dot\xi + F_{\text{gravity}}.
\]
Poiche' \(y=x+\xi\), si ha
\[
m(\ddot\xi+\ddot x)=-k\xi-b\dot\xi+F_{\text{gravity}}.
\]
Nel caso ideale in cui l'asse sensibile sia ortogonale alla gravita` e quindi \(F_{\text{gravity}}=0\), la dinamica si riduce a
\[
m\ddot\xi+b\dot\xi+k\xi=-m\ddot x.
\]

Se invece il sensore e` inclinato di un angolo \(\gamma\), compare una componente gravitazionale lungo l'asse:
\[
F_{\text{gravity}}=-mg\sin\gamma,
\]
e quindi
\[
m\ddot\xi+b\dot\xi+k\xi=-m\ddot x-mg\sin\gamma.
\]
La slide lo interpreta come una accelerazione disturbante equivalente:
\[
a^{\text{dist}}=-g\sin\gamma.
\]

Questa osservazione e` importantissima. Un accelerometro ideale non misura ``solo'' l'accelerazione cinematica del corpo: misura la \emph{forza specifica}, cioe` il risultato delle forze non gravitazionali piu` l'effetto della proiezione della gravita` quando si cambia orientamento. Per questo gli accelerometri sono utili anche come inclinometri a bassa dinamica.

\slidefig{images10/page-09.png}{Dinamica dell'accelerometro: funzione di trasferimento, guadagno statico e banda naturale.}

Passando in Laplace, la slide ricava:
\[
G(s)=\frac{\xi(s)}{\ddot x(s)}
=-\frac{m}{k}\frac{k}{ms^2+bs+k}
=-\frac{m}{ms^2+bs+k}.
\]
Il guadagno statico vale
\[
G_0=-\frac{m}{k},
\]
mentre la pulsazione naturale e`
\[
\omega_n=\sqrt{\frac{k}{m}}.
\]
Se si introduce anche il rapporto di smorzamento
\[
\zeta=\frac{b}{2\sqrt{km}},
\]
si capisce bene il compromesso progettuale:
\begin{itemize}
\item aumentare \(m/k\) aumenta la sensibilita` statica;
\item aumentare \(k/m\) aumenta la banda e in generale la resistenza alla saturazione;
\item lo smorzamento \(b\) governa overshoot, risonanza e rapidita` del transitorio.
\end{itemize}

La slide lo dice in modo molto concreto: scegliendo \(m\), \(b\) e \(k\) si fissano banda, transitorio, range massimo e sensibilita`. In un accelerometro, queste grandezze non si possono ottimizzare tutte insieme senza compromessi.

\begin{keybox}
\textbf{Idea dinamica.} Un accelerometro non e` un semplice trasduttore statico forza-spostamento. E` un sistema dinamico del secondo ordine, e la sua utilita` dipende da dove cade la banda utile rispetto alla frequenza naturale e da quanto bene e` smorzato.
\end{keybox}

\subsection{Tecnologie di accelerometri}

\slidefig{images10/page-10.png}{Accelerometri a pendolo vincolato: la massa di prova e` un pendolo mantenuto vicino alla posizione di equilibrio.}

Gli \emph{restrained pendulum accelerometers} sono accelerometri di precisione in cui la proof mass e` un pendolo. Un pick-off misura la sua piccola rotazione o traslazione relativa, mentre un attuatore elettromagnetico applica la forza necessaria a riportarlo verso la posizione nominale.

\slidepair{images10/page-11.png}{images10/page-12.png}{Pendulum accelerometer a retroazione: il segnale di force-rebalance diventa la misura utile.}

Questa architettura e` molto vicina all'idea di \emph{force rebalance} gia` vista per i giroscopi. La massa non viene lasciata muovere liberamente: l'elettronica in anello chiuso riduce il suo spostamento e usa la corrente di feedback come misura dell'accelerazione.

Il vantaggio e` notevole:
\begin{itemize}
\item la massa si muove poco, quindi il comportamento resta piu` lineare;
\item si riducono errori di cross-axis e problemi geometrici;
\item si puo` estendere il range di misura senza richiedere grandi spostamenti meccanici.
\end{itemize}
Il prezzo da pagare e` una maggiore complessita` elettronica e costruttiva.

\slidefig{images10/page-13.png}{Esempio di accelerometro a pendolo biassiale con pick-off ottico.}

La variante mostrata in slide usa una massa collegata a una fibra ottica e una lettura ottica biassiale. La deformazione in due direzioni viene letta su un array fotosensibile, ottenendo una misura di accelerazione su due assi. Il principio fisico resta identico: una massa di prova, una cedevolezza meccanica e una lettura della deflessione.

\slidepair{images10/page-14.png}{images10/page-15.png}{Accelerometri al quarzo e MEMS capacitivi: due strade molto diverse per trasformare l'accelerazione in segnale elettrico.}

Negli accelerometri al quarzo mostrati nella slide, due travi vibranti lavorano in modo differenziale. A carico nullo vibrano alla stessa frequenza di risonanza; quando si applica una accelerazione, una trave si tende e l'altra si comprime:
\[
\Delta f = f_1-f_2 \propto a.
\]
La misura non e` quindi uno spostamento statico, ma una differenza di frequenza tra due risonatori. Questo aiuta la stabilita` e il rigetto dei disturbi comuni.

Nel MEMS capacitivo, invece, una massa microfabbricata si sposta di una piccola quantita` e modifica la capacita` di elettrodi differenziali. Vale la relazione base
\[
C=\frac{\varepsilon A}{d},
\]
quindi un piccolo spostamento che cambia \(d\) o l'area sovrapposta produce una variazione:
\[
\Delta C = C_1-C_2 \propto \xi \propto a.
\]
Questa e` una delle architetture piu` diffuse in assoluto, perche' combina costo basso, integrazione elevata e possibilita` di avere tre assi nello stesso chip.

\slidefig{images10/page-16.png}{Esempio di layout microfabbricato di accelerometro MEMS capacitivo integrato con elettronica sullo stesso die.}

La slide fotografica serve a ricordare che dietro il modello massa-molla c'e` una struttura reale fatta di ancoraggi, sospensioni, elettrodi differenziali e circuiti di lettura. Nei MEMS moderni la parte meccanica e quella elettronica sono fortemente co-progettate.

\slidepair{images10/page-17.png}{images10/page-18.png}{Accelerometri MEMS capacitivi triassiali: geometrie diverse per assi in piano e fuori piano.}

Per misurare tre assi sullo stesso die si usano geometrie diverse. Le accelerazioni \(x\) e \(y\) vengono spesso misurate con movimenti nel piano del chip e strutture a dita interdigitate. L'asse \(z\), invece, richiede una sensibilita` fuori piano e quindi strutture capacitive differenti, spesso basate su piastre superiore/inferiore o su variazioni di distanza verticali.

Il messaggio della slide e` importante: \emph{same die} non significa \emph{same sensor replicated three times}. Gli assi hanno vincoli geometrici diversi, e quindi anche rumore, banda, sensibilita` e cross-axis coupling possono essere diversi da asse ad asse.

\slidepair{images10/page-19.png}{images10/page-20.png}{Accelerometri termici: la grandezza utile non e` lo spostamento di una massa, ma lo sbilanciamento di temperatura nel gas.}

Negli accelerometri termici non si sfrutta una massa solida sospesa. Si riscalda il gas dentro una microcavita` con un elemento resistivo centrale, e si misurano con termopile o sensori termici le differenze di temperatura tra lati opposti.

Quando non c'e` accelerazione, il profilo termico e` simmetrico. Quando il sensore accelera, il gas caldo e meno denso si sposta in una direzione e quello piu` freddo nell'altra, creando una differenza di temperatura:
\[
\Delta T \propto a.
\]
Questa famiglia di sensori ha un vantaggio molto interessante: assenza di parti meccaniche solide in movimento, quindi robustezza elevata. Il limite, sottolineato anche in slide, e` la banda ridotta, dell'ordine di poche decine di hertz.

\slidefig{images10/page-21.png}{Confronto qualitativo di accuratezza tra accelerometri meccanici e accelerometri quartz/MEMS.}

La slide finale riassume la metrica piu` importante per confrontare gli accelerometri: stabilita` del bias e dello scale factor. Come per i giroscopi, il problema non e` solo ``quanto e` sensibile'' il sensore, ma quanto restano stabili nel tempo e con le condizioni operative lo zero e il coefficiente di conversione.

Il grafico suggerisce che gli accelerometri meccanici di alta gamma possono raggiungere prestazioni molto elevate, adatte alla navigazione. Gli accelerometri al quarzo e molti MEMS coprono invece un range piu` ampio che va dall'automotive all'industrial e, in alcuni casi, fino al tactical. Questo spiega perche' nella pratica convivano ancora tecnologie cosi` diverse: non servono tutte allo stesso problema.

\begin{keybox}
\textbf{Confronto finale.} Gli accelerometri a pendolo con retroazione massimizzano precisione e linearita`, ma sono costosi e complessi. I MEMS capacitivi dominano quando contano costo, integrazione e compattezza. I termici sacrificano banda per guadagnare robustezza e assenza di parti mobili solide. Non esiste un accelerometro ``migliore'' in assoluto: esiste il migliore per il profilo dinamico e per il budget di errore del sistema.
\end{keybox}

\subsection*{Sintesi finale della Lezione 11}

L'undicesima lezione completa il quadro delle misure di velocita` propria introducendo il Doppler. La grandezza misurata non e` una distanza o una pressione, ma uno shift di frequenza:
\[
f_d=\frac{2v\cos\theta}{\lambda}.
\]
Questo rende possibile misurare velocita` in modo contactless con laser, radar e sonar, fino al caso del DVL per i veicoli subacquei.

La seconda parte introduce gli accelerometri come sensori inerziali basati sul moto relativo di una proof mass. Il modello fondamentale e`
\[
m\ddot\xi+b\dot\xi+k\xi=-m\ddot x,
\]
con la correzione gravitazionale
\[
m\ddot\xi+b\dot\xi+k\xi=-m\ddot x-mg\sin\gamma
\]
quando il sensore e` inclinato. Questa e` la ragione per cui gli accelerometri sono utili sia per misurare dinamica sia per stimare l'assetto a bassa frequenza.

Infine, la rassegna tecnologica chiarisce i compromessi: pendoli con force rebalance per alta precisione, quarzo per stabilita`, MEMS capacitivi per integrazione e costo, termici per robustezza ma con banda ridotta. La lezione, in sostanza, mostra che ``misurare accelerazione'' non significa avere un solo sensore standard, ma un'intera famiglia di soluzioni con fisiche e compromessi molto diversi.

% =========================================================
% LEZIONE 12
% =========================================================

\clearpage
\section{Lezione 12}

\subsection{Visione generale della lezione}

Con questa lezione il corso apre un nuovo blocco tematico: si passa dai sensori agli \emph{attuatori}. Fin qui il problema principale era capire come misurare bene posizione, velocita` e accelerazione; da ora in poi il punto di vista si ribalta, perche' ci interessa capire come il sistema di controllo riesca davvero a \emph{produrre} moto, forza e coppia nel mondo fisico.

La prima parte introduce il concetto di attuatore e insiste su un messaggio molto importante per il controllo: l'attuatore non e` un elemento ideale che esegue perfettamente il comando del controllore. Ha saturazioni, zone morte, isteresi, limiti di banda, ritardi, vincoli energetici e limiti termici. Tutto questo influenza direttamente prestazioni, stabilita` e accuratezza del sistema chiuso.

La seconda parte entra nel caso piu` classico della robotica e dell'automazione, il \emph{motore DC brushed}. Le slide mostrano le leggi fisiche fondamentali, l'architettura minima statore-rotore-commutatore e la derivazione della coppia di una singola spira immersa in un campo magnetico. Il risultato chiave e` che la coppia nasce dall'interazione tra corrente elettrica e campo magnetico, ma senza commutazione resterebbe oscillante; e` il commutatore che rende possibile la rotazione continua.

\subsection{Ingresso nel blocco sugli attuatori}

\slidepair{images11/page-01.png}{images11/page-02.png}{Apertura del blocco ``Actuators'' e panoramica degli argomenti: caratteristiche generali, motori elettrici, driver e altre famiglie di attuazione.}

Le prime due slide hanno soprattutto un valore di orientamento. Dopo il lungo blocco sui sensori, il corso chiarisce che l'attuatore e` l'elemento che converte energia, tipicamente elettrica, in moto meccanico utile. In robotica questo significa trasformare una decisione di controllo in una coppia su un giunto, in una forza su un asse lineare, in una pressione su un attuatore pneumatico oppure in una portata comandata in un sistema idraulico.

L'outline mostra anche un aspetto metodologico importante: parlare di attuatori non significa parlare solo del motore in senso stretto. Bisogna includere l'elettronica di pilotaggio, l'interfaccia di potenza, gli organi di trasmissione, gli eventuali riduttori e, in molti casi, il dispositivo composto finale che il controllore vede davvero.

\slidefig{images11/page-03.png}{Definizione di attuatore e suo ruolo come interfaccia tra controllo e mondo fisico.}

La terza slide formula in modo molto netto l'idea centrale: un attuatore e` il punto di contatto tra algoritmi di controllo e realta` fisica. Questo e` il motivo per cui, in un progetto ben fatto, non si puo` separare completamente la sintesi del controllore dalla fisica dell'attuazione.

Un attuatore non e` un generatore ideale di forza o di coppia. Se il controllore chiede un'azione incompatibile con i limiti del dispositivo, l'attuatore non puo` eseguirla. Da qui nasce la frase, concettualmente giustissima, che la slide mette in evidenza: il controllore puo` essere buono solo quanto glielo permette l'attuatore. In pratica, anche il miglior regolatore perde efficacia se l'organo di attuazione satura, e` lento o introduce forti non linearita`.

\slidefig{images11/page-04.png}{Caratteristiche statiche e dinamiche degli attuatori: due famiglie di limiti da tenere sempre distinte.}

La quarta slide divide i problemi in due categorie:
\begin{itemize}
\item \textbf{caratteristiche statiche}, cioe` limiti presenti anche in condizioni lentamente variabili, come saturazione, zona morta, isteresi e guadagno non lineare;
\item \textbf{caratteristiche dinamiche}, cioe` limiti che emergono nel tempo e in frequenza, come banda, slew rate, ritardi, tempo di salita e tempo di assestamento.
\end{itemize}

Questa distinzione e` molto utile anche per leggere i problemi sperimentali. Se un asse non riesce a raggiungere il riferimento massimo, spesso il problema e` statico, per esempio una saturazione. Se invece arriva ma sempre in ritardo, oppure non segue bene profili rapidi, il problema e` dinamico.

\subsection{Non linearita` statiche: saturazione, zona morta, isteresi}

\slidepair{images11/page-05.png}{images11/page-06.png}{Saturazione dell'attuatore: il comando richiesto cresce, ma l'uscita reale non puo` superare i limiti fisici del dispositivo.}

La saturazione e` il limite piu` intuitivo e, spesso, il piu` pericoloso per il controllo. Se indichiamo con \(u_{in}\) il comando richiesto dal controllore e con \(u_{out}\) l'azione effettivamente disponibile, possiamo scrivere in forma compatta:
\[
u_{out}=\operatorname{sat}(u_{in}).
\]
Finche' \(|u_{in}|\) resta nel range disponibile, l'attuatore segue il comando. Oltre quel range, l'uscita si blocca al valore massimo o minimo fisicamente realizzabile.

Nel caso di un motore, i limiti possono essere dovuti a tensione massima, corrente massima del driver o coppia massima disponibile. La conseguenza pratica e` che il sistema smette di comportarsi linearmente. La slide lo dice in modo molto efficace: quando l'attuatore e` saturo, l'uscita resta costante anche se il controllore continua ad aumentare il comando. Dal punto di vista del plant, questo equivale a perdere autorita` di controllo.

Una conseguenza classica e` il \emph{windup} dell'integratore: il regolatore continua ad accumulare errore come se il comando potesse ancora crescere, ma l'attuatore e` gia` al limite. Quando il sistema rientra in zona lineare, il controllore si trova ``carico'' e produce sovraelongazioni o transitori molto lenti da smaltire.

\slidepair{images11/page-07.png}{images11/page-08.png}{Zona morta e isteresi: due non linearita` statiche che degradano precisione e ripetibilita`.}

La \emph{dead zone} descrive un intervallo di ingresso in cui l'uscita resta nulla:
\[
u_{out}=0
\qquad \text{per} \qquad
|u_{in}|<u_D.
\]
Fisicamente puo` nascere da attrito statico, giochi meccanici, soglie del driver o piccoli backlash nella trasmissione. Per il controllo, il problema e` serio soprattutto vicino allo zero: il regolatore chiede una piccola correzione, ma l'attuatore non produce nessun effetto finche' il comando non supera la soglia.

Questo spiega due fenomeni ricorrenti:
\begin{itemize}
\item scarsa precisione di regolazione attorno al riferimento;
\item possibili piccoli cicli limite, dovuti al continuo entrare e uscire dalla zona morta.
\end{itemize}

L'isteresi e` diversa, ma altrettanto insidiosa: l'uscita non dipende solo dal valore attuale dell'ingresso, ma anche dalla sua storia precedente. A parita` di comando, si possono ottenere uscite diverse a seconda che il segnale stia crescendo o diminuendo. Questo riduce la ripetibilita` e rende piu` difficile costruire un modello semplice e affidabile del dispositivo.

\begin{keybox}
\textbf{Messaggio pratico.} Saturazione, zona morta e isteresi non sono dettagli secondari: sono spesso il vero motivo per cui un sistema reale si discosta molto da quello simulato con modelli lineari ideali.
\end{keybox}

\subsection{Vincoli dinamici degli attuatori}

\slidepair{images11/page-09.png}{images11/page-10.png}{Slew rate e banda: anche un attuatore ``forte'' puo` essere troppo lento per seguire comandi rapidi.}

Lo \emph{slew rate} e` il limite sulla rapidita` con cui l'uscita dell'attuatore puo` cambiare:
\[
\left|\frac{du}{dt}\right|\le \dot u_{\max}.
\]
Questo limite compare, per esempio, nei driver di potenza, nei circuiti di pilotaggio della corrente o nei servosistemi con vincoli di velocita` interna. Dal punto di vista meccanico, uno slew rate limitato significa che l'attuatore non puo` sviluppare istantaneamente grandi variazioni di forza o coppia; quindi la dinamica del carico risulta piu` lenta di quanto il controllore vorrebbe.

La \emph{bandwidth} introduce la stessa idea, ma nel dominio della frequenza. Se \(G(s)\) e` la funzione di trasferimento dell'attuatore, la banda \(\omega_B\) viene tipicamente definita al punto di \(-3\) dB:
\[
|G(j\omega_B)|=\frac{|G(0)|}{\sqrt{2}}.
\]
Un attuatore con banda alta riesce a seguire variazioni rapide del comando; un attuatore con banda bassa appare invece ``pigro'' e attenua o ritarda le componenti veloci.

Le cause fisiche possono essere sia elettriche sia meccaniche. Nei motori compaiono costanti di tempo elettriche degli avvolgimenti, inerzie, attriti e limiti del driver. Il punto importante e` che la banda dell'attuatore impone un tetto naturale alla banda del controllo in anello chiuso: non si puo` pretendere di chiudere un loop molto piu` veloce dell'organo che deve applicare l'azione correttiva.

\slidefig{images11/page-11.png}{Tempo di salita, assestamento, overshoot e ritardo puro: lettura nel tempo delle prestazioni dinamiche dell'attuatore.}

Per un ingresso a gradino, la risposta reale dell'attuatore viene riassunta da grandezze classiche:
\begin{itemize}
\item \emph{rise time}, che misura quanto rapidamente l'uscita parte e raggiunge la zona del valore desiderato;
\item \emph{settling time}, cioe` il tempo necessario a stabilizzarsi entro una fascia accettabile;
\item \emph{overshoot}, che quantifica l'eventuale superamento del target;
\item \emph{time delay}, cioe` il ritardo tra comando e inizio della risposta.
\end{itemize}

Il ritardo si puo` modellare, in prima approssimazione, come un fattore esponenziale:
\[
G_{\text{delay}}(s)=e^{-T_d s}.
\]
Questo termine e` semplice da scrivere ma critico da gestire, perche' toglie margine di fase e rende piu` delicato aumentare le prestazioni del controllore.

Le slide sottolineano anche un aspetto molto concreto: questi tempi non dipendono solo dall'attuatore, ma dal sistema attuatore + carico. Cambiare inerzia, attrito o rigidezza del meccanismo collegato al motore significa cambiare la risposta che il controllore osservera` davvero.

\subsection{Energia, rendimento e limiti termici}

\slidepair{images11/page-12.png}{images11/page-13.png}{Rendimento energetico e vincoli termici: la potenza non convertita in lavoro utile diventa calore da dissipare.}

Un attuatore non e` solo un convertitore di moto, ma anche un convertitore di potenza. La grandezza piu` immediata per misurarne la qualita` energetica e` il rendimento:
\[
\eta=\frac{P_{\text{mecc}}}{P_{\text{in}}}.
\]
Il numeratore rappresenta la potenza meccanica effettivamente disponibile sul carico; il denominatore e` la potenza elettrica assorbita. La differenza tra le due non scompare: si trasforma in perdite.

Le slide citano le tre famiglie di perdita piu` classiche:
\begin{itemize}
\item riscaldamento Joule nei conduttori;
\item attriti meccanici;
\item perdite magnetiche.
\end{itemize}

Per la robotica mobile questa parte e` fondamentale, perche' il rendimento non influenza solo la temperatura, ma anche autonomia e dimensionamento della batteria.

La slide successiva spinge il ragionamento al livello fisico corretto: la potenza persa diventa calore interno, e quindi compare un vincolo di temperatura. Un attuatore non puo` lavorare indefinitamente alla massima potenza continua, perche' supererebbe i limiti dei materiali, dei MOSFET, degli avvolgimenti o dell'isolamento elettrico.

Per questo nei datasheet si trovano quasi sempre limiti multipli:
\begin{itemize}
\item corrente massima;
\item temperatura massima;
\item eventuali curve di \emph{derating}, cioe` riduzione delle prestazioni continuative al crescere della temperatura o del tempo di utilizzo.
\end{itemize}

\begin{keybox}
\textbf{Idea energetica.} Se il controllore chiede molta prestazione, l'attuatore non risponde solo in termini di coppia o velocita`: risponde anche con assorbimento, perdite e aumento di temperatura. Per questo la progettazione del controllo e quella della potenza non possono essere separate.
\end{keybox}

\subsection{Motore DC brushed: principio fisico e coppia elementare}

\slidefig{images11/page-14.png}{Leggi fisiche di base del motore DC brushed: forza di Lorentz e legge di Faraday.}

La seconda meta` della lezione entra nel motore DC brushed partendo dalle due leggi fisiche essenziali.

La prima e` la forza di Lorentz su un conduttore percorso da corrente in un campo magnetico:
\[
\mathbf F=I\,\mathbf l\times \mathbf B.
\]
Questa relazione spiega come da una corrente elettrica si generi una forza meccanica laterale sul conduttore. Nel motore, la geometria e` organizzata in modo che queste forze producano una coppia attorno all'asse del rotore.

La seconda e` la legge di Faraday:
\[
e=-\frac{d\Phi}{dt}.
\]
Quando un conduttore si muove in un campo magnetico, oppure quando varia il flusso concatenato, nasce una tensione indotta. Nel motore questa tensione prende il nome di \emph{back-EMF} ed e` il motivo per cui la parte elettrica del motore non puo` essere separata dalla sua velocita` di rotazione.

Le due leggi dicono quindi le due meta` del problema:
\begin{itemize}
\item corrente nel campo magnetico \(\Rightarrow\) forza \(\Rightarrow\) coppia;
\item moto nel campo magnetico \(\Rightarrow\) tensione indotta.
\end{itemize}

\slidepair{images11/page-15.png}{images11/page-16.png}{Architettura minima del motore DC e derivazione della coppia di una singola spira immersa nel campo magnetico.}

La slide sull'architettura minima identifica i tre blocchi essenziali:
\begin{itemize}
\item \textbf{statore}, che genera un campo magnetico quasi costante;
\item \textbf{rotore} o \textbf{armatura}, che contiene gli avvolgimenti attraversati da corrente;
\item \textbf{commutatore con spazzole}, che inverte la connessione elettrica al momento giusto.
\end{itemize}

La derivazione della slide successiva considera una singola spira. Ogni lato attivo, lungo \(l\), immerso in un campo \(B\) e percorso da corrente \(i\), subisce una forza di modulo
\[
F=Bil.
\]
Le due forze sui lati opposti formano una coppia. Se \(r\) e` la distanza dall'asse del rotore e \(\theta\) l'angolo tra normale alla spira e campo, il braccio efficace vale \(r\sin\theta\), quindi:
\[
\tau=2rF\sin\theta.
\]
Sostituendo \(F=Bil\) e introducendo l'area della spira \(A=2rl\), si ottiene:
\[
\tau=BiA\sin\theta.
\]
Con \(N\) spire identiche il risultato diventa:
\[
\tau=NBiA\sin\theta.
\]

Questa formula e` importante non solo per il motore DC classico, ma per tutta l'elettromeccanica: la coppia nasce da un accoppiamento tra corrente, campo magnetico e geometria della macchina.

\slidepair{images11/page-17.png}{images11/page-18.png}{Ruolo del commutatore e andamento della coppia con una sola spira nel caso ideale di commutazione ogni mezzo giro.}

Se non esistesse il commutatore, la corrente resterebbe sempre nello stesso verso elettrico e la coppia cambierebbe segno dopo mezzo giro: il rotore tenderebbe ad accelerare in una meta` della rotazione e a frenare nell'altra. Il commutatore risolve esattamente questo problema invertendo il verso della corrente ogni \(180^\circ\) meccanici.

In questo modo la coppia netta di una singola spira diventa sempre positiva:
\[
\tau_{\text{net}}(\theta)=\tau_{\max}|\sin\theta|.
\]
La slide finale lo mostra con un grafico molto chiaro: la coppia non e` costante, ma pulsante. Ci sono massimi attorno a \(90^\circ\) e \(270^\circ\), mentre in prossimita` di \(0^\circ\), \(180^\circ\) e \(360^\circ\) la coppia torna a zero.

Questo e` il motivo per cui i motori reali non usano una sola spira. Si impiegano piu` avvolgimenti e una commutazione distribuita, in modo da ridurre il ripple di coppia e rendere il moto piu` regolare.

\begin{keybox}
\textbf{Idea chiave sul motore DC brushed.} Il motore non ``crea'' coppia in modo misterioso: converte corrente in forza tramite Lorentz, e la disposizione geometrica delle forze genera coppia. Il commutatore fa il passo decisivo: sincronizza elettricamente la corrente con la posizione del rotore, cosi` la coppia media resta motrice e la rotazione puo` continuare.
\end{keybox}

\subsection*{Sintesi finale della Lezione 12}

La dodicesima lezione apre il blocco sugli attuatori e mette subito a fuoco il punto centrale per il controllo: l'attuatore e` l'elemento che rende fisicamente possibile il comando, ma proprio per questo introduce limiti e non idealita` che il controllore non puo` ignorare.

Le non linearita` statiche piu` importanti sono saturazione, zona morta e isteresi. I principali vincoli dinamici sono invece slew rate, banda, tempi di salita/assestamento e ritardi. Dal punto di vista energetico entra in gioco il rendimento
\[
\eta=\frac{P_{\text{mecc}}}{P_{\text{in}}},
\]
mentre dal punto di vista fisico e affidabilistico il limite termico impedisce di lavorare indefinitamente alla massima potenza.

La parte finale introduce il motore DC brushed a partire dalle leggi di Lorentz e Faraday:
\[
\mathbf F=I\,\mathbf l\times \mathbf B,
\qquad
e=-\frac{d\Phi}{dt}.
\]
Per una spira si ottiene la coppia
\[
\tau=NBiA\sin\theta,
\]
che, con commutazione ideale, diventa
\[
\tau_{\text{net}}(\theta)=\tau_{\max}|\sin\theta|.
\]
In altre parole, la lezione mostra che il tema dell'attuazione non e` ``dare corrente a un motore'', ma capire come limiti fisici, energetici ed elettromagnetici entrino direttamente nel comportamento dinamico del sistema controllato.

% =========================================================
% LEZIONE 13
% =========================================================

\clearpage
\section{Lezione 13}

\subsection{Visione generale della lezione}

La tredicesima lezione completa il modello del \emph{motore DC brushed} e prepara il passaggio ai motori passo-passo. La Lezione 12 aveva mostrato come una singola spira possa generare coppia, ma aveva anche lasciato aperto un problema evidente: con una sola spira, anche usando un commutatore ideale, la coppia resta pulsante e si annulla in alcune posizioni. La prima parte della lezione risolve proprio questo punto introducendo piu` avvolgimenti e mostrando perche' un motore reale puo` essere modellato, in prima approssimazione, con la relazione semplice
\[
\tau = k_t i.
\]

La seconda parte introduce la \emph{back-EMF}, cioe` la tensione indotta dalla rotazione del rotore. Questo e` il passaggio che trasforma il motore DC da semplice convertitore corrente-coppia a sistema elettromeccanico accoppiato: la corrente genera coppia, ma la rotazione genera una tensione che si oppone alla corrente. Da qui nascono il modello elettrico, il modello meccanico, la caratteristica coppia-velocita`, il punto di lavoro e la funzione di trasferimento.

La parte finale aggiunge due aspetti pratici importanti: i limiti termici del motore e l'introduzione agli \emph{stepper motor}. Gli stepper vengono presentati come attuatori discreti, capaci di posizionarsi per passi successivi tramite sequenze di eccitazione delle fasi.

\subsection{Ripple di coppia e ruolo di piu` avvolgimenti}

\slidepair{images12/part1-01.png}{images12/part1-02.png}{Con una sola spira e un commutatore ideale la coppia resta positiva, ma non e` costante e si annulla in posizioni critiche.}

La lezione riparte dal risultato finale della lezione precedente. Con una sola spira, il commutatore inverte la corrente ogni mezzo giro e rende la coppia sempre motrice:
\[
\tau_{\text{net}}(\theta)=\tau_{\max}|\sin\theta|.
\]
Questa forma e` gia` molto meglio della coppia alternata positiva-negativa che si avrebbe senza commutazione, perche' consente la rotazione continua. Pero` non e` ancora una coppia buona dal punto di vista del controllo.

Il grafico mostra infatti due difetti:
\begin{itemize}
\item la coppia e` fortemente ondulata, quindi la velocita` tende a oscillare anche con carico costante;
\item in corrispondenza di alcuni angoli la coppia e` nulla, quindi il motore ha posizioni in cui non riesce a ripartire bene.
\end{itemize}

Questo e` il motivo fisico per cui un motore reale non usa una sola spira. La singola spira serve per capire il principio, ma non e` una soluzione pratica se si vuole una coppia regolare.

\slidepair{images12/part1-03.png}{images12/part1-04.png}{Aggiungendo avvolgimenti sfasati, il commutatore seleziona quello che produce la coppia maggiore e il ripple diminuisce.}

Con due avvolgimenti sfasati di \(90^\circ\), il commutatore puo` selezionare l'avvolgimento che, in quel momento, fornisce la coppia maggiore:
\[
\tau(\theta)=\tau_0\max\left(|\sin\theta|,|\cos\theta|\right).
\]
La coppia minima non cade piu` a zero come nel caso a singolo avvolgimento. Il ripple resta presente, ma diventa piu` piccolo.

Con tre avvolgimenti sfasati di \(120^\circ\), l'idea e` la stessa:
\[
\tau(\theta)=\tau_0\max\left(
|\sin\theta|,
\left|\sin\left(\theta+\frac{2\pi}{3}\right)\right|,
\left|\sin\left(\theta+\frac{4\pi}{3}\right)\right|
\right).
\]
In generale, aumentando il numero di avvolgimenti e distribuendo meglio le commutazioni, la coppia risultante diventa piu` piatta.

\slidepair{images12/part1-05.png}{images12/part1-06.png}{Modello con \(N\) avvolgimenti e approssimazione del motore reale come generatore di coppia proporzionale alla corrente.}

Nel caso generale con \(N\) avvolgimenti, il commutatore ha due spazzole e \(2N\) lamelle. In forma ideale, la coppia disponibile puo` essere scritta come:
\[
\tau(\theta)=\tau_0\max_{k=0,\ldots,N-1}
\left|\sin\left(\theta+\frac{2\pi k}{N}\right)\right|.
\]
Questa formula va letta qualitativamente: il commutatore fa in modo che, tra gli avvolgimenti disponibili, resti attivo quello piu` favorevole alla produzione di coppia.

La slide aggiunge anche un dettaglio costruttivo importante. Nella pratica non si aumenta indefinitamente il numero di avvolgimenti principali; spesso si usano pochi avvolgimenti, tipicamente da tre a sei, ma si distribuiscono le spire e le lamelle in modo da rendere la commutazione piu` dolce. Alcuni conduttori collegati alla stessa coppia di lamelle possono essere avvolti anche in posizioni laterali, cosi` le loro coppie si sommano e la caratteristica totale si appiattisce.

Il risultato modellistico e` molto comodo:
\[
\tau(\theta)\simeq \tau_0=NBiA=k_t i.
\]
Per un campo magnetico fissato, la coppia sviluppata dal motore e` quindi proporzionale alla corrente di armatura. Tutta la complessita` geometrica del motore viene riassunta nella costante di coppia \(k_t\).

\begin{keybox}
\textbf{Idea chiave.} Piu` avvolgimenti non servono a cambiare il principio fisico della coppia, ma a renderlo usabile: trasformano una coppia pulsante in una coppia quasi continua, che puo` essere modellata come \(\tau=k_t i\).
\end{keybox}

\subsection{Back-EMF: la tensione generata dalla rotazione}

\slidepair{images12/part1-07.png}{images12/part1-08.png}{Legge di Faraday e flusso magnetico concatenato a una spira rotante.}

La \emph{back electromotive force}, o \emph{back-EMF}, nasce dalla legge di Faraday. Se il flusso magnetico concatenato con una spira cambia nel tempo, nella spira compare una tensione indotta:
\[
e=-\frac{d\Phi}{dt}.
\]
Per una bobina con \(N\) spire:
\[
e=-N\frac{d\Phi}{dt}.
\]

Il flusso magnetico e` definito come
\[
\Phi=\int \mathbf B\cdot d\mathbf A.
\]
Nel caso semplice di campo uniforme e spira piana di area \(A\), se \(\theta\) e` l'angolo tra il campo magnetico e la normale alla spira:
\[
\Phi=BA\cos\theta.
\]

\slidepair{images12/part1-09.png}{images12/part1-10.png}{Se la spira ruota con velocita` angolare \(\omega\), il flusso varia e genera una tensione proporzionale alla velocita`.}

Se il rotore ruota con velocita` angolare \(\omega\), si puo` scrivere
\[
\theta(t)=\omega t.
\]
Il flusso diventa:
\[
\Phi(t)=BA\cos(\omega t).
\]
Applicando Faraday:
\[
e(t)=-N\frac{d}{dt}\left(BA\cos(\omega t)\right)
    =NBA\omega\sin(\omega t).
\]
Per una singola bobina la back-EMF e` quindi sinusoidale. Nel motore DC brushed, pero`, il commutatore ``raddrizza'' internamente questa tensione alternata, in modo simile a quanto faceva con la coppia. Nel modello medio si ottiene:
\[
e\simeq k_e\omega.
\]

La back-EMF e` una tensione proporzionale alla velocita` e di segno opposto rispetto alla tensione che spinge la corrente. Per questo viene chiamata anche contro-elettromotrice: piu` il motore gira veloce, piu` genera una tensione che riduce la corrente disponibile a parita` di alimentazione.

\subsection{Modello elettromeccanico del motore DC}

\slidepair{images12/part1-11.png}{images12/part1-12.png}{Modello elettrico e modello meccanico del motore DC brushed.}

Dal lato elettrico, l'avvolgimento di armatura viene modellato come una resistenza \(R\), un'induttanza \(L\) e un generatore di back-EMF in serie. Applicando Kirchhoff:
\[
V=Ri+L\frac{di}{dt}+e.
\]
Sostituendo \(e=k_e\omega\):
\[
V=Ri+L\frac{di}{dt}+k_e\omega.
\]

Questa equazione mostra il primo accoppiamento: la parte meccanica rientra nel circuito elettrico tramite \(\omega\). A velocita` elevata, anche con la stessa tensione \(V\), resta meno margine per far crescere la corrente.

Dal lato meccanico si applica la seconda legge di Newton in rotazione. Indicando con \(J\) l'inerzia complessiva, con \(b\) l'attrito viscoso e con \(\tau_L\) la coppia resistente esterna:
\[
J\frac{d\omega}{dt}=\tau-b\omega-\tau_L.
\]
Poiche' \(\tau=k_t i\), si ottiene:
\[
J\frac{d\omega}{dt}=k_t i-b\omega-\tau_L.
\]

\slidepair{images12/part1-13.png}{images12/part1-14.png}{Sistema elettromeccanico completo e relazione tra costante di coppia \(k_t\) e costante di back-EMF \(k_e\).}

Il modello completo e` quindi:
\[
\begin{cases}
V=Ri+L\dfrac{di}{dt}+k_e\omega,\\[0.4em]
J\dfrac{d\omega}{dt}=k_t i-b\omega-\tau_L.
\end{cases}
\]
E` un sistema accoppiato: la corrente produce coppia, ma la velocita` produce back-EMF e quindi modifica la corrente.

Le slide chiariscono anche la relazione tra \(k_t\) e \(k_e\). In un motore ideale, senza perdite, la potenza elettrica associata alla back-EMF e` uguale alla potenza meccanica convertita:
\[
P_{\text{mecc}}=\tau\omega=k_t i\,\omega,
\qquad
P_{\text{elect}}=e i=k_e\omega\,i.
\]
Uguagliando le due espressioni:
\[
k_t=k_e.
\]
Questa uguaglianza vale quando si usano unita` SI coerenti. Dal punto di vista fisico, e` una manifestazione della reciprocita` della macchina elettrica: lo stesso accoppiamento che trasforma corrente in coppia trasforma velocita` in tensione.

\begin{infobox}
Nel modello di controllo conviene ricordare sempre entrambe le direzioni dell'accoppiamento: \(k_t\) dice quanta coppia si ottiene da una certa corrente, \(k_e\) dice quanta tensione opposta nasce da una certa velocita`.
\end{infobox}

\subsection{Comportamento a regime e caratteristica coppia-velocita`}

\slidepair{images12/part1-15.png}{images12/part1-16.png}{A regime, la relazione tra coppia disponibile e velocita` e` lineare.}

A regime stazionario, corrente e velocita` sono costanti:
\[
\frac{di}{dt}=0,
\qquad
\frac{d\omega}{dt}=0.
\]
Le equazioni diventano:
\[
V=Ri+k_e\omega,
\qquad
k_t i=b\omega+\tau_L.
\]

Trascurando temporaneamente l'attrito viscoso nella costruzione della caratteristica motore, dalla prima equazione si ricava:
\[
i=\frac{V-k_e\omega}{R}
  =\frac{V}{R}-\frac{k_e}{R}\omega.
\]
Poiche' \(\tau=k_t i\):
\[
\tau=k_t\frac{V}{R}-k_t\frac{k_e}{R}\omega.
\]
La coppia disponibile decresce linearmente con la velocita`. Questa e` una delle caratteristiche piu` importanti del motore DC: a tensione fissata, molta velocita` significa molta back-EMF, quindi meno corrente e meno coppia residua.

\slidepair{images12/part1-17.png}{images12/part1-18.png}{Rappresentazione grafica della caratteristica coppia-velocita`: coppia di stallo e corrente massima a velocita` nulla.}

La retta coppia-velocita` ha due intercette fondamentali. A velocita` nulla, la back-EMF e` nulla, quindi:
\[
i_{\text{stall}}=\frac{V}{R},
\qquad
\tau_{\text{stall}}=k_t\frac{V}{R}.
\]
Questa e` la \emph{stall torque}, cioe` la coppia massima che il motore puo` sviluppare a rotore bloccato. E` anche la condizione di corrente massima, quindi e` pericolosa termicamente: tutta la potenza elettrica assorbita diventa perdita Joule.

Quando invece la coppia richiesta si annulla, idealmente si ottiene la velocita` massima a vuoto:
\[
\omega_0=\frac{V}{k_e}.
\]
In pratica la velocita` a vuoto reale e` un po' piu` bassa, perche' servono comunque corrente e coppia per vincere attriti e perdite interne.

\slidepair{images12/part1-19.png}{images12/part1-20.png}{La massima potenza meccanica si ha a meta` della velocita` massima e a meta` della coppia di stallo.}

La potenza meccanica utile e`:
\[
P_{\text{out}}=\tau\omega.
\]
Poiche' \(\tau\) decresce linearmente con \(\omega\), il prodotto \(\tau\omega\) e` una parabola. Nel modello ideale la potenza massima si trova a meta` della retta:
\[
\omega=\frac{\omega_0}{2},
\qquad
\tau=\frac{\tau_{\text{stall}}}{2}.
\]

Questa informazione e` molto utile per leggere i datasheet. Il punto di massima potenza non coincide con il punto di massima efficienza, e non coincide nemmeno con il punto piu` sicuro termicamente. Un motore puo` erogare molta potenza per brevi intervalli, ma l'uso continuativo dipende da corrente, raffreddamento e temperatura.

La slide sull'efficienza ricorda la definizione:
\[
\eta=\frac{P_{\text{out}}}{P_{\text{in}}}
    =\frac{\tau\omega}{Vi}.
\]
Allo stallo, l'efficienza e` nulla perche' \(\omega=0\). A vuoto, anche la potenza meccanica utile tende a zero. La zona utile sta quindi in mezzo, ma la scelta del punto operativo dipende dal carico e dal ciclo di lavoro.

\slidepair{images12/part1-21.png}{images12/part1-22.png}{La caratteristica puo` essere disegnata con assi scambiati; cio` non cambia il contenuto fisico.}

Alcuni grafici mettono la coppia sull'asse orizzontale e la velocita` sull'asse verticale, altri fanno il contrario. La slide segnala questa convenzione per evitare confusione. Il contenuto fisico resta lo stesso: a tensione fissata, coppia e velocita` sono legate da una retta a pendenza negativa.

Una scrittura utile e`:
\[
\tau=\tau_{\text{stall}}\left(1-\frac{\omega}{\omega_0}\right).
\]
Invertendo:
\[
\omega=\omega_0\left(1-\frac{\tau}{\tau_{\text{stall}}}\right).
\]
L'attenzione va quindi posta sulle intercette e sul punto di lavoro, non sull'orientamento grafico degli assi.

\slidepair{images12/part1-23.png}{images12/part1-24.png}{Punto di lavoro e traslazione della caratteristica al variare della tensione di alimentazione.}

Il punto di lavoro si trova dove la coppia fornita dal motore uguaglia la coppia richiesta dal carico piu` gli attriti:
\[
k_t i=b\omega+\tau_L.
\]
Graficamente e` l'intersezione tra la caratteristica del motore e la caratteristica del carico.

Cambiare la tensione di alimentazione \(V\) trasla la retta coppia-velocita`. Infatti:
\[
\tau=k_t\frac{V}{R}-k_t\frac{k_e}{R}\omega.
\]
Il termine di intercetta \(k_tV/R\) cresce con \(V\), quindi aumentando la tensione aumentano sia la coppia di stallo sia la velocita` a vuoto ideale. Questa e` la base intuitiva del controllo in tensione di un motore DC: modificando \(V\), si sposta la caratteristica e quindi cambia il punto di equilibrio con il carico.

\begin{keybox}
\textbf{Lettura della caratteristica.} La retta coppia-velocita` nasce dalla back-EMF: il motore veloce genera piu` tensione opposta, assorbe meno corrente e quindi produce meno coppia.
\end{keybox}

\subsection{Aspetti termici del motore DC}

\slidepair{images12/part2-01.png}{images12/part2-02.png}{Perdite Joule, temperatura degli avvolgimenti e ruolo delle resistenze termiche.}

Il limite termico e` uno dei piu` importanti nella scelta di un motore. La principale sorgente di perdita e` spesso la perdita Joule negli avvolgimenti:
\[
P_J=Ri^2.
\]
A questa si aggiungono perdite nel ferro, dovute a correnti parassite e isteresi, e perdite meccaniche per attrito.

La temperatura cresce perche' la potenza dissipata diventa calore. Una rappresentazione semplice usa la resistenza termica, in analogia con una resistenza elettrica:
\[
\Delta T=P_J\left(R_{\text{th1}}+R_{\text{th2}}\right).
\]
La temperatura degli avvolgimenti a regime diventa:
\[
T_{\text{winding}}=T_{\text{amb}}+\Delta T.
\]
Un dissipatore o un miglior accoppiamento termico riducono la resistenza termica verso l'ambiente e quindi abbassano la temperatura a parita` di corrente.

\slidepair{images12/part2-03.png}{images12/part2-04.png}{Effetto della temperatura su resistenza degli avvolgimenti e costante di coppia dei magneti.}

La temperatura non e` solo una conseguenza del funzionamento: cambia anche i parametri del motore. La resistenza del rame cresce con la temperatura:
\[
R_{\text{warm}}=R\left(1+\alpha_{\text{cu}}\Delta T\right),
\qquad
\alpha_{\text{cu}}\simeq 0.0039\,K^{-1}.
\]
Quindi, a caldo, a parita` di tensione, il motore assorbe meno corrente. Allo stesso tempo aumenta la potenza dissipata per ogni ampere richiesto, perche' \(P_J=Ri^2\).

Anche il campo generato dai magneti permanenti varia con la temperatura. Per magneti al neodimio, la slide usa:
\[
k_{t,\text{warm}}=k_t\left(1+\alpha_{\text{mag}}\Delta T\right),
\qquad
\alpha_{\text{mag}}\simeq -0.0011\,K^{-1}.
\]
Il coefficiente e` negativo: scaldando il motore, i magneti perdono prestazione e la costante di coppia diminuisce. Se si superano certe temperature, compare anche il rischio di demagnetizzazione permanente.

Dal punto di vista del controllo, questo significa che un motore caldo non e` identico allo stesso motore a freddo. La stessa tensione non produce la stessa corrente, e la stessa corrente non produce esattamente la stessa coppia.

\subsection{Funzione di trasferimento del motore DC}

\slidepair{images12/part2-05.png}{images12/part2-06.png}{Dalle equazioni differenziali al diagramma a blocchi del motore DC.}

Per ricavare la funzione di trasferimento si parte dal modello senza coppia di carico esterna:
\[
V(t)=Ri(t)+L\frac{di(t)}{dt}+k_e\omega(t),
\]
\[
J\frac{d\omega(t)}{dt}+b\omega(t)=k_t i(t).
\]
Applicando la trasformata di Laplace:
\[
V(s)=(R+Ls)I(s)+k_e\Omega(s),
\]
\[
(Js+b)\Omega(s)=k_t I(s).
\]

Il diagramma a blocchi mostra bene l'accoppiamento: il ramo diretto converte tensione in corrente e corrente in coppia, mentre la velocita` torna indietro come feedback negativo di back-EMF.

\slidepair{images12/part2-07.png}{images12/part2-08.png}{Funzione di trasferimento completa tra tensione applicata e velocita` angolare.}

Eliminando \(I(s)\), si ottiene:
\[
\frac{\Omega(s)}{V(s)}
=
\frac{k_t}
{(Ls+R)(Js+b)+k_tk_e}.
\]
Sviluppando il denominatore:
\[
\frac{\Omega(s)}{V(s)}
=
\frac{k_t}
{LJ s^2+(Lb+RJ)s+(Rb+k_tk_e)}.
\]

Se l'attrito viscoso viene trascurato, cioe` \(b\simeq 0\), la funzione diventa:
\[
\frac{\Omega(s)}{V(s)}
=
\frac{k_t}
{LJ s^2+RJ s+k_tk_e}.
\]
Questo e` un secondo ordine, perche' sono presenti due accumuli di energia: l'induttanza elettrica \(L\) e l'inerzia meccanica \(J\).

\slidepair{images12/part2-09.png}{images12/part2-10.png}{Separazione tra costante di tempo elettrica e costante di tempo meccanica; modello approssimato del motore.}

Le radici del denominatore, nel caso \(b=0\), sono:
\[
s_{1,2}
=
\frac{-RJ\pm\sqrt{(RJ)^2-4LJk_tk_e}}{2LJ}.
\]
Se la dinamica elettrica e` molto piu` rapida di quella meccanica:
\[
\frac{L}{R}\ll \frac{RJ}{k_tk_e},
\]
allora i poli si approssimano come:
\[
s_{1,2}\simeq -\frac{R}{L},\quad -\frac{k_tk_e}{RJ}.
\]
Le costanti di tempo corrispondenti sono:
\[
\tau_e=\frac{L}{R},
\qquad
\tau_m=\frac{RJ}{k_tk_e}.
\]

Quando \(\tau_e\ll\tau_m\), la corrente raggiunge quasi subito il suo valore di equilibrio rispetto alla velocita`. Si puo` quindi porre:
\[
L\frac{di}{dt}\simeq 0,
\qquad
i(t)\simeq \frac{V(t)-k_e\omega(t)}{R}.
\]
Sostituendo nella dinamica meccanica:
\[
J\frac{d\omega}{dt}
+
\left(b+\frac{k_tk_e}{R}\right)\omega
=\frac{k_t}{R}V.
\]
La funzione di trasferimento approssimata e`:
\[
\frac{\Omega(s)}{V(s)}
=
\frac{k_t/R}
{Js+b+k_tk_e/R}
=
\frac{k_t}
{RJ s+Rb+k_tk_e}.
\]
La costante di tempo approssimata diventa:
\[
\tau_m^{\prime}
=
\frac{J}{b+k_tk_e/R}
=
\frac{RJ}{Rb+k_tk_e},
\]
che coincide con \(\tau_m\) se \(b=0\).

\begin{keybox}
\textbf{Modello utile per il controllo.} Se l'induttanza e` veloce rispetto alla meccanica, il motore DC puo` essere trattato come un primo ordine tra tensione e velocita`, con una retroazione interna dovuta alla back-EMF.
\end{keybox}

\slidefig{images12/part2-11.png}{Esempi di motori DC brushed e loro componenti: magneti permanenti, rotore, commutatore e spazzole.}

La slide fotografica serve a collegare il modello alla costruzione fisica. Nei piccoli motori DC si ritrovano gli stessi blocchi modellati nelle equazioni:
\begin{itemize}
\item i magneti permanenti dello statore generano il campo \(B\);
\item il rotore contiene le bobine percorse da corrente;
\item il commutatore e le spazzole realizzano la commutazione meccanica;
\item l'albero trasmette la coppia al carico.
\end{itemize}
Le spazzole e il commutatore sono anche una sorgente pratica di usura, rumore elettrico e limiti di vita utile. Per questo, in molte applicazioni moderne, i motori brushed vengono sostituiti da motori brushless; il principio elettromagnetico resta simile, ma la commutazione e` elettronica invece che meccanica.

\subsection{Introduzione agli stepper motor}

\slidetriple{images12/part2-12.png}{images12/part2-13.png}{images12/part2-14.png}{Principio di funzionamento generale degli stepper motor e costruzione fisica di statore e rotore.}

Gli \emph{stepper motor} sono motori progettati per muoversi per passi discreti. Il rotore si porta in posizioni di equilibrio successive allineandosi con il campo magnetico generato dallo statore oppure cercando di minimizzare la riluttanza del circuito magnetico.

Le slide distinguono tre famiglie principali:
\begin{itemize}
\item \textbf{Permanent Magnet} (PM), con rotore a magnete permanente;
\item \textbf{Variable Reluctance} (VR), in cui il rotore tende alla posizione di minima riluttanza;
\item \textbf{Hybrid} (HB), che combina magnete permanente e dentature per ottenere passi piccoli e buona coppia.
\end{itemize}

Le caratteristiche comuni sono:
\begin{itemize}
\item moto in passi;
\item controllo di posizione semplice in anello aperto, almeno se non si perdono passi;
\item coppia di mantenimento a rotore fermo;
\item velocita` determinata dalla frequenza degli impulsi;
\item verso di rotazione determinato dall'ordine di eccitazione delle fasi.
\end{itemize}

Dal punto di vista costruttivo, lo statore porta avvolgimenti di fase che generano un campo magnetico sequenziale. A ogni stato di eccitazione corrisponde un asse magnetico preferito; passando allo stato successivo, l'equilibrio si sposta e il rotore compie un passo.

\slidepair{images12/part2-15.png}{images12/part2-16.png}{Motore passo-passo a magnete permanente: attrazione/repulsione magnetica e avanzamento tra posizioni di equilibrio.}

Nel motore passo-passo a magnete permanente, lo statore ha poli salienti con avvolgimenti di fase, mentre il rotore e` un magnete cilindrico o a coppa con poli nord e sud alternati. Quando una fase viene alimentata, un polo dello statore diventa nord e quello opposto sud; il rotore subisce una coppia di attrazione/repulsione e si allinea con il campo.

Quando l'eccitazione passa alla fase successiva, il campo statorico ruota di un angolo fissato dalla geometria. Il rotore segue il nuovo equilibrio. Nel caso elementare mostrato, con due fasi statoriche e due poli rotorici, ci sono quattro posizioni stabili per giro e quindi:
\[
\alpha_s=90^\circ.
\]

Questa modalita` di funzionamento rende lo stepper molto intuitivo per il posizionamento: se ogni impulso produce un passo, contando gli impulsi si conosce la posizione comandata. Il limite e` che la posizione reale coincide con quella comandata solo se il motore non perde passi.

\slidepair{images12/part2-16.png}{images12/part2-17.png}{Regola geometrica del passo per un motore PM stepper e esempio con \(m=2\), \(N_r=6\).}

Per un PM stepper, la slide introduce:
\[
m=\text{numero di fasi statoriche},
\qquad
N_r=\text{numero di poli rotorici},
\qquad
p_r=\frac{N_r}{2}.
\]
La regola geometrica proposta per l'angolo di passo e`:
\[
\alpha_s=\frac{360^\circ}{mN_r}.
\]
Nell'esempio:
\[
m=2,
\qquad
N_r=6,
\qquad
\alpha_s=30^\circ.
\]

Il segno della corrente decide quale polo statorico diventa nord o sud e quindi anche il verso dello spostamento verso il polo rotorico piu` vicino. Questa e` la base per ottenere rotazione controllata: non si applica una tensione continua qualunque, ma una sequenza ordinata di eccitazioni di fase.

I motori stepper reali, soprattutto ibridi, usano dentature e molte piu` posizioni magnetiche, arrivando a passi tipici molto piu` piccoli, per esempio \(1.8^\circ\) nei motori da 200 passi/giro. Inoltre, con il \emph{microstepping}, le correnti nelle fasi vengono modulate in modo sinusoidale per ottenere posizioni intermedie e ridurre vibrazioni e ripple di coppia.

\begin{keybox}
\textbf{Differenza di filosofia.} Nel motore DC brushed il comando modifica una caratteristica continua coppia-velocita`; nello stepper il comando e` una sequenza di stati magnetici che sposta il punto di equilibrio del rotore passo dopo passo.
\end{keybox}

\subsection*{Sintesi finale della Lezione 13}

La lezione completa il quadro del motore DC brushed. La coppia pulsante della singola spira viene resa quasi continua tramite piu` avvolgimenti e commutazione distribuita, permettendo il modello medio:
\[
\tau=k_t i.
\]
La rotazione genera pero` una tensione opposta:
\[
e=k_e\omega,
\]
che accoppia parte elettrica e parte meccanica. Il modello completo e`:
\[
\begin{cases}
V=Ri+L\dfrac{di}{dt}+k_e\omega,\\[0.4em]
J\dfrac{d\omega}{dt}=k_t i-b\omega-\tau_L.
\end{cases}
\]

A regime, la back-EMF produce la caratteristica lineare coppia-velocita`:
\[
\tau=k_t\frac{V}{R}-k_t\frac{k_e}{R}\omega.
\]
Da qui si leggono coppia di stallo, velocita` a vuoto, punto di massima potenza e punto operativo con il carico.

La parte termica ricorda che la corrente non e` gratis: \(P_J=Ri^2\) scalda gli avvolgimenti, modifica \(R\), riduce le prestazioni dei magneti e puo` imporre limiti molto piu` severi di quelli dinamici. Infine, l'introduzione agli stepper mostra una famiglia diversa di attuatori elettrici: invece di produrre moto continuo tramite una coppia media, lo stepper realizza posizioni discrete tramite l'allineamento successivo del rotore con il campo statorico.

% =========================================================
% LEZIONE 14
% =========================================================

\clearpage
\section{Lezione 14}

\subsection{Visione generale della lezione}

La quattordicesima lezione completa il blocco sugli \emph{stepper motor} e apre quello sui \emph{brushless DC motor}. Il filo conduttore e` la commutazione: negli stepper la sequenza di eccitazione sposta a scatti la posizione di equilibrio del rotore; nei BLDC la commutazione elettronica sostituisce il commutatore meccanico del motore DC brushed e deve essere sincronizzata con la posizione del rotore.

La prima parte riprende gli stepper da dove si era fermata la Lezione 13. Dopo il motore a magnete permanente vengono introdotti i motori a riluttanza variabile e ibridi. La differenza fisica e` importante: nel PM stepper la coppia nasce soprattutto dall'interazione tra magnete permanente e campo statorico, nel VR stepper nasce dalla tendenza del circuito magnetico a ridurre la riluttanza, mentre nello hybrid stepper le due idee vengono combinate per ottenere passi piccoli e buona coppia di mantenimento.

La seconda parte entra nei motori BLDC. Dal punto di vista costruttivo il rotore contiene magneti permanenti e lo statore contiene tre fasi. Dal punto di vista del controllo, pero`, il punto cruciale e` che non esistono spazzole: le fasi devono essere alimentate da un inverter, scegliendo tensioni e correnti in funzione della posizione elettrica del rotore. La lezione costruisce quindi il modello trifase, introduce le back-EMF di fase, ricava la coppia elettromagnetica da un bilancio di potenza e chiude con la commutazione trapezoidale a sei settori tramite sensori Hall.

\subsection{Stepper a riluttanza variabile}

\slidetriple{images13/part1-01.png}{images13/part1-02.png}{images13/part1-03.png}{Motore passo-passo a riluttanza variabile: statore dentato, rotore ferromagnetico e sequenza di eccitazione delle fasi.}

Il motore passo-passo a riluttanza variabile, indicato spesso come \emph{VR stepper}, non ha magneti permanenti nel rotore. Il rotore e` realizzato in materiale ferromagnetico dolce e presenta denti salienti. Lo statore, anch'esso dentato, contiene avvolgimenti concentrati sulle fasi.

Il principio di funzionamento e` diverso da quello del PM stepper. Il rotore non possiede una polarita` magnetica propria; quindi non viene attratto perche' ha un nord o un sud gia` fissato. Quando una fase statorica viene alimentata, il campo magnetico prodotto dagli avvolgimenti cerca il percorso a riluttanza minima. Il rotore si muove allora nella posizione in cui i suoi denti si allineano con i denti statorici eccitati, perche' in quella configurazione il flusso attraversa piu` ferro e meno traferro.

La riluttanza magnetica gioca qui un ruolo analogo alla resistenza elettrica: un circuito magnetico con molto traferro ha riluttanza alta, uno con molto materiale ferromagnetico allineato ha riluttanza bassa. La coppia nasce quindi dalla tendenza del sistema a massimizzare l'induttanza vista dalla fase eccitata, o equivalentemente a minimizzare la riluttanza del percorso magnetico.

Quando si alimenta la fase successiva, la posizione di minima riluttanza si sposta. Il rotore segue questo nuovo minimo e compie un passo. In questo modo la rotazione non nasce da una tensione continua applicata al motore, ma da una sequenza discreta di stati magnetici. La direzione dipende dall'ordine della sequenza: per esempio, una sequenza \(A,B,C,A,\ldots\) puo` produrre rotazione in un verso, mentre \(A,C,B,A,\ldots\) produce il verso opposto.

\begin{keybox}
\textbf{Idea chiave del VR stepper.} Il rotore non deve essere magnetizzato: basta che sia ferromagnetico. La posizione stabile e` quella in cui il circuito magnetico ha minima riluttanza, quindi il controllo consiste nello spostare questo minimo da una fase alla successiva.
\end{keybox}

\subsection{Regole geometriche per il passo del VR stepper}

\slidepair{images13/part1-04.png}{images13/part1-05.png}{Regola del numero di denti per lo stepper a riluttanza variabile ed esempio con \(N_s=8\), \(N_r=6\).}

Le slide successive trasformano il principio fisico in una regola geometrica. Indichiamo con:
\[
N_s=\text{numero di denti statorici},
\qquad
N_r=\text{numero di denti rotorici}.
\]
Per ottenere moto passo-passo, \(N_s\) e \(N_r\) devono essere diversi. Se fossero uguali e perfettamente allineabili tutti insieme, l'eccitazione delle fasi non produrrebbe lo spostamento progressivo desiderato dei minimi di riluttanza.

L'angolo di passo fondamentale viene dato dalla relazione:
\[
\alpha_s
=360^\circ
\left|
\frac{N_s-N_r}{N_sN_r}
\right|.
\]
La formula dice una cosa molto intuitiva: se statore e rotore hanno molti denti e il loro numero e` vicino, lo sfasamento geometrico tra un allineamento e il successivo e` piccolo, quindi il passo diventa piccolo. In progettazione si scelgono dunque \(N_s\) e \(N_r\) abbastanza grandi e vicini tra loro, ma compatibili con il raggruppamento delle fasi e con la costruzione meccanica.

Nell'esempio delle slide:
\[
N_s=8,
\qquad
N_r=6.
\]
Allora:
\[
\alpha_s
=360^\circ\frac{|8-6|}{8\cdot 6}
=360^\circ\frac{2}{48}
=15^\circ.
\]
La figura mostra anche la sequenza delle fasi. In quel caso, per una rotazione antioraria si usa la sequenza \(A,D,C,B,A,\ldots\), mentre per la rotazione oraria si usa \(A,B,C,D,A,\ldots\). Ancora una volta, la posizione comandata non e` imposta direttamente da un sensore di posizione, ma dall'ipotesi che il rotore riesca a seguire la sequenza magnetica.

\subsection{Stepper ibrido: costruzione e motivo della sua diffusione}

\slidetriple{images13/part1-06.png}{images13/part1-07.png}{images13/part1-08.png}{Motore passo-passo ibrido: magnete permanente assiale, coppie di denti rotorici sfalsate e statore dentato.}

Il motore passo-passo ibrido combina le due famiglie precedenti. Da una parte ha un magnete permanente nel rotore, quindi conserva una componente di coppia legata all'attrazione/repulsione magnetica. Dall'altra usa dentature di rotore e statore, quindi sfrutta anche la tendenza alla minima riluttanza.

La struttura tipica e` composta da:
\begin{itemize}
\item un magnete permanente magnetizzato assialmente all'interno del rotore;
\item due coppe rotoriche dentate, una associata al polo nord e una al polo sud;
\item due serie di denti rotorici sfalsate di mezzo passo di dente;
\item uno statore dentato con avvolgimenti di fase.
\end{itemize}

Se il rotore ha \(N_r\) denti, il passo dentale rotorico e`:
\[
\alpha_R=\frac{360^\circ}{N_r}.
\]
Le fasi dello statore sono meccanicamente sfalsate di un quarto di questo passo. Questa e` la piccola finezza geometrica che rende possibile ottenere passi molto piccoli senza aumentare in modo esagerato il numero di fasi elettriche.

Il termine \emph{hybrid} non e` quindi solo commerciale: il motore produce coppia sia per effetto del magnete permanente sia per effetto di riluttanza. Il risultato pratico e` molto utile:
\begin{itemize}
\item angoli di passo piccoli, tipicamente \(1.8^\circ\) o \(0.9^\circ\);
\item coppia di mantenimento elevata;
\item buon compromesso tra precisione, costo e semplicita` di pilotaggio.
\end{itemize}

Per questo gli hybrid stepper sono molto comuni in macchine CNC leggere, stampanti 3D, piccoli assi automatici e sistemi dove e` accettabile lavorare in anello aperto pur mantenendo una buona risoluzione di posizione comandata.

\subsection{Principio di funzionamento dello stepper ibrido}

\slidepair{images13/part1-09.png}{images13/part1-10.png}{Funzionamento dello stepper ibrido: l'equilibrio si sposta di un quarto di passo dentale a ogni commutazione di fase.}

Quando una fase statorica viene eccitata, i denti del rotore e dello statore possono attrarsi o respingersi a seconda della polarita` della fase e del polo rotorico coinvolto. Il rotore si muove verso una posizione di equilibrio in cui:
\begin{itemize}
\item i denti statorici eccitati sono allineati con i denti rotorici compatibili;
\item la riluttanza tra poli eccitati e rotore e` minima;
\item il flusso magnetico attraversa il percorso piu` favorevole.
\end{itemize}

La slide riassume questo spostamento dicendo che l'equilibrio cambia di \(\pm \alpha_R/4\), dove \(\alpha_R\) e` il passo dentale del rotore. Il segno dipende dalla polarita` della fase alimentata e dalla direzione della sequenza.

L'esempio con due fasi e tre denti rotorici serve a visualizzare il meccanismo. La sequenza e` del tipo:
\[
A^+,\quad B^+,\quad A^-,\quad B^-,
\]
e poi si ripete. Ogni commutazione sposta il rotore verso il successivo allineamento favorevole. Se \(\alpha_R=120^\circ\), il passo di commutazione e`:
\[
\frac{\alpha_R}{4}=\frac{120^\circ}{4}=30^\circ.
\]

\slidepair{images13/part1-11.png}{images13/part1-12.png}{Regola del passo per lo stepper ibrido e ruolo della dentatura nella riduzione dell'angolo di passo.}

Per uno stepper ibrido con \(m\) fasi e \(N_r\) denti rotorici, la dimensione del passo intero e`:
\[
\alpha_s=\frac{360^\circ}{2mN_r}.
\]
Il fattore \(2\) compare perche' il rotore ibrido ha due sezioni dentate, sfalsate di mezzo passo. Nell'esempio classico:
\[
m=2,
\qquad
N_r=50,
\]
da cui:
\[
\alpha_s=\frac{360^\circ}{2\cdot 2\cdot 50}=1.8^\circ.
\]
Questo corrisponde a:
\[
\frac{360^\circ}{1.8^\circ}=200
\]
passi per giro.

Il messaggio progettuale e` importante: ridurre il passo non richiede necessariamente piu` fasi elettriche. Si puo` lavorare sulla forma meccanica dei denti di statore e rotore. In altre parole, l'hybrid stepper usa la geometria come moltiplicatore di risoluzione.

\slidetriple{images13/part1-13.png}{images13/part1-14.png}{images13/part1-15.png}{Sequenza di funzionamento dello stepper ibrido: stato iniziale e primi spostamenti del rotore.}

Le immagini successive mostrano una sequenza reale di funzionamento. I denti bianchi anteriori sono associati a un polo, quelli blu posteriori al polo opposto. Alimentando la fase \(B\), per esempio, un lato dello statore diventa sud e l'altro nord; i denti rotorici compatibili vengono attratti verso i poli opposti.

Quando si passa alla fase \(A\), il campo statorico cambia orientamento. I denti che prima erano allineati non sono piu` nella posizione di minima energia, e il rotore avanza di un altro quarto di passo dentale. Il processo continua invertendo opportunamente la polarita` delle fasi:
\[
B,\quad A,\quad -B,\quad -A,\quad B,\ldots
\]
oppure con una notazione equivalente legata ai terminali delle fasi.

\slidepair{images13/part1-16.png}{images13/part1-17.png}{Dopo quattro commutazioni il rotore ha avanzato di un passo dentale; ogni singolo stato produce un quarto di passo.}

Il punto da fissare e` che quattro stati elettrici consecutivi portano il rotore ad avanzare di un passo dentale completo. Ogni commutazione produce quindi un quarto del passo dei denti. Ripetendo la sequenza si ottiene rotazione continua; invertendo l'ordine degli stati si ottiene il verso opposto.

Questa descrizione aiuta anche a capire perche' lo stepper non e` un motore ``digitale ideale''. La sequenza comanda posizioni di equilibrio, ma il rotore ha inerzia, smorzamento, attrito e un carico applicato. Tra una posizione e la successiva c'e` una dinamica meccanica vera, con oscillazioni possibili e con il rischio di non raggiungere l'allineamento prima della commutazione successiva.

\subsection{Problemi pratici degli stepper: perdita di passo e coppia disponibile}

\slidepair{images13/part1-18.png}{images13/part1-19.png}{Gli stepper sono comodi in anello aperto, ma il rotore puo` perdere sincronismo se carico, inerzia o velocita` sono eccessivi.}

Gli stepper sono interessanti perche' possono essere pilotati in anello aperto. Il controllore manda impulsi e assume che ogni impulso produca un passo meccanico. Questa proprieta` semplifica molto il sistema: non serve necessariamente un encoder per conoscere la posizione comandata.

La stessa proprieta`, pero`, e` anche il limite principale. Il controllore conosce il numero di passi richiesti, ma non misura automaticamente se il rotore li ha davvero eseguiti. Se il carico e` troppo alto, se l'inerzia e` elevata, se l'attrito e` sfavorevole o se la frequenza di passo cresce troppo rapidamente, il rotore puo` non seguire la sequenza. Si parla allora di \emph{loss of synchronism} o perdita di passo.

Dal punto di vista statico, lo stepper produce una coppia di richiamo solo se il rotore e` spostato dalla posizione di equilibrio istantanea. Una modellazione locale molto usata e`:
\[
T=T_{\max}\sin\delta,
\]
dove \(\delta\) e` il \emph{load angle}, cioe` lo scostamento rispetto all'equilibrio comandato. Se il carico richiede una coppia troppo alta, \(\delta\) cresce oltre la zona sostenibile e il rotore salta verso un equilibrio vicino.

Le slide distinguono due grandezze:
\begin{itemize}
\item \textbf{holding torque}, la massima coppia statica che il motore alimentato puo` sopportare senza abbandonare il passo comandato;
\item \textbf{detent torque}, la coppia residua a motore non alimentato, dovuta soprattutto ai magneti permanenti e alla geometria. E` significativa nei PM e negli hybrid, molto piccola nei VR.
\end{itemize}

\subsection{Pull-in, pull-out e limiti dinamici di coppia}

\slidepair{images13/part1-20.png}{images13/part1-21.png}{Pull-in torque e riduzione della coppia media al crescere della frequenza di passo.}

La \emph{pull-in torque} e` la massima coppia di carico con cui il motore puo` partire da fermo, sincronizzarsi immediatamente con una frequenza di impulsi applicata e fermarsi di nuovo senza perdere passi. E` quindi un limite di avviamento diretto: il rotore deve accelerare abbastanza rapidamente da raggiungere ogni nuovo equilibrio prima che arrivi il comando successivo.

La slide collega questo limite anche alla dinamica elettrica della fase. La corrente in una bobina non cresce istantaneamente, perche' l'avvolgimento ha induttanza e perche' compare back-EMF. In forma qualitativa:
\[
v=Ri+L\frac{di}{dt}+e_{\text{bemf}}.
\]
A bassa frequenza di passo, la corrente ha tempo di salire quasi fino al valore massimo, quindi il flusso e la coppia sono elevati. A frequenza alta, invece, la fase viene commutata prima che la corrente abbia raggiunto il valore desiderato: il flusso medio diminuisce e la coppia disponibile cala.

\slidepair{images13/part1-22.png}{images13/part1-23.png}{Pull-out torque e curva coppia-frequenza: zona start/stop, slew range e massima frequenza di marcia.}

La \emph{pull-out torque} e` diversa: e` la massima coppia che il motore puo` sostenere mentre sta gia` ruotando in sincronismo a una certa frequenza di passo. In genere, alla stessa frequenza:
\[
T_{\text{pull-out}}>T_{\text{pull-in}}.
\]
Il motivo e` che un rotore gia` in moto ha energia cinetica e sta gia` seguendo la sequenza; non deve partire istantaneamente da fermo. Puo` quindi tollerare una coppia maggiore rispetto al caso di avvio diretto.

Il grafico di coppia-frequenza separa due regioni operative. La regione \emph{start/stop}, sotto la curva di pull-in, e` quella in cui il motore puo` partire e fermarsi direttamente senza rampe. La regione tra pull-in e pull-out viene spesso chiamata \emph{slew range}: il motore puo` lavorarci, ma deve entrarci tramite una rampa di accelerazione. Se si supera la curva di pull-out, il motore perde sincronismo.

Questo e` un punto molto pratico: uno stepper non va comandato solo scegliendo una frequenza finale, ma anche progettando il profilo di accelerazione. Un salto brusco a una frequenza alta puo` fallire anche se, a regime, quella stessa frequenza sarebbe sostenibile.

\begin{keybox}
\textbf{Regola pratica.} Negli stepper la coppia non e` costante con la velocita`. Crescendo la frequenza di passo, la corrente ha meno tempo per stabilizzarsi, la back-EMF aumenta e la coppia media diminuisce. Per questo servono rampe e margine rispetto alle curve di pull-in e pull-out.
\end{keybox}

\subsection{Full-step, half-step, microstepping e confronto tra famiglie}

\slidepair{images13/part1-24.png}{images13/part1-25.png}{Modalita` di passo e confronto qualitativo tra stepper PM, VR e Hybrid.}

La slide finale sugli stepper introduce tre modi di pilotaggio.

Nel \emph{full-step} il rotore avanza del passo base:
\[
\alpha_{\text{full}}=\alpha_s.
\]
Nel \emph{half-step} si eccitano opportunamente due fasi insieme, creando un equilibrio intermedio tra due posizioni di full-step:
\[
\alpha_{\text{half}}=\frac{\alpha_s}{2}.
\]
Nel \emph{microstepping}, invece, le correnti di fase vengono modulate in modo approssimativamente sinusoidale. L'obiettivo e` far muovere il campo risultante in incrementi piu` piccoli:
\[
\alpha_{\mu\text{step}}\approx \frac{\alpha_s}{N_\mu},
\]
dove \(N_\mu\) e` il numero di microstep per passo intero.

Il microstepping non aumenta magicamente la precisione assoluta del motore: riduce vibrazioni, risonanze e ripple di coppia, e permette un moto piu` fluido. La precisione reale dipende comunque da attriti, carico, linearita` delle correnti, dentatura, driver e posizione effettiva del rotore.

Il confronto conclusivo tra PM, VR e Hybrid puo` essere letto cosi`:
\begin{itemize}
\item il PM ha buona coppia di mantenimento, ma passi piu` grandi;
\item il VR e` semplice e puo` raggiungere buone velocita`, ma ha coppia piu` bassa e detent torque quasi nullo;
\item l'Hybrid offre i passi piu` piccoli, la coppia migliore e il microstepping piu` efficace.
\end{itemize}

Per applicazioni di posizionamento generiche, l'hybrid e` spesso la scelta piu` equilibrata. Il VR resta interessante quando contano semplicita`, robustezza e velocita`, mentre il PM e` adatto a soluzioni economiche e meno spinte.

\subsection{Motore BLDC: motivazione e componenti principali}

\slidetriple{images13/part2-01.png}{images13/part2-02.png}{images13/part2-03.png}{Introduzione ai motori brushless DC, statore trifase, rotore a magneti permanenti e commutazione elettronica.}

La seconda parte della lezione apre i \emph{brushless DC motor}. Il nome e` gia` una descrizione del salto rispetto al motore DC brushed: le spazzole e il commutatore meccanico scompaiono. Questo porta vantaggi molto concreti:
\begin{itemize}
\item maggiore efficienza;
\item alta densita` di coppia;
\item minore manutenzione;
\item migliore capacita` di lavorare ad alte velocita`;
\item minore usura meccanica.
\end{itemize}

Per questo i BLDC sono molto diffusi in droni, veicoli elettrici, robotica, automazione industriale, pompe e compressori. In pratica, un BLDC trifase viene alimentato da un inverter che applica tensioni alle tre fasi in funzione della posizione del rotore.

La struttura base contiene uno statore con tre avvolgimenti, indicati come \(a,b,c\), e un rotore con magneti permanenti. Le slide assumono una connessione a stella. In questo caso, per la legge di Kirchhoff sul nodo di stella:
\[
i_a+i_b+i_c=0.
\]
Questa relazione riduce da tre a due il numero di correnti indipendenti. E` una semplificazione importante nel controllo, perche' la terza corrente e` determinata dalle altre due.

Il principio fisico resta vicino a quello gia` visto nei motori elettrici: le correnti statoriche generano un campo magnetico rotante, i magneti permanenti del rotore tendono ad allinearsi con quel campo e da questo allineamento nasce coppia. La differenza decisiva e` che la commutazione non e` piu` meccanica. Deve essere realizzata elettronicamente dal driver.

Le slide notano anche che, in molti contesti, si parla sia di BLDC sia di PMSM, \emph{Permanent Magnet Synchronous Motor}. La distinzione dettagliata e` fuori dallo scopo della lezione; qui interessa soprattutto il modello di macchina trifase a magneti permanenti e il problema della commutazione.

\subsection{Equazioni elettriche trifase e flussi concatenati}

\slidepair{images13/part2-04.png}{images13/part2-05.png}{Equazioni di fase del BLDC e modello dei flussi concatenati.}

Per ogni fase si scrive un'equazione di tensione:
\[
v_a=R_s i_a+\frac{d\lambda_a}{dt},
\qquad
v_b=R_s i_b+\frac{d\lambda_b}{dt},
\qquad
v_c=R_s i_c+\frac{d\lambda_c}{dt}.
\]
In forma vettoriale:
\[
\mathbf v_{abc}
=R_s\mathbf i_{abc}
+\frac{d\boldsymbol\lambda_{abc}}{dt}.
\]

Il termine \(\boldsymbol\lambda_{abc}\) contiene i flussi concatenati delle tre fasi. Le slide lo scompongono in tre contributi:
\begin{itemize}
\item autoinduttanza della fase;
\item mutue induttanze tra fasi;
\item flusso prodotto dai magneti permanenti del rotore.
\end{itemize}

Un modello compatto e`:
\[
\boldsymbol\lambda_{abc}
=\mathbf L_{abc}\mathbf i_{abc}
+\boldsymbol\lambda_m(\theta_e),
\]
dove \(\theta_e\) e` la posizione elettrica del rotore. La posizione elettrica non coincide sempre con la posizione meccanica: se il rotore ha \(p\) coppie polari,
\[
\theta_e=p\theta_m.
\]
Questo significa che, per un giro meccanico del rotore, le grandezze elettriche possono compiere piu` periodi.

Per una macchina simmetrica con magneti superficiali, la matrice di induttanza puo` essere approssimata come:
\[
\mathbf L_{abc}=
\begin{bmatrix}
L & M & M\\
M & L & M\\
M & M & L
\end{bmatrix}.
\]
La parte davvero delicata e` il flusso dei magneti \(\boldsymbol\lambda_m(\theta_e)\), perche' dipende dalla posizione del rotore. Da qui nasce la necessita` di conoscere, o stimare, la posizione rotorica per controllare correttamente il motore.

\subsection{Back-EMF di fase e modello ridotto}

\slidetriple{images13/part2-06.png}{images13/part2-07.png}{images13/part2-08.png}{Dalle equazioni complete alla forma ridotta con induttanza equivalente e back-EMF di fase.}

Derivando il flusso concatenato si ottengono termini dovuti alla variazione delle correnti e termini dovuti alla variazione del flusso dei magneti. Questi ultimi sono le back-EMF di fase:
\[
e_a=\frac{d\lambda_{m,a}}{dt},
\qquad
e_b=\frac{d\lambda_{m,b}}{dt},
\qquad
e_c=\frac{d\lambda_{m,c}}{dt}.
\]
Sono l'analogo trifase della back-EMF del motore DC brushed: quando il rotore si muove, i magneti permanenti inducono tensioni nelle fasi statoriche.

Usando il vincolo della connessione a stella:
\[
i_a+i_b+i_c=0
\qquad \Rightarrow \qquad
\dot i_a+\dot i_b+\dot i_c=0,
\]
i termini di mutua induttanza possono essere riordinati. Per la fase \(a\), per esempio, dalla simmetria con mutua \(M\) si ottiene:
\[
M\dot i_b+M\dot i_c=-M\dot i_a.
\]
Quindi:
\[
v_a=R_s i_a+(L_s-M)\dot i_a+e_a.
\]
Definendo l'induttanza equivalente:
\[
L_{eq}=L_s-M,
\]
il modello ridotto di fase diventa:
\[
\begin{cases}
v_a=R_s i_a+L_{eq}\dot i_a+e_a,\\
v_b=R_s i_b+L_{eq}\dot i_b+e_b,\\
v_c=R_s i_c+L_{eq}\dot i_c+e_c.
\end{cases}
\]

Questa forma e` molto utile per il controllo: ogni fase assomiglia a un circuito \(R\)-\(L\) con una sorgente di back-EMF in serie. Pero` le tre fasi restano collegate dal vincolo \(i_a+i_b+i_c=0\) e dalla dipendenza delle back-EMF dalla posizione del rotore.

\subsection{Coppia elettromagnetica e dinamica meccanica del BLDC}

\slidepair{images13/part2-09.png}{images13/part2-10.png}{Generazione della coppia BLDC tramite bilancio di potenza e dinamica meccanica del rotore.}

La coppia viene ricavata con un ragionamento energetico. La potenza elettromagnetica associata alle back-EMF e`:
\[
p_e=e_a i_a+e_b i_b+e_c i_c.
\]
Trascurando le perdite, questa potenza coincide con la potenza meccanica convertita:
\[
p_e=p_m=T_e\omega_m.
\]
Quindi:
\[
T_e=
\frac{e_a i_a+e_b i_b+e_c i_c}{\omega_m}.
\]

Questa formula e` molto istruttiva: la coppia non dipende solo dal modulo delle correnti, ma da come le correnti vengono allineate con le back-EMF. Se una fase viene alimentata quando la sua back-EMF e` favorevole, contribuisce positivamente alla coppia; se viene alimentata nel momento sbagliato, puo` produrre ripple, perdita di efficienza o addirittura frenatura.

La dinamica meccanica del rotore e`:
\[
J\frac{d\omega_m}{dt}
=T_e-T_L-B\omega_m,
\]
dove \(J\) e` l'inerzia totale, \(T_L\) la coppia di carico e \(B\) il coefficiente di attrito viscoso. La posizione meccanica evolve come:
\[
\frac{d\theta_m}{dt}=\omega_m.
\]
La posizione elettrica, come gia` visto, e`:
\[
\theta_e=p\theta_m.
\]

\slidefig{images13/part2-11.png}{Sintesi del modello BLDC: equazioni elettriche, dinamica meccanica e dipendenza dei flussi dalla posizione rotorica.}

Il riepilogo delle slide mette insieme le equazioni principali:
\[
\mathbf v_{abc}
=R_s\mathbf i_{abc}
+\frac{d\boldsymbol\lambda_{abc}}{dt},
\]
\[
J\frac{d\omega_m}{dt}=T_e-T_L-B\omega_m,
\qquad
\omega_m=\frac{d\theta_m}{dt},
\qquad
\theta_e=p\theta_m,
\]
con:
\[
\boldsymbol\lambda_{abc}
=\mathbf L_{abc}\mathbf i_{abc}
+\boldsymbol\lambda_m(\theta_e).
\]

La frase piu` importante e` che \(\boldsymbol\lambda_{abc}\) non e` costante: dipende dalla posizione rotorica. Di conseguenza, il controllo in anello chiuso di un BLDC deve conoscere la posizione del rotore. Questa informazione puo` arrivare da sensori Hall, encoder, resolver, oppure da metodi sensorless che stimano la posizione a partire da correnti e tensioni.

\begin{keybox}
\textbf{Differenza rispetto al brushed DC.} Nel motore brushed la commutazione viene fatta meccanicamente dal collettore. Nel BLDC la stessa funzione deve essere svolta dall'elettronica di potenza, che deve sapere dove si trova il rotore per alimentare le fasi giuste nel momento giusto.
\end{keybox}

\subsection{Controllo BLDC e commutazione trapezoidale}

\slidepair{images13/part2-12.png}{images13/part2-13.png}{Controllo BLDC e idea della back-EMF trapezoidale: attivare le fasi quando la back-EMF e` sul plateau.}

Controllare un BLDC significa scegliere le tensioni di fase in modo da ottenere le correnti necessarie e quindi la coppia desiderata alla velocita` desiderata. In forma concettuale, l'inverter deve imporre \(\mathbf v_{abc}\) per ottenere \(\mathbf i_{abc}\), sapendo pero` che le back-EMF e i flussi dipendono da \(\theta_e\).

Questo puo` diventare complesso. Per questo alcune macchine sono progettate in modo che la forma della back-EMF renda il controllo piu` semplice. Le slide distinguono:
\begin{itemize}
\item BLDC con back-EMF trapezoidale;
\item macchine con back-EMF sinusoidale, spesso trattate con controllo sinusoidale o vettoriale.
\end{itemize}

Nel caso trapezoidale, le tre back-EMF hanno forme quasi trapezoidali sfasate di \(120^\circ\) elettrici. L'idea e` alimentare le fasi quando le rispettive back-EMF sono sui tratti piatti, cioe` quando il contributo alla coppia e` quasi costante. Tipicamente bastano sei stati di commutazione per giro elettrico.

\slidefig{images13/part2-14.png}{Nel BLDC trapezoidale si alimentano due fasi alla volta e si lascia la terza non alimentata, ottenendo una coppia idealmente quasi piatta.}

In un settore della sequenza, per esempio, le back-EMF delle fasi \(A\) e \(B\) possono essere sui plateau, mentre la fase \(C\) attraversa il tratto variabile. Una scelta tipica e`:
\[
i_A=I,
\qquad
i_B=-I,
\qquad
i_C=0.
\]
La fase \(C\) viene lasciata flottante o comunque non usata per produrre coppia in quel settore. Passando al settore successivo, cambia la coppia di fasi attive. Dopo sei settori la sequenza si ripete.

Questa tecnica viene chiamata spesso \emph{six-step commutation}. E` semplice, robusta e adatta a molti azionamenti, ma produce piu` ripple di coppia rispetto a un controllo sinusoidale ben realizzato. Il suo grande vantaggio e` che richiede una misura di posizione relativamente grossolana: basta sapere in quale dei sei settori elettrici si trova il rotore.

\subsection{Sensori Hall e decodifica dei sei settori}

\slidepair{images13/part2-15.png}{images13/part2-16.png}{Sensori Hall a \(120^\circ\) elettrici e decodifica delle combinazioni logiche per la commutazione a sei passi.}

Per conoscere il settore elettrico del rotore si possono usare tre sensori a effetto Hall, disposti tipicamente a \(120^\circ\) elettrici. Ogni sensore produce un livello logico alto o basso a seconda del polo magnetico che gli passa davanti. La terna di segnali Hall identifica una delle sei zone utili del ciclo elettrico.

Il driver decodifica quindi la combinazione dei tre segnali e sceglie quali fasi alimentare:
\[
(H_A,H_B,H_C)
\longrightarrow
\text{settore elettrico}
\longrightarrow
\text{coppia di fasi attive}.
\]
Nella tabella grafica della slide compaiono anche stati ad alta impedenza, indicati con \(Z\): sono i momenti in cui una fase non viene pilotata direttamente.

Questa architettura rende il BLDC molto diverso dallo stepper. Nello stepper si invia una sequenza e si spera che il rotore rimanga agganciato al campo; nel BLDC, invece, la sequenza di commutazione viene sincronizzata con la posizione effettiva del rotore. Il motore resta sincrono per costruzione: se il rotore accelera o rallenta, l'elettronica deve adattare il momento in cui cambia settore.

\subsection*{Sintesi finale della Lezione 14}

La prima parte della lezione completa gli stepper. Nel VR stepper la coppia nasce dalla minimizzazione della riluttanza, e il passo dipende dalla differenza tra denti statorici e rotorici:
\[
\alpha_s
=360^\circ
\left|
\frac{N_s-N_r}{N_sN_r}
\right|.
\]
Nello stepper ibrido, invece, magnete permanente e riluttanza lavorano insieme; la dentatura permette passi piccoli, con:
\[
\alpha_s=\frac{360^\circ}{2mN_r}.
\]
La trattazione pratica chiarisce poi che lo stepper non e` ideale: holding torque, detent torque, pull-in, pull-out, frequenza di passo, induttanza degli avvolgimenti e profili di accelerazione determinano se il motore resta sincronizzato o perde passi.

La seconda parte apre i BLDC. Il modello trifase parte da:
\[
\mathbf v_{abc}
=R_s\mathbf i_{abc}
+\frac{d\boldsymbol\lambda_{abc}}{dt},
\qquad
\boldsymbol\lambda_{abc}
=\mathbf L_{abc}\mathbf i_{abc}
+\boldsymbol\lambda_m(\theta_e).
\]
La coppia si legge dal bilancio di potenza:
\[
T_e=
\frac{e_a i_a+e_b i_b+e_c i_c}{\omega_m},
\]
mentre la dinamica meccanica e`:
\[
J\dot\omega_m=T_e-T_L-B\omega_m.
\]

Il punto concettuale piu` importante e` che il BLDC elimina il commutatore meccanico, ma non elimina il problema della commutazione: lo sposta nell'elettronica. Per produrre coppia in modo efficiente, il driver deve conoscere la posizione elettrica del rotore e alimentare le fasi coerentemente con le back-EMF. Nel BLDC trapezoidale questa idea diventa una commutazione a sei passi, spesso basata su tre sensori Hall.

% =========================================================
% LEZIONE 15
% =========================================================

\clearpage
\section{Lezione 15}

\subsection{Visione generale della lezione}

Questa lezione chiude il discorso sui brushless e introduce poi il problema pratico dei driver. Il filo logico e` questo: nel BLDC trapezoidale la commutazione a sei passi e` semplice, ma nella macchina reale lascia ripple di coppia. Se invece la macchina ha back-EMF sinusoidale, si puo` comandarla con correnti sinusoidali e ottenere una coppia piu` regolare. Per farlo in modo controllabile si usano le trasformate di Clarke e Park, che portano naturalmente al \emph{Field Oriented Control}.

Nella seconda parte si cambia livello: non si guarda piu` solo il modello del motore, ma il modo in cui l'elettronica genera fisicamente le tensioni richieste dal controllo. Il caso usato come riferimento e` il motore DC brushed pilotato con ponte H e PWM, perche' permette di capire in modo molto diretto il legame tra tensione applicata, corrente, coppia e verso di rotazione.

\subsection{Dal BLDC trapezoidale al controllo sinusoidale}

\slidefig{images14/page-01.png}{Slide 414: nel BLDC trapezoidale reale la commutazione delle fasi non e` istantanea, quindi la coppia non resta perfettamente piatta.}

Nella lezione precedente il controllo trapezoidale era stato descritto nel caso ideale: due fasi attive, una fase lasciata inattiva, back-EMF trapezoidale e correnti quasi rettangolari. L'idea e` alimentare le fasi proprio quando le rispettive back-EMF sono sul plateau, in modo che la potenza elettromagnetica
\[
p_e=e_a i_a+e_b i_b+e_c i_c
\]
resti il piu` possibile costante. Poiche' la coppia si ottiene da
\[
T_e=\frac{p_e}{\omega_m},
\]
una potenza convertita quasi costante produce una coppia quasi costante.

La slide mette pero` in evidenza il limite pratico: le correnti di fase non possono cambiare istantaneamente. Ogni fase e` un circuito con resistenza, induttanza e back-EMF:
\[
v=Ri+L\frac{di}{dt}+e_{\text{bemf}}.
\]
Quando il driver cambia settore di commutazione, una corrente deve scendere, un'altra deve salire e la terza cambia ruolo. Durante questo intervallo il prodotto \(e_k i_k\) delle varie fasi non mantiene il valore ideale, quindi la coppia ha ondulazioni.

Il punto da fissare e` che la commutazione trapezoidale e` ottima quando contano semplicita`, robustezza e costo, ma non quando serve moto molto fluido. Il ripple alla frequenza di commutazione puo` diventare vibrazione, rumore acustico e peggioramento della regolazione a bassa velocita`.

\slidepair{images14/page-02.png}{images14/page-03.png}{Slide 420--421: motivazione del controllo sinusoidale e prima trasformazione, da \(abc\) al piano statorico \(\alpha\beta\).}

La slide 420 non dice semplicemente di usare onde sinusoidali al posto di onde rettangolari. Il punto e` piu` preciso: se il motore e` costruito in modo che la back-EMF sia sinusoidale in funzione dell'angolo rotorico, allora conviene generare correnti sinusoidali opportune. In questo modo il prodotto tra back-EMF e corrente puo` rimanere molto regolare, e quindi anche la coppia migliora.

I vantaggi sono quelli indicati nelle slide:
\begin{itemize}
\item minore ripple di coppia;
\item controllo piu` fine a bassa velocita`;
\item migliore risposta dinamica.
\end{itemize}
Il prezzo da pagare e` la complessita` del controllo. Nel riferimento fisico delle tre fasi, infatti, le correnti desiderate sono sinusoidi sfasate di \(120^\circ\) elettrici. Un regolatore che dovesse inseguire direttamente tre sinusoidi variabili con la posizione rotorica sarebbe scomodo e poco trasparente.

Qui entrano le due trasformazioni della lezione:
\[
abc \longrightarrow \alpha\beta \longrightarrow dq.
\]
La trasformata di Clarke sostituisce le tre grandezze di fase con due componenti ortogonali fisse allo statore. La trasformata di Park ruota poi queste due componenti nel riferimento solidale al rotore. L'obiettivo finale e` trasformare tre sinusoidi in due grandezze quasi costanti, molto piu` facili da regolare.

\subsection{Trasformata di Clarke: dalle tre fasi al vettore statorico}

\slidefig{images14/page-04.png}{Slide 422: la Clarke rappresenta le tre fasi con un unico vettore equivalente nel piano statorico \(\alpha\beta\).}

Per un sistema trifase bilanciato vale:
\[
x_a+x_b+x_c=0.
\]
Questo significa che le tre variabili non sono indipendenti: conoscendone due, la terza e` fissata. La Clarke sfrutta proprio questa ridondanza e proietta la terna \(abc\) su due assi ortogonali fissi allo statore, \(\alpha\) e \(\beta\). Per una generica grandezza \(x\), che puo` essere corrente, tensione o flusso concatenato:
\[
\begin{bmatrix}
x_\alpha\\[0.2em]
x_\beta
\end{bmatrix}
=
\frac{2}{3}
\begin{bmatrix}
1 & -\frac{1}{2} & -\frac{1}{2}\\[0.3em]
0 & \frac{\sqrt 3}{2} & -\frac{\sqrt 3}{2}
\end{bmatrix}
\begin{bmatrix}
x_a\\[0.2em]
x_b\\[0.2em]
x_c
\end{bmatrix}.
\]
Nel caso delle correnti:
\[
\begin{bmatrix}
i_\alpha\\[0.2em]
i_\beta
\end{bmatrix}
=
\frac{2}{3}
\begin{bmatrix}
1 & -\frac{1}{2} & -\frac{1}{2}\\[0.3em]
0 & \frac{\sqrt 3}{2} & -\frac{\sqrt 3}{2}
\end{bmatrix}
\begin{bmatrix}
i_a\\[0.2em]
i_b\\[0.2em]
i_c
\end{bmatrix}.
\]

Geometricamente, le correnti di fase generano un campo statorico risultante. Invece di seguire separatamente \(i_a,i_b,i_c\), si segue il vettore:
\[
\mathbf i_s=
\begin{bmatrix}
i_\alpha\\ i_\beta
\end{bmatrix}.
\]
Questo vettore vive nel piano dello statore. Se il motore gira a velocita` costante e le correnti sono sinusoidali, \(\mathbf i_s\) ruota nel piano \(\alpha\beta\). La Clarke quindi semplifica il problema da tre variabili a due, ma non elimina ancora la rotazione: le componenti \(i_\alpha\) e \(i_\beta\) restano variabili nel tempo.

\begin{keybox}
\textbf{Lettura della Clarke.} Non e` una magia matematica fine a se stessa: e` il passaggio dalla descrizione ``per fase'' alla descrizione ``per vettore di campo statorico''. Da qui in avanti si controlla dove punta e quanto vale il vettore di corrente prodotto dallo statore.
\end{keybox}

\subsection{Trasformata di Park e modello nel riferimento dq}

\slidepair{images14/page-05.png}{images14/page-06.png}{Slide 423--424: la Park ruota il riferimento dell'angolo elettrico \(\theta_e\) e permette di riscrivere le equazioni nel riferimento \(dq\).}

La trasformata di Park compie il passo decisivo: osserva il vettore statorico da un riferimento che ruota con il rotore. La forma usata nelle slide e`:
\[
\begin{bmatrix}
x_d\\[0.2em]
x_q
\end{bmatrix}
=
\begin{bmatrix}
\cos\theta_e & \sin\theta_e\\[0.2em]
-\sin\theta_e & \cos\theta_e
\end{bmatrix}
\begin{bmatrix}
x_\alpha\\[0.2em]
x_\beta
\end{bmatrix}.
\]
Applicata alle correnti:
\[
\begin{bmatrix}
i_d\\[0.2em]
i_q
\end{bmatrix}
=
\begin{bmatrix}
\cos\theta_e & \sin\theta_e\\[0.2em]
-\sin\theta_e & \cos\theta_e
\end{bmatrix}
\begin{bmatrix}
i_\alpha\\[0.2em]
i_\beta
\end{bmatrix}.
\]

Il riferimento \(dq\) ruota sincronicamente con il campo magnetico rotorico. Con la convenzione adottata nella slide 423:
\[
\text{asse }d \parallel \text{flusso del rotore},
\qquad
\text{asse }q \perp \text{flusso del rotore}.
\]
Questa distinzione e` il cuore del controllo vettoriale. La corrente \(i_d\) e` la componente lungo il flusso, mentre \(i_q\) e` la componente in quadratura, cioe` quella che produce coppia nella macchina a magneti superficiali.

Nel riferimento \(dq\), le equazioni di tensione assumono la forma:
\[
v_d=R_s i_d+L_d\frac{di_d}{dt}-\omega_e L_q i_q,
\]
\[
v_q=R_s i_q+L_q\frac{di_q}{dt}+\omega_e L_d i_d+\omega_e\lambda_m.
\]
Qui \(L_d\) e \(L_q\) sono le induttanze viste lungo gli assi \(d\) e \(q\), \(\omega_e=\dot\theta_e\) e` la velocita` elettrica, e \(\lambda_m\) e` il flusso concatenato dei magneti permanenti.

Queste equazioni sono piu` ricche di due semplici circuiti \(R\)-\(L\). I termini
\[
-\omega_e L_q i_q,
\qquad
\omega_e L_d i_d,
\qquad
\omega_e\lambda_m
\]
nascono dal fatto che il riferimento sta ruotando. In particolare, il flusso dei magneti e` fisso nel riferimento del rotore, ma la sua variazione vista dallo statore genera back-EMF; per questo il termine \(\omega_e\lambda_m\) compare nell'equazione \(q\). Detto in modo operativo: nel riferimento \(dq\) il flusso non oscilla piu` con \(\theta_e\), ma la velocita` di rotazione continua a entrare nelle tensioni.

\slidepair{images14/page-07.png}{images14/page-08.png}{Slide 425--426: nel riferimento \(dq\) il flusso rotorico diventa fisso e le sinusoidi a regime diventano grandezze costanti.}

Le slide 425 e 426 spiegano perche' questo cambio di coordinate e` utile. Nel riferimento \(abc\) le correnti sono tre sinusoidi. Nel riferimento \(\alpha\beta\) sono un vettore rotante. Nel riferimento \(dq\), se il riferimento ruota alla stessa velocita` elettrica del rotore, quel vettore viene visto quasi fermo.

Questo e` analogo a osservare un oggetto che gira mettendosi su una piattaforma che gira con lui: cio` che prima cambiava continuamente diventa quasi costante. Per il controllo e` una semplificazione enorme, perche' un regolatore PI lavora molto meglio su riferimenti costanti che su sinusoidi.

La lettura fisica proposta dalle slide e` quella di due avvolgimenti virtuali:
\begin{itemize}
\item uno sull'asse \(d\), allineato al flusso rotorico;
\item uno sull'asse \(q\), ortogonale al flusso rotorico.
\end{itemize}
Nella macchina a magneti superficiali ideale, la corrente sull'asse \(d\) non produce coppia utile; la coppia viene dalla corrente in quadratura \(i_q\). Il termine
\[
L_d i_d+\lambda_m
\]
si puo` leggere come il flusso complessivo concatenato all'asse \(d\): contiene il flusso dei magneti e l'eventuale contributo prodotto da \(i_d\). Nella strategia base si evita di usare \(i_d\), perche' non conviene spendere corrente in una componente che non genera direttamente coppia.

\subsection{Coppia nel riferimento dq e idea del FOC}

\slidepair{images14/page-09.png}{images14/page-10.png}{Slide 427--428: la coppia nel riferimento \(dq\) porta direttamente alla struttura del Field Oriented Control.}

La coppia elettromagnetica nel riferimento \(dq\) e`:
\[
T_e=\frac{3}{2}p\left(\lambda_m i_q+(L_d-L_q)i_di_q\right),
\]
dove \(p\) e` il numero di coppie polari. La formula contiene due contributi. Il primo,
\[
\frac{3}{2}p\lambda_m i_q,
\]
e` la coppia dovuta all'interazione tra il flusso dei magneti permanenti e la corrente in quadratura. Il secondo,
\[
\frac{3}{2}p(L_d-L_q)i_di_q,
\]
e` la coppia di riluttanza, presente quando la macchina ha salienza, cioe` quando \(L_d\) e \(L_q\) sono diversi.

Per un motore a magneti superficiali si assume spesso:
\[
L_d\simeq L_q,
\]
percio` il secondo termine scompare e rimane:
\[
T_e=\frac{3}{2}p\lambda_m i_q.
\]
Questa e` la formula piu` importante della parte FOC: a flusso magnetico fissato, la coppia e` proporzionale a \(i_q\). Quindi, se si vuole una certa coppia \(T_e^\star\), si sceglie un riferimento \(i_q^\star\) proporzionale a quella coppia. Nel caso base, per massimizzare la coppia prodotta a parita` di corrente, si pone:
\[
i_d^\star=0,
\qquad
i_q^\star\propto T_e^\star.
\]

Questa e` la preparazione al \emph{Field Oriented Control}. Il FOC non controlla direttamente tre correnti sinusoidali. Misura le correnti reali, le trasforma in \(dq\), controlla \(i_d\) e \(i_q\), poi trasforma indietro le tensioni richieste. La catena operativa e`:
\begin{itemize}
\item misura o stima di \(i_a,i_b,i_c\), spesso con due soli sensori perche' \(i_a+i_b+i_c=0\);
\item Clarke: \(abc\rightarrow\alpha\beta\);
\item Park: \(\alpha\beta\rightarrow dq\);
\item regolazione di \(i_d\) e \(i_q\) con controllori di corrente;
\item calcolo delle tensioni \(v_d^\star,v_q^\star\);
\item Park inversa e Clarke inversa per tornare alle tensioni di fase;
\item generazione pratica delle tensioni con l'inverter.
\end{itemize}

\begin{keybox}
\textbf{Idea essenziale del FOC.} Il motore brushless sinusoidale viene fatto assomigliare, dal punto di vista del controllo, a un motore DC: una componente di corrente gestisce il flusso, l'altra gestisce la coppia. Nel caso a magneti superficiali, il flusso e` gia` fornito dai magneti e la coppia si controlla quasi tutta con \(i_q\).
\end{keybox}

\subsection{Schema completo del FOC}

\slidetriple{images14/page-11.png}{images14/page-12.png}{images14/page-13.png}{Slide 429--431: confronto tra riferimenti, trasformate inverse e schema a blocchi completo del FOC con SVPWM.}

La slide 429 serve a fissare visivamente il risultato delle trasformate. Le stesse correnti possono essere guardate in tre modi:
\begin{itemize}
\item in \(abc\): tre sinusoidi di fase;
\item in \(\alpha\beta\): un vettore rotante nello statore;
\item in \(dq\): due componenti quasi costanti, se la trasformata usa l'angolo elettrico corretto.
\end{itemize}
Questo spiega perche' i regolatori del FOC sono messi nel riferimento \(dq\). Un riferimento di coppia costante diventa un \(i_q^\star\) costante; un riferimento di flusso nullo o fissato diventa un \(i_d^\star\) costante.

La slide 430 mostra il ritorno dal mondo virtuale \(dq\) al mondo fisico delle tre fasi. Si parte da:
\[
i_d^\star=0,
\qquad
i_q^\star\propto T_e^\star.
\]
I regolatori generano \(v_d^\star,v_q^\star\). Queste tensioni non possono essere applicate direttamente, perche' l'inverter e` collegato alle fasi fisiche \(a,b,c\). Si applica allora la Park inversa:
\[
\begin{bmatrix}
x_\alpha\\[0.2em]
x_\beta
\end{bmatrix}
=
\begin{bmatrix}
\cos\theta_e & -\sin\theta_e\\[0.2em]
\sin\theta_e & \cos\theta_e
\end{bmatrix}
\begin{bmatrix}
x_d\\[0.2em]
x_q
\end{bmatrix}.
\]
Poi la Clarke inversa:
\[
\begin{bmatrix}
x_a\\[0.2em]
x_b\\[0.2em]
x_c
\end{bmatrix}
=
\begin{bmatrix}
1 & 0\\[0.2em]
-\frac{1}{2} & \frac{\sqrt 3}{2}\\[0.2em]
-\frac{1}{2} & -\frac{\sqrt 3}{2}
\end{bmatrix}
\begin{bmatrix}
x_\alpha\\[0.2em]
x_\beta
\end{bmatrix}.
\]
Alla fine si ottengono riferimenti di tensione trifase. Durante la rotazione, queste tensioni risultano sinusoidali, ma il regolatore non le ha inseguite come sinusoidi: le ha prodotte trasformando indietro due comandi quasi costanti.

Lo schema della slide 431 aggiunge il dettaglio dell'attuazione. I PI di corrente calcolano le tensioni nel riferimento \(dq\); le trasformate inverse producono un vettore di tensione nello statore; il blocco SVM/SVPWM traduce quel vettore in duty cycle per i MOSFET dell'inverter. La SVM non e` un nuovo modello del motore: e` una tecnica di modulazione per sintetizzare nel modo migliore possibile la tensione richiesta usando un bus DC e interruttori elettronici.

\begin{infobox}
\textbf{Ruolo dell'angolo elettrico.} La Park funziona solo se si conosce \(\theta_e\). Se l'angolo e` sbagliato, il controllore mescola \(i_d\) e \(i_q\): una parte della corrente che si pensava di usare per la coppia finisce sull'asse di flusso, e viceversa.
\end{infobox}

\subsection{Misura dell'angolo elettrico ed esempi di brushless}

\slidepair{images14/page-14.png}{images14/page-15.png}{Slide 432--433: misura dell'angolo elettrico e tipologie costruttive di motori brushless.}

Per usare la trasformata di Park bisogna conoscere l'angolo elettrico:
\[
\theta_e=p\theta_m+\theta_0,
\]
dove \(p\) e` il numero di coppie polari e \(\theta_0\) rappresenta l'offset tra sensore e assi magnetici della macchina. Questo dettaglio e` importante: non basta misurare un angolo meccanico qualsiasi, bisogna sapere come quell'angolo si allinea al flusso rotorico.

Le slide distinguono tre casi:
\begin{itemize}
\item nei BLDC trapezoidali bastano sensori Hall discreti, perche' serve solo identificare il settore di commutazione;
\item nei brushless sinusoidali con FOC serve una misura piu` fine, spesso un encoder magnetico o ottico;
\item nel controllo sensorless la posizione viene stimata dalle grandezze elettriche, in particolare dalla back-EMF.
\end{itemize}
Il limite del sensorless e` chiaro: la back-EMF e` proporzionale alla velocita`. A rotore fermo o a bassissima velocita`, il segnale utile e` debole, quindi la stima della posizione diventa difficile. Per questo applicazioni che richiedono coppia precisa da fermo usano spesso un sensore di posizione.

La slide con gli esempi costruttivi collega il controllo alla meccanica reale. Un brushless industriale puo` integrare direttamente l'encoder, proprio per rendere possibile il FOC. I motori ad albero cavo sono utili quando bisogna far passare cavi o organi meccanici lungo l'asse. Negli \emph{outrunner} il rotore e` esterno e lo statore interno; negli \emph{inrunner} il rotore e` interno e lo statore esterno. Gli outrunner tendono ad avere piu` coppia a basso regime grazie al raggio maggiore, mentre gli inrunner sono spesso piu` adatti ad alte velocita`.

\subsection{Driver: dal segnale di controllo alla potenza}

\slidetriple{images14/page-16.png}{images14/page-17.png}{images14/page-18.png}{Slide 434--436: tutti i motori richiedono tensioni controllate; nel brushed DC il driver regola tensione, corrente, coppia e verso.}

La seconda parte della lezione sposta l'attenzione sui driver. Il controllore non puo` generare direttamente coppia meccanica: puo` generare segnali logici o riferimenti numerici. Il driver trasforma questi segnali in tensioni e correnti di potenza.

Il principio e` comune a tutte le famiglie viste:
\begin{itemize}
\item nel DC brushed, cambiando la tensione di armatura si cambia il punto di lavoro in velocita`, mentre la coppia dipende dalla corrente;
\item nello stepper, il driver eccita le fasi in sequenza per generare posizioni di equilibrio successive;
\item nel brushless, l'inverter genera tensioni trifase per ottenere le correnti che producono coppia.
\end{itemize}

Per il motore DC brushed, un controllore logico non puo` essere collegato direttamente all'armatura. Non ha abbastanza corrente, non gestisce tensioni elevate, non permette in modo diretto inversione e frenatura, e non offre le protezioni richieste da un carico induttivo. Il driver e` quindi l'interfaccia tra controllore, alimentazione DC e motore.

Il modello richiamato dalla slide 436 e`:
\[
v_a=R_a i_a+L_a\frac{di_a}{dt}+k_e\omega,
\qquad
T_e=k_t i_a.
\]
La coppia e` proporzionale alla corrente. La corrente, pero`, dipende dalla tensione applicata al netto della back-EMF \(k_e\omega\). A regime, trascurando l'induttanza:
\[
v_a\simeq R_a i_a+k_e\omega.
\]
Quindi, se il motore ruota gia` velocemente, una parte della tensione disponibile serve solo a compensare la back-EMF. Per ottenere una certa coppia, il driver deve applicare una tensione tale da far circolare la corrente richiesta. Il segno della tensione e della corrente determina il segno della coppia.

\subsection{Ponte H: segno della tensione, frenatura e coasting}

\slidefig{images14/page-19.png}{Slide 437: il ponte H permette di applicare al motore tensione positiva, tensione negativa, frenatura elettrica o coasting.}

Il circuito tipico per pilotare un motore DC in entrambi i versi e` il ponte H. I quattro interruttori permettono di collegare i terminali del motore al bus DC con polarita` diverse. Accendendo una diagonale si ottiene una tensione di un segno:
\[
S_1=1,\ S_4=1 \quad \Rightarrow \quad \text{moto verso destra},
\]
accendendo l'altra diagonale si ottiene il segno opposto:
\[
S_2=1,\ S_3=1 \quad \Rightarrow \quad \text{moto verso sinistra}.
\]
Con tutti gli interruttori spenti il motore va in \emph{coasting}: l'armatura e` sostanzialmente libera e il motore rallenta per attrito e carico. Cortocircuitando opportunamente i terminali si ottiene invece \emph{brake}, cioe` frenatura elettrica. Nella slide compare la parola ``Break'', ma nel contesto il significato e` frenatura, quindi \emph{brake}.

Questa tabella non serve solo a scegliere il verso. Prepara l'idea del PWM: la tensione media non viene regolata con un alimentatore analogico variabile, ma alternando rapidamente stati del ponte. Il motore, grazie alla propria induttanza e inerzia, non vede il dettaglio di ogni commutazione, ma soprattutto il valore medio e il ripple sovrapposto.

\subsection{PWM: tensione media e scelta della frequenza}

\slidepair{images14/page-20.png}{images14/page-21.png}{Slide 438--439: il PWM regola la tensione media tramite il duty cycle; le frequenze tipiche sono dell'ordine di \(10\text{--}40\,\text{kHz}\).}

Nella maggior parte delle applicazioni il bus DC ha valore fissato \(V_{dc}\), mentre il motore richiede una tensione efficace variabile. Il PWM risolve il problema commutando rapidamente tra:
\begin{itemize}
\item uno stato ON, in cui il motore e` collegato all'alimentazione con una certa polarita`;
\item uno stato OFF o di ricircolo, in cui la corrente continua a fluire attraverso i percorsi consentiti dal ponte.
\end{itemize}
Se il periodo di switching e` \(T_s\) e il tempo di accensione e` \(T_{on}\), il duty cycle e`:
\[
D=\frac{T_{on}}{T_s}.
\]
Nel caso piu` semplice, il valore medio della tensione sul motore e`:
\[
v_m\simeq D V_{dc}.
\]
La relazione da ricordare e` questa: il duty cycle non cambia il valore istantaneo del bus DC, ma cambia la frazione di tempo in cui il bus viene applicato al motore. Da qui nasce una tensione media controllabile.

La frequenza di switching e`:
\[
f_s=\frac{1}{T_s}.
\]
Le slide indicano come valore tipico \(10\text{--}40\,\text{kHz}\). Non e` una legge universale: e` un compromesso. Aumentare \(f_s\) riduce ripple di corrente e rumore acustico, ma aumenta perdite di commutazione, riscaldamento del driver e problemi di compatibilita` elettromagnetica. Ridurre \(f_s\) semplifica il lavoro dei MOSFET, ma rende piu` evidente il ripple.

\slidepair{images14/page-22.png}{images14/page-23.png}{Slide 440--441: il motore filtra il PWM; la frequenza va scelta confrontandola con le costanti di tempo elettrica e meccanica.}

Il motore non segue perfettamente ogni impulso del PWM. L'induttanza dell'armatura filtra la corrente e l'inerzia meccanica filtra la velocita`. Per questo la slide dice che il motore agisce come filtro: il moto meccanico vede soprattutto la componente media, mentre il ripple resta sovrapposto alla corrente e alla coppia.

Dalla relazione:
\[
v_m\simeq D V_{dc},
\qquad
D=\frac{T_{on}}{T_s},
\]
si ottiene:
\[
T_{on}=T_s\frac{v_m}{V_{dc}}.
\]
Il verso di moto non dipende dal duty cycle da solo, ma dalla configurazione del ponte H; il duty imposta il modulo medio della tensione, mentre la coppia di interruttori scelta imposta il segno.

L'ultima slide richiama le costanti di tempo del motore:
\[
\tau_e=\frac{L}{R}\simeq 3\,\mu\text{s},
\qquad
\tau_m=\frac{RJ}{k_tk_e}\simeq 8.2\,\text{ms}.
\]
La prima descrive quanto rapidamente puo` variare la corrente, la seconda quanto rapidamente cambia la velocita`. La slide usa la dinamica piu` lenta come riferimento per scegliere una frequenza PWM sufficientemente alta. Usando \(1/\tau_m\) come frequenza caratteristica:
\[
\frac{1}{\tau_m}\simeq \frac{1}{8.2\,\text{ms}}\simeq 122\,\text{Hz}.
\]
Andare circa una decade sopra porta all'ordine di grandezza:
\[
f_{\text{PWM}}>1.2\,\text{kHz}.
\]
Poi, nella pratica, si scelgono spesso frequenze ben piu` alte, coerenti con l'intervallo \(10\text{--}40\,\text{kHz}\), per ridurre rumore e ripple.

\begin{infobox}
\textbf{Attenzione alle unita`.} Prima di usare \(L/R\) o \(RJ/(k_tk_e)\) bisogna convertire tutto in unita` SI coerenti. Microhenry, millisecondi, mN\,m e g\,cm\(^2\) non possono essere combinati direttamente.
\end{infobox}

\subsection*{Sintesi finale della Lezione 15}

La prima idea della lezione e` che il BLDC trapezoidale reale non produce coppia perfettamente piatta, perche' le commutazioni di corrente richiedono tempo. Il controllo sinusoidale nasce per ridurre questo ripple quando la macchina ha back-EMF sinusoidale.

La seconda idea e` che il problema trifase si semplifica cambiando riferimento:
\[
abc \longrightarrow \alpha\beta \longrightarrow dq.
\]
La Clarke passa dalle tre fasi a un vettore statorico; la Park osserva quel vettore da un riferimento che ruota con il rotore. In \(dq\), le grandezze sinusoidali diventano quasi costanti e la coppia della macchina a magneti superficiali diventa:
\[
T_e=\frac{3}{2}p\lambda_m i_q.
\]
Da qui segue la strategia base:
\[
i_d^\star=0,
\qquad
i_q^\star\propto T_e^\star.
\]
Il FOC misura le correnti reali, le trasforma in \(dq\), controlla \(i_d\) e \(i_q\), poi torna alle tensioni di fase che l'inverter sintetizza, spesso tramite SVPWM.

La terza idea e` che il controllo deve diventare potenza elettrica reale. Per il DC brushed, il ponte H sceglie polarita`, frenatura o coasting; il PWM regola il valore medio della tensione:
\[
D=\frac{T_{on}}{T_s},
\qquad
v_m\simeq D V_{dc}.
\]
La frequenza PWM deve essere abbastanza alta rispetto alle dinamiche del motore da ridurre ripple e rumore, ma non cosi` alta da rendere dominanti perdite e stress di commutazione.

% =========================================================
% APPENDICE
% =========================================================

\clearpage
\phantomsection
\section*{Appendice}
\addcontentsline{toc}{section}{Appendice}

\phantomsection
\subsection*{Formulario sintetico}
\addcontentsline{toc}{subsection}{Formulario sintetico}

Questa raccolta finale non sostituisce il testo, ma aiuta a ripassare rapidamente le formule centrali del documento.

{\renewcommand{\arraystretch}{1.18}
\begin{longtable}{>{\raggedright\arraybackslash}p{0.24\linewidth}>{\raggedright\arraybackslash}p{0.60\linewidth}}
\toprule
\textbf{Tema} & \textbf{Formula chiave} \\
\midrule
Campionamento & $x[k]=x(kT_s)$ \\
Nyquist e banda utile & $f_s>2f_{\max}$, \quad in pratica spesso $f_s\approx 5\text{--}20\,f_B$ \\
Quantizzazione & $q=\dfrac{R}{2^N}$, \quad $e_q\in\left[-\dfrac{q}{2},\dfrac{q}{2}\right]$, \quad $\sigma_q^2=\dfrac{q^2}{12}$ \\
SNR di quantizzazione & $\text{SNR}_{\text{ideal}}\approx 6.02\,N+1.76\ \text{dB}$ \\
Calibrazione affine & $y=ax+b$, \quad $\hat\theta=(X^TX)^{-1}X^Ty$ \\
Compensazione affine inversa & $\hat x=\dfrac{y-b}{a}$ \\
Modello completo di misura & $y=f(x,\theta)+b+n$ \\
Rumore colorato & $n_{k+1}=\alpha n_k+v_k$ \\
Bias drift & $b_{k+1}=b_k+w_k$ \\
Sensore del primo ordine & $G(s)=\dfrac{K}{\tau s+1}$ \\
Primo ordine nel tempo & $\tau \dot y(t)+y(t)=K\,u(t)$, \quad $y(t)=K\,u_0\left(1-e^{-t/\tau}\right)$ \\
Accuratezza percentuale & $\varepsilon_f=100\dfrac{x_m-x_v}{x_{FS}}$, \quad $\varepsilon_a=100\dfrac{x_m-x_v}{x_v}$ \\
Precisione campionaria & $s=\sqrt{\dfrac{1}{N-1}\sum_{i=1}^{N}(x_i-\bar x)^2}$ \\
Encoder incrementale & $\theta[k]=\dfrac{2\pi}{N_c}\,n[k]$, \quad $N_c=N_e,\,2N_e,\,4N_e$ in X1, X2, X4 \\
Encoder assoluto & $\Delta\theta=\dfrac{2\pi}{2^N}$ \\
Interpolazione ottica & $x_{\text{frac}}=\dfrac{p}{2\pi}\operatorname{atan2}(d_s,d_c)$, \quad $x=kp+x_{\text{frac}}$ \\
Frange di Moir\'e & $\Lambda=\dfrac{p_1p_2}{|p_1-p_2|}\approx\dfrac{p^2}{|\Delta p|}$ \\
LVDT/RVDT lineare locale & $V_{\text{out}}(t)\propto x\sin(\omega t)$ oppure $V_{\text{out}}(t)\propto \theta\sin(\omega t)$ \\
Resolver demodulato & $V_{s,DC}=K\sin\theta$, \quad $V_{c,DC}=K\cos\theta$, \quad $\theta=\operatorname{atan2}(V_{s,DC},V_{c,DC})$ \\
Tracking resolver-digitale & $\varepsilon=S\cos\hat\theta-C\sin\hat\theta$, \quad $\varepsilon \propto \sin(\theta-\hat\theta)\approx \theta-\hat\theta$, \quad $\dot{\hat\theta}=K_p\varepsilon+K_i\int \varepsilon\,dt$ \\
Resolver brushless & $V_{\text{rotor}}=k\,V_{\text{primary}}$, \quad $V_{\text{rotor}}(t)=V_0\sin(\omega t+\phi_e)$ \\
Synchro & $V_{s1}=K\sin(\omega t)\cos(\theta-0^\circ)$, \quad $V_{s2}=K\sin(\omega t)\cos(\theta-120^\circ)$, \quad $V_{s3}=K\sin(\omega t)\cos(\theta+120^\circ)$ \\
Errore synchro & $V_{out}=K\sin(\theta_{TX}-\theta_{CT})$ \\
Encoder magnetico assoluto & $B_x=B_0\cos\theta$, \quad $B_y=B_0\sin\theta$, \quad $\theta=\operatorname{atan2}(B_y,B_x)$ \\
Effetto Hall & $V_H=\dfrac{I\,B}{q\,n\,t}$ \\
Sensore magnetico lineare & $V_s=A\sin\!\left(\dfrac{2\pi x}{p}\right)$, \quad $V_c=A\cos\!\left(\dfrac{2\pi x}{p}\right)$, \quad $x=\dfrac{p}{2\pi}\operatorname{atan2}(V_s,V_c)+kp$ \\
Ultrasonico ToF & $c \approx 331 + 0.6\,T(^\circ\text{C})$, \quad $d=\dfrac{c\,t_{\text{echo}}}{2}$, \quad $\hat v \approx \dfrac{\Delta d}{\Delta t}$ \\
Triangolazione laser & $z \approx \dfrac{b\,f}{u}$, \quad $\hat v_z \approx \dfrac{\Delta z}{\Delta t}$ \\
Faraday su spira rotante & $\Phi(t)=BA\cos(\omega t)$, \quad $e(t)=NBA\omega\sin(\omega t)$ \\
Tachogeneratore DC & $V_{\text{out}}(t)\approx K\omega + \varepsilon(t)$ \\
Tachogeneratore AC & $V_{\text{out}}(t)=K\omega\sin(\omega t)$ \\
Momento angolare giroscopico & $\mathbf{L}=I\boldsymbol{\omega}$, \quad $\boldsymbol{\tau}=\dfrac{d\mathbf{L}}{dt}$, \quad $\omega_p=\dfrac{\tau}{L}$ \\
Precessione della trottola & $\tau=mgr\sin\phi$, \quad $\omega_p=\dfrac{\tau}{L\sin\phi}=\dfrac{mgr}{I\omega}$ \\
Rate gyro meccanico & $J_y\ddot\theta+b\dot\theta+k\theta=M_u-h\dot\psi$, \quad a regime $k\theta=-h\dot\psi$ \\
Dinamica giroscopio 2 assi & $s=0,0,\pm j\dfrac{h}{\sqrt{J_yJ_z}}$, \quad $\omega_0=\dfrac{h}{\sqrt{J_yJ_z}}$, \quad $\dot\psi_{\text{ss}}=\dfrac{M_{y0}}{h}$ \\
Rate gyro a rebalance & $G(s)=\dfrac{1}{J_y s^2}$, \quad $M_b=K(s)(\theta_{\text{REF}}-\theta)$, \quad se $G(s)K(s)\gg 1$, $M_b\simeq -h\dot\psi$ \\
MEMS vibrante a Coriolis & $x=X\sin(\omega_d t)$, \quad $\dot x=X\omega_d\cos(\omega_d t)$, \quad $a_y=-2\Omega_z\dot x$, \quad $m\ddot y+c_y\dot y+k_y y=-2m\Omega_z\dot x$ \\
Attuazione elettrostatica MEMS & $U_c=\dfrac{1}{2}CV^2$, \quad $F_x=-\dfrac{\partial U_c}{\partial x}$, \quad $C_{\text{tot}}\simeq N\varepsilon_0\dfrac{WX}{z}$, \quad $|F_x|\simeq N\dfrac{1}{2}\dfrac{k\varepsilon_0W}{z}V^2$ \\
Effetto Sagnac & $\Delta t=\dfrac{4A}{c^2}\omega$, \quad $\Delta L=c\Delta t=\dfrac{4A}{c}\omega$ \\
Fiber-Optic Gyro & $\Delta\phi=2\pi\dfrac{\Delta L}{\lambda}=\dfrac{8\pi A}{c\lambda}\omega$, \quad con $N$ spire $\Delta\phi=\dfrac{8\pi A N}{c\lambda}\omega$ \\
Ring Laser Gyro & $\nu\lambda=c$, \quad $\Delta\nu=-\dfrac{\Delta\lambda}{\lambda}\nu$, \quad $\Delta\nu=-4\dfrac{A\nu}{cL}\omega$ \\
Velocita` da ruote & $v=R\omega$, \quad $v=\dfrac{1}{N}\sum_i v_i$, \quad $v_i=R_i\omega_i$ \\
Bernoulli e tubo di Pitot & $p+\dfrac{1}{2}\rho v^2+\rho gh=\text{costante}$, \quad $p_0+\dfrac{1}{2}\rho v_0^2=p_T$, \quad $v_0=\sqrt{\dfrac{2(p_T-p_0)}{\rho}}$ \\
Doppler velocimetrico & $f_d=\dfrac{2v\cos\theta}{\lambda}=\dfrac{2vf_{\text{emit}}\cos\theta}{c}$, \quad $v=\dfrac{f_d\lambda}{2\cos\theta}$ \\
Accelerometro massa-molla & $\xi=y-x$, \quad $m\ddot\xi+b\dot\xi+k\xi=-m\ddot x$, \quad $\xi_{\text{ss}}\simeq -\dfrac{m}{k}a$ \\
Accelerometro con gravita` & $m\ddot\xi+b\dot\xi+k\xi=-m\ddot x-mg\sin\gamma$, \quad $a^{\text{dist}}=-g\sin\gamma$ \\
Accelerometro capacitivo MEMS & $C=\dfrac{\varepsilon A}{d}$, \quad $\Delta C=C_1-C_2\propto \xi \propto a$ \\
Accelerometro termico & $\Delta T \propto a$ \\
Saturazione attuatore & $u_{out}=\operatorname{sat}(u_{in})$ \\
Zona morta & $u_{out}=0$ per $|u_{in}|<u_D$ \\
Slew rate & $\left|\dfrac{du}{dt}\right|\le \dot u_{\max}$ \\
Banda attuatore & $|G(j\omega_B)|=\dfrac{|G(0)|}{\sqrt{2}}$ \quad (punto di $-3$ dB) \\
Rendimento attuatore & $\eta=\dfrac{P_{\text{mecc}}}{P_{\text{in}}}$ \\
Forza di Lorentz e Faraday & $\mathbf F=I\,\mathbf l\times \mathbf B$, \quad $e=-\dfrac{d\Phi}{dt}$ \\
Coppia di una spira in motore DC & $F=Bil$, \quad $\tau=2rF\sin\theta=BiA\sin\theta$, \quad $A=2rl$ \\
Motore DC con $N$ spire e commutazione ideale & $\tau=NBiA\sin\theta$, \quad $\tau_{\text{net}}(\theta)=\tau_{\max}|\sin\theta|$ \\
Motore DC reale e back-EMF & $\tau=k_t i$, \quad $e=k_e\omega$, \quad in SI ideale $k_t=k_e$ \\
Modello elettromeccanico motore DC & $V=Ri+L\dot i+k_e\omega$, \quad $J\dot\omega=k_t i-b\omega-\tau_L$ \\
Caratteristica coppia-velocita` motore DC & $\tau=k_t\dfrac{V}{R}-k_t\dfrac{k_e}{R}\omega$, \quad $\tau_{\text{stall}}=k_t\dfrac{V}{R}$, \quad $\omega_0=\dfrac{V}{k_e}$ \\
Potenza ed efficienza motore DC & $P_{\text{out}}=\tau\omega$, \quad $\eta=\dfrac{\tau\omega}{Vi}$, \quad $P_{\max}$ a $\omega=\omega_0/2$ e $\tau=\tau_{\text{stall}}/2$ \\
Termica motore DC & $P_J=Ri^2$, \quad $\Delta T=P_J(R_{\text{th1}}+R_{\text{th2}})$, \quad $R_{\text{warm}}=R(1+\alpha_{\text{cu}}\Delta T)$ \\
Funzione di trasferimento motore DC & $\dfrac{\Omega(s)}{V(s)}=\dfrac{k_t}{(Ls+R)(Js+b)+k_tk_e}$, \quad se $\tau_e\ll\tau_m$: $\dfrac{\Omega(s)}{V(s)}\simeq\dfrac{k_t}{RJs+Rb+k_tk_e}$ \\
Motore PM stepper & $p_r=\dfrac{N_r}{2}$, \quad $\alpha_s=\dfrac{360^\circ}{mN_r}$ \\
Motore VR stepper & $\alpha_s=360^\circ\left|\dfrac{N_s-N_r}{N_sN_r}\right|$ \\
Motore hybrid stepper & $\alpha_R=\dfrac{360^\circ}{N_r}$, \quad $\alpha_s=\dfrac{360^\circ}{2mN_r}$ \\
Coppia statica stepper & $T=T_{\max}\sin\delta$, \quad $T_{\text{pull-out}}>T_{\text{pull-in}}$ a pari frequenza di passo \\
Microstepping & $\alpha_{\text{full}}=\alpha_s$, \quad $\alpha_{\text{half}}=\dfrac{\alpha_s}{2}$, \quad $\alpha_{\mu\text{step}}\approx\dfrac{\alpha_s}{N_\mu}$ \\
Vincolo trifase a stella BLDC & $i_a+i_b+i_c=0$, \quad $\dot i_a+\dot i_b+\dot i_c=0$ \\
Modello elettrico BLDC & $\mathbf v_{abc}=R_s\mathbf i_{abc}+\dfrac{d\boldsymbol\lambda_{abc}}{dt}$, \quad $\boldsymbol\lambda_{abc}=\mathbf L_{abc}\mathbf i_{abc}+\boldsymbol\lambda_m(\theta_e)$ \\
Modello di fase BLDC ridotto & $v_k=R_s i_k+L_{eq}\dot i_k+e_k$, \quad $L_{eq}=L_s-M$, \quad $k\in\{a,b,c\}$ \\
Coppia e dinamica BLDC & $T_e=\dfrac{e_ai_a+e_bi_b+e_ci_c}{\omega_m}$, \quad $J\dot\omega_m=T_e-T_L-B\omega_m$, \quad $\theta_e=p\theta_m$ \\
Trasformata di Clarke & $\begin{bmatrix}x_\alpha\\x_\beta\end{bmatrix}=\dfrac{2}{3}\begin{bmatrix}1&-\frac12&-\frac12\\0&\frac{\sqrt3}{2}&-\frac{\sqrt3}{2}\end{bmatrix}\begin{bmatrix}x_a\\x_b\\x_c\end{bmatrix}$ \\
Trasformata di Park & $\begin{bmatrix}x_d\\x_q\end{bmatrix}=\begin{bmatrix}\cos\theta_e&\sin\theta_e\\-\sin\theta_e&\cos\theta_e\end{bmatrix}\begin{bmatrix}x_\alpha\\x_\beta\end{bmatrix}$ \\
Inversa Park/Clarke & $\begin{bmatrix}x_\alpha\\x_\beta\end{bmatrix}=\begin{bmatrix}\cos\theta_e&-\sin\theta_e\\\sin\theta_e&\cos\theta_e\end{bmatrix}\begin{bmatrix}x_d\\x_q\end{bmatrix}$, \quad $\begin{bmatrix}x_a\\x_b\\x_c\end{bmatrix}=\begin{bmatrix}1&0\\-\frac12&\frac{\sqrt3}{2}\\-\frac12&-\frac{\sqrt3}{2}\end{bmatrix}\begin{bmatrix}x_\alpha\\x_\beta\end{bmatrix}$ \\
Modello brushless in $dq$ & $v_d=R_s i_d+L_d\dot i_d-\omega_eL_qi_q$, \quad $v_q=R_s i_q+L_q\dot i_q+\omega_eL_di_d+\omega_e\lambda_m$ \\
Coppia e FOC brushless & $T_e=\dfrac{3}{2}p\left(\lambda_m i_q+(L_d-L_q)i_di_q\right)$, \quad se $L_d\simeq L_q$: $T_e=\dfrac{3}{2}p\lambda_m i_q$, \quad $i_d^\star=0$ \\
PWM per driver DC & $D=\dfrac{T_{on}}{T_s}$, \quad $f_s=\dfrac{1}{T_s}$, \quad $v_m\simeq D V_{dc}$, \quad $T_{on}=T_s\dfrac{v_m^\star}{V_{dc}}$ \\
Costanti di tempo motore DC & $\tau_e=\dfrac{L}{R}$, \quad $\tau_m=\dfrac{RJ}{k_tk_e}$, \quad scegliere $f_{\text{PWM}}$ abbastanza maggiore della dinamica piu` lenta \\
Metodo M & $\omega_m=\dfrac{2\pi}{T_c}\dfrac{m}{PPR}$, \quad $N=\dfrac{60}{T_c}\dfrac{m}{PPR}$ \\
Metodo T & $T_c=m_cT_{\text{clock}}$, \quad $\omega_m=\dfrac{2\pi}{PPR\,m_cT_{\text{clock}}}$ \\
Metodo M/T & $T=T_c+m_2T_2$, \quad $\omega_m=\dfrac{2\pi}{PPR}\dfrac{m_1}{T_c+m_2T_2}$ \\
Modello a velocita` costante & $\dot\theta=\omega$, \quad $\dot\omega=0$ \\
Osservatore di Luenberger & $\dot{\hat\theta}=\hat\omega+l_1(\theta-\hat\theta)$, \quad $\dot{\hat\omega}=l_2(\theta-\hat\theta)$, \quad $L=\begin{bmatrix}2/\varepsilon\\ 1/\varepsilon^2\end{bmatrix}$ \\
\bottomrule
\end{longtable}
}

\phantomsection
\subsection*{Confronto rapido dei sensori di posizione}
\addcontentsline{toc}{subsection}{Confronto rapido dei sensori di posizione}

Come supporto allo studio, la tabella seguente raccoglie in un colpo d’occhio i compromessi principali tra le tecnologie introdotte nelle Lezioni 4--7.

{\renewcommand{\arraystretch}{1.18}
\begin{longtable}{@{}>{\raggedright\arraybackslash}p{0.15\linewidth}>{\raggedright\arraybackslash}p{0.14\linewidth}>{\raggedright\arraybackslash}p{0.14\linewidth}>{\raggedright\arraybackslash}p{0.21\linewidth}>{\raggedright\arraybackslash}p{0.22\linewidth}@{}}
\toprule
\textbf{Sensore} & \textbf{Variabile} & \textbf{Assoluto / Incrementale} & \textbf{Punti forti} & \textbf{Limiti principali} \\
\midrule
\endfirsthead
\toprule
\textbf{Sensore} & \textbf{Variabile} & \textbf{Assoluto / Incrementale} & \textbf{Punti forti} & \textbf{Limiti principali} \\
\midrule
\endhead
Potenziometro & Posizione lineare o angolare & Assoluto & Semplice, economico, uscita immediata, facile da interfacciare & Usura, contatto strisciante, rumore di contatto, vita utile limitata \\
Encoder incrementale & Posizione e velocit\`a tramite conteggio & Incrementale & Alta risoluzione, uscita digitale, ottimo per servoazionamenti & Necessita homing o riferimento iniziale; se perde conteggio perde posizione \\
Encoder assoluto & Posizione angolare entro un giro; nei multi-turn anche numero di giri & Assoluto & Posizione disponibile subito, niente homing, buona integrazione digitale & Costo pi\`u alto, complessit\`a maggiore, sensibilit\`a alla qualit\`a della codifica \\
Encoder magnetico relativo & Posizione angolare tramite conteggio di poli o denti & Incrementale & Robusto a sporco, olio e vibrazioni; semplice e contactless & Richiede homing; risoluzione tipicamente inferiore agli encoder ottici \\
Sensore lineare magnetico & Posizione lineare & Incrementale o assoluto & Robusto all’ambiente, compatto, interpolabile con segnali seno-coseno & Sensibile a distorsioni del campo, deriva termica e invecchiamento del magnete \\
Encoder ottico interpolato & Posizione lineare o angolare con stima di fase & Tipicamente incrementale con riferimento opzionale & Risoluzione molto alta, buona precisione metrica, ottima qualit\`a del segnale se ben condizionato & Richiede elettronica analogica accurata; sensibile a sfasamenti, offset e allineamenti \\
LVDT & Posizione lineare & Assoluto locale attorno allo zero & Alta linearit\`a, assenza di contatto, robustezza, ripetibilit\`a elevata & Range limitato, alimentazione AC, demodulazione necessaria, costo non basso \\
RVDT & Posizione angolare su intervallo limitato & Assoluto locale attorno allo zero & Robusto, contactless, adatto ad ambienti severi & Range angolare limitato, elettronica di eccitazione e demodulazione \\
Resolver & Posizione angolare continua su $360^\circ$ & Assoluto continuo & Molto robusto, tollera sporco, vibrazioni, temperatura; ottimo in ambito industriale e aerospaziale & Uscite analogiche modulate; richiede RDC o demodulazione dedicata; risoluzione nativa inferiore ai migliori encoder ottici \\
Synchro & Posizione angolare e trasmissione remota dell’angolo & Assoluto analogico & Utile in sistemi master-slave, robusto in ambienti severi, adatto a trasmissione elettrica dell’informazione angolare & Ingombrante, poco integrato con l’elettronica digitale moderna, spesso usa slip rings \\
Encoder magnetico assoluto & Posizione angolare & Assoluto & Posizione disponibile all’avvio, compatto, economico, robusto a polvere e olio & Sensibile a eccentricita` del magnete, offset e distorsioni del campo; precisione tipicamente inferiore ai migliori encoder ottici \\
\bottomrule
\end{longtable}
}

\phantomsection
\subsection*{Confronto rapido di sensori e metodi per velocita`, distanza e accelerazione}
\addcontentsline{toc}{subsection}{Confronto rapido di sensori e metodi per velocita`, distanza e accelerazione}

Per completare il quadro, la tabella seguente riassume le tecnologie e i metodi introdotti nelle Lezioni 7--11 per la misura di velocita`, distanza e accelerazione.

{\renewcommand{\arraystretch}{1.18}
\begin{longtable}{@{}>{\raggedright\arraybackslash}p{0.17\linewidth}>{\raggedright\arraybackslash}p{0.13\linewidth}>{\raggedright\arraybackslash}p{0.18\linewidth}>{\raggedright\arraybackslash}p{0.20\linewidth}>{\raggedright\arraybackslash}p{0.18\linewidth}@{}}
\toprule
\textbf{Tecnologia / Metodo} & \textbf{Variabile} & \textbf{Punti forti} & \textbf{Limiti principali} & \textbf{Uso tipico} \\
\midrule
\endfirsthead
\toprule
\textbf{Tecnologia / Metodo} & \textbf{Variabile} & \textbf{Punti forti} & \textbf{Limiti principali} & \textbf{Uso tipico} \\
\midrule
\endhead
Tachogeneratore DC & Velocit\`a angolare & Uscita analogica diretta, buona linearit\`a, segno del moto immediato & Spazzole, usura, ripple, rumore di commutazione, prestazioni scarse a bassa velocit\`a & Servoazionamenti tradizionali e feedback analogico \\
Tachogeneratore AC & Velocit\`a angolare & Niente collettore, minore usura, ampiezza del segnale proporzionale alla velocit\`a & Uscita AC da rettificare o contare; condizionamento necessario & Misura diretta robusta della velocit\`a \\
Giroscopio / rate gyro & Velocit\`a angolare propria & Misura inerziale diretta, utile per assetto, stabilizzazione e navigazione & Richiede taratura, smorzamento e gestione di deriva e disturbi; nei meccanici classici e` complesso & Avionica, IMU, veicoli mobili e controllo d'assetto \\
FOG / RLG & Velocit\`a angolare propria & Nessuna parte mobile; misura ottica basata su effetto Sagnac; alta affidabilit\`a & Richiede componenti ottici, elettronica interferometrica o conteggio di battimento; costo e compensazioni non banali & Navigazione inerziale, girobussole, avionica e sistemi triassiali ad alte prestazioni \\
MEMS vibratorio & Velocit\`a angolare propria & Piccolo, economico, basso consumo, integrabile in IMU compatte & Bias drift, rumore e sensibilit\`a a temperatura/vibrazioni; richiede calibrazione e compensazione & Smartphone, droni, robot mobili, stabilizzazione e automotive \\
Ruote + encoder & Velocit\`a lineare propria & Alta banda, basso costo, funziona indoor, ottimo per controllo locale & Slittamento, raggio effettivo variabile, errori su terreno irregolare e in curva & Robot mobili, AGV, veicoli terrestri \\
Tubo di Pitot & Velocit\`a rispetto all'aria & Misura direttamente l'airspeed utile al controllo aerodinamico & Dipende da densit\`a e modello di flusso; sensibile a ghiaccio, ostruzioni e prese statiche errate & Aeronautica, autopiloti, sonde aria \\
Doppler laser / radar & Velocit\`a relativa lungo il fascio & Misura contactless, alta risoluzione in frequenza, utile anche ad alta velocit\`a & Misura solo la componente radiale; richiede bersaglio/riflessione e buon allineamento & Radar automotive, misura industriale di velocit\`a, vibrazioni con LDV \\
DVL & Velocit\`a propria rispetto al fondale & Buono per navigazione subacquea, usa pi\`u fasci e stima vettoriale della velocit\`a & Richiede fondale osservabile e stazionario; prestazioni dipendono dall'acustica e dalla quota & AUV, ROV, navigazione inerziale assistita underwater \\
Accelerometro inerziale & Accelera\-zione / forza specifica & Misura inerziale diretta, fondamentale nelle IMU, buona banda e risposta rapida & Risente della gravita`, del bias e delle vibrazioni; doppia integrazione molto rumorosa per ricavare posizione & IMU, assetto, stabilizzazione, airbag, robotica mobile \\
Accelerometro MEMS capacitivo & Accelera\-zione & Molto compatto, economico, facile da integrare su 3 assi nello stesso die & Sensibile a deriva termica, bias e variazioni di processo; i modelli economici hanno rumore non trascurabile & Smartphone, droni, wearable, automotive, elettronica embedded \\
Accelerometro termico & Accelera\-zione & Nessuna massa solida mobile, robusto agli urti e alla fabbricazione & Banda bassa, forte dipendenza dal gas interno e dalla gestione termica & Misure low-frequency, applicazioni robuste, ambienti severi \\
Encoder + metodo M & Velocit\`a da posizione & Campionamento fisso, facile da integrare nel controllo digitale & Quantizzazione forte a bassa velocit\`a; compromesso tra errore e banda & Azionamenti con velocit\`a medio-alte \\
Encoder + metodo T / M/T & Velocit\`a da posizione & Migliore accuratezza a bassa velocit\`a; M/T offre compromesso migliore & Campionamento non uniforme o maggiore complessit\`a; stima ancora media e non istantanea & Basse velocit\`a o stime pi\`u accurate da encoder \\
Osservatore / filtro di stato & Posizione e velocit\`a stimate & Filtra rumore, regolarizza la stima e sfrutta il modello dinamico & Va tarato; introduce compromessi tra banda, rumore e ritardo; dipende dal modello & Servo ad alte prestazioni e stima robusta di velocit\`a \\
Ultrasonico ToF & Distanza & Economico, assoluto, robusto, adatto a ostacoli e livello & Bassa banda, fascio ampio, sensibilit\`a a temperatura e geometria del target & Robotica mobile semplice, serbatoi, prossimit\`a \\
Triangolazione laser & Distanza & Rapido, non-contact, buona precisione a corto-medio raggio & Sensibile a riflettivit\`a della superficie e luce ambiente & Evitamento ostacoli ravvicinati, assi lineari e macchine industriali \\
\bottomrule
\end{longtable}
}

\phantomsection
\subsection*{Riferimento al materiale originale}
\addcontentsline{toc}{subsection}{Riferimento al materiale originale}

Questo elaborato è una traduzione e rielaborazione estesa del PDF caricato dall’utente, relativo al corso ``Tecnologie per il Controllo dei Robot'' dell’Università di Pisa.

\end{document}
